% Généralités sur les fonctions numériques — Mathématiques 1re Tech — Cours signé OnBuch+
% Programme officiel MINESEC. Compiler avec : tectonic N12-generalites-fonctions-numeriques.tex
\newcommand{\DOCMATIERE}{Mathématiques}
\newcommand{\DOCNIVEAU}{1re Tech}
\newcommand{\DOCMODULE}{Mathématiques – classe de Première technique}
\newcommand{\DOCLECON}{N12}
\newcommand{\DOCTITRE}{Généralités sur les fonctions numériques}
% OnBuch+ — préambule « cours signés » ENRICHI (compilé avec tectonic / XeLaTeX)
% Système visuel « Pop » d'OnBuch+ : fond crème (jamais blanc), bordures encre
% épaisses, ombres « dures » décalées, titres Archivo Black en capitales.
% Version MATHÉMATIQUES : graphes (pgfplots/tikz), boîtes pédagogiques
% (définition, propriété, méthode, attention, le savais-tu, exercice résolu,
% exercice, TP, bilan), page de garde et signature OnBuch+ ancrée en bas.
%
% Macros attendues dans main.tex AVANT \input{preamble} :
%   \DOCMATIERE  \DOCNIVEAU  \DOCMODULE  \DOCLECON (numéro)  \DOCTITRE
\documentclass[11pt]{article}

\usepackage{fontspec}
\usepackage[french]{babel}
\usepackage[a4paper,margin=2cm,top=3.4cm,bottom=2.8cm,headheight=1.4cm,footskip=1.3cm]{geometry}
\usepackage{amsmath,amssymb,amsfonts}
\usepackage{mathtools} % \xmapsto, \coloneqq... souvent utilisés par les modèles
\usepackage{cancel} % simplifications barrées
% \abs{z} (module, valeur absolue) et \norm{u} (norme), utilisés par les modèles
\providecommand{\abs}{}\let\abs\relax\DeclarePairedDelimiter{\abs}{\lvert}{\rvert}
\providecommand{\norm}{}\let\norm\relax\DeclarePairedDelimiter{\norm}{\lVert}{\rVert}
\usepackage[table]{xcolor} % [table] : fournit aussi \rowcolors (alternance de lignes)
\usepackage{pagecolor}
\usepackage{tikz}
\usetikzlibrary{arrows.meta,positioning,calc,shapes.geometric,shapes.misc,shapes.arrows,decorations.pathmorphing,decorations.pathreplacing,decorations.markings,patterns,shadows,mindmap,trees,fit,backgrounds,matrix,intersections,angles,quotes}
\usepackage{pgfplots}
\usepackage{pgf-pie} % diagrammes circulaires \pie (statistiques)
\pgfplotsset{compat=1.18}
\usepgfplotslibrary{fillbetween,groupplots} % aires entre courbes (calcul intégral)
\usepackage{enumitem}
\usepackage{fancyhdr}
% [most] charge toutes les bibliothèques tcolorbox (listings, minted, raster,
% documentation...) et gonfle inutilement la mémoire TeX sur un document long
% (~300 boîtes) : on ne charge que ce qui est réellement utilisé.
\usepackage[skins,breakable]{tcolorbox}
\usepackage{graphicx}
\usepackage{colortbl}
\usepackage{booktabs}
\usepackage{tabularx}
\usepackage{diagbox} % case d'en-tête diagonale des tableaux à double entrée
% Colonnes extensibles centrée / gauche / droite (C, L, R), souvent utilisées par les modèles
\newcolumntype{C}{>{\centering\arraybackslash}X}
\newcolumntype{L}{>{\raggedright\arraybackslash}X}
\newcolumntype{R}{>{\raggedleft\arraybackslash}X}
\usepackage{array}
\usepackage{multicol}
\usepackage{siunitx}
% parse-numbers=false : accepte aussi \SI{2,5 \times 10^{-3}}{...}
\sisetup{locale=FR,output-decimal-marker={,},group-minimum-digits=4,parse-numbers=false}
\DeclareSIUnit\ohm{\ensuremath{\symup{\Omega}}} % U+2126 absent de la police
\usepackage{qrcode}
% \degree : commande fréquemment utilisée par les modèles pour un angle ou une
% température en dehors de \SI{}{} (non fournie par les paquets chargés).
\providecommand{\degree}{^{\circ}}
% \qed (fin de démonstration), utilisé par les modèles sans amsthm
\providecommand{\qed}{\hfill$\square$}
% \barre : double barre des valeurs interdites dans un tableau de variations (tkz-tab)
\providecommand{\barre}{\ensuremath{\|}}
% environnement proof (amsthm non chargé) utilisé par les modèles
\ifdefined\proof\else\newenvironment{proof}[1][Démonstration]{\par\noindent\textit{#1.}\ }{\qed\par}\fi
\usepackage{lastpage}
\usepackage{hyperref}
\hypersetup{hidelinks}

% ============================================================
% Palette « Pop » OnBuch+ — mêmes hex que la classe P (plus_ui.dart)
% ============================================================
\definecolor{poporange}{HTML}{F59321}
\definecolor{popdark}{HTML}{E07A0C}
\definecolor{popcream}{HTML}{FFF6E8}
\definecolor{popink}{HTML}{1C1714}
\definecolor{poppurple}{HTML}{7A5AE0}
\definecolor{popgreen}{HTML}{2E9E6A}
\definecolor{poppink}{HTML}{E0457B}
\definecolor{popblue}{HTML}{2F7DE1}
\definecolor{popgold}{HTML}{C98A12}
\definecolor{popmuted}{HTML}{8A7F76}
\definecolor{popline}{HTML}{EADFD2}
\definecolor{popgray}{HTML}{8A7F76}
\colorlet{popgrey}{popgray}
% Alias tolérants (le modèle nomme parfois les couleurs différemment)
\colorlet{popwhite}{white}
\colorlet{popblack}{popink}
\colorlet{popred}{poppink}
\colorlet{popyellow}{popgold}
% Teintes claires pour fonds de tableaux / zones
\colorlet{popblueL}{popblue!10!white}
\colorlet{poporangeL}{poporange!14!white}
\colorlet{popgreenL}{popgreen!12!white}
\colorlet{poppurpleL}{poppurple!12!white}
\colorlet{poppinkL}{poppink!12!white}
\colorlet{popgoldL}{popgold!14!white}

\pagecolor{popcream}

% ============================================================
% Typographie
% ============================================================
\defaultfontfeatures{Ligatures=TeX}
\setmainfont{PlusJakartaSans-Medium.ttf}[Path=../../../fonts/,BoldFont=PlusJakartaSans-Bold.ttf]
% Maths en sans-serif (Fira Math) avec chiffres et lettres droites de Plus
% Jakarta, pour que les formules se fondent dans le texte.
\usepackage[math-style=ISO,bold-style=ISO]{unicode-math}
\setmathfont{FiraMath-Regular.otf}[Path=../../../fonts/]
\setmathfont{PlusJakartaSans-Medium.ttf}[Path=../../../fonts/,range={up/num,up/{Latin,latin},\mathup}]
\usepackage{icomma} % virgule décimale sans espace en mode math
% Caractères mathématiques Unicode absents des polices de texte (Plus Jakarta,
% Archivo Black) : rendus par leur équivalent mathématique, en texte comme en
% formule — sinon ils s'affichent en carré vide (ex. « Arithmétique dans ℤ »).
\usepackage{newunicodechar}
\newunicodechar{ℝ}{\ensuremath{\mathbb{R}}}
\newunicodechar{ℤ}{\ensuremath{\mathbb{Z}}}
\newunicodechar{ℂ}{\ensuremath{\mathbb{C}}}
\newunicodechar{ℕ}{\ensuremath{\mathbb{N}}}
\newunicodechar{ℚ}{\ensuremath{\mathbb{Q}}}
\newunicodechar{𝔻}{\ensuremath{\mathbb{D}}}
\newunicodechar{α}{\ensuremath{\alpha}}
\newunicodechar{β}{\ensuremath{\beta}}
\newunicodechar{γ}{\ensuremath{\gamma}}
\newunicodechar{δ}{\ensuremath{\delta}}
\newunicodechar{ε}{\ensuremath{\varepsilon}}
\newunicodechar{θ}{\ensuremath{\theta}}
\newunicodechar{λ}{\ensuremath{\lambda}}
\newunicodechar{μ}{\ensuremath{\mu}}
\newunicodechar{σ}{\ensuremath{\sigma}}
\newunicodechar{τ}{\ensuremath{\tau}}
\newunicodechar{φ}{\ensuremath{\varphi}}
\newunicodechar{ψ}{\ensuremath{\psi}}
\newunicodechar{ω}{\ensuremath{\omega}}
\newunicodechar{Σ}{\ensuremath{\Sigma}}
% majuscules produites par \MakeUppercase dans les titres (α-aminés -> Α)
\newunicodechar{Α}{\ensuremath{\alpha}}
\newunicodechar{Β}{\ensuremath{\beta}}
\newunicodechar{Γ}{\ensuremath{\Gamma}}
\newunicodechar{Δ}{\ensuremath{\Delta}}
\newunicodechar{Ε}{\ensuremath{\varepsilon}}
\newunicodechar{Θ}{\ensuremath{\Theta}}
\newunicodechar{Λ}{\ensuremath{\Lambda}}
\newunicodechar{Μ}{\ensuremath{\mu}}
\newunicodechar{Τ}{\ensuremath{\tau}}
\newunicodechar{Φ}{\ensuremath{\Phi}}
\newunicodechar{Ψ}{\ensuremath{\Psi}}
\newunicodechar{Ω}{\ensuremath{\Omega}}
\newunicodechar{↦}{\ensuremath{\mapsto}}
\newunicodechar{∈}{\ensuremath{\in}}
\newunicodechar{∉}{\ensuremath{\notin}}
\newunicodechar{∧}{\ensuremath{\wedge}}
\newunicodechar{⇔}{\ensuremath{\Leftrightarrow}}
\newunicodechar{⟺}{\ensuremath{\Longleftrightarrow}}
\newunicodechar{⇒}{\ensuremath{\Rightarrow}}
\newunicodechar{⟹}{\ensuremath{\Longrightarrow}}
\newunicodechar{∘}{\ensuremath{\circ}}
\newunicodechar{∪}{\ensuremath{\cup}}
\newunicodechar{∩}{\ensuremath{\cap}}
\newunicodechar{≡}{\ensuremath{\equiv}}
\newunicodechar{⊂}{\ensuremath{\subset}}
\newunicodechar{⊥}{\ensuremath{\perp}}
\newunicodechar{∃}{\ensuremath{\exists}}
\newunicodechar{∀}{\ensuremath{\forall}}
\newunicodechar{∛}{\ensuremath{\sqrt[3]{\,}}}
\newunicodechar{∜}{\ensuremath{\sqrt[4]{\,}}}
\newunicodechar{⃗}{\ensuremath{{}^{\to}}}
\newunicodechar{ⁿ}{\ifmmode^{n}\else\textsuperscript{\ensuremath{n}}\fi}
\newunicodechar{ˣ}{\ifmmode^{x}\else\textsuperscript{\ensuremath{x}}\fi}
\newunicodechar{ᵇ}{\ifmmode^{b}\else\textsuperscript{\ensuremath{b}}\fi}
\newunicodechar{ʳ}{\ifmmode^{r}\else\textsuperscript{\ensuremath{r}}\fi}
\newunicodechar{ᵘ}{\ifmmode^{u}\else\textsuperscript{\ensuremath{u}}\fi}
\newunicodechar{ʸ}{\ifmmode^{y}\else\textsuperscript{\ensuremath{y}}\fi}
\newunicodechar{ᵃ}{\ifmmode^{a}\else\textsuperscript{\ensuremath{a}}\fi}
\newunicodechar{ᵉ}{\ifmmode^{e}\else\textsuperscript{\ensuremath{e}}\fi}
\newunicodechar{ᵏ}{\ifmmode^{k}\else\textsuperscript{\ensuremath{k}}\fi}
\newunicodechar{ᵖ}{\ifmmode^{p}\else\textsuperscript{\ensuremath{p}}\fi}
\newunicodechar{ᴺ}{\ifmmode^{N}\else\textsuperscript{\ensuremath{N}}\fi}
\newunicodechar{ᶜ}{\ifmmode^{c}\else\textsuperscript{\ensuremath{c}}\fi}
\newunicodechar{ʰ}{\ifmmode^{h}\else\textsuperscript{\ensuremath{h}}\fi}
\newunicodechar{ᵠ}{\ifmmode^{q}\else\textsuperscript{\ensuremath{q}}\fi}
\newunicodechar{ₙ}{\ifmmode_{n}\else\textsubscript{\ensuremath{n}}\fi}
\newunicodechar{ₐ}{\ifmmode_{a}\else\textsubscript{\ensuremath{a}}\fi}
\newunicodechar{ᵢ}{\ifmmode_{i}\else\textsubscript{\ensuremath{i}}\fi}
\newunicodechar{ᵣ}{\ifmmode_{r}\else\textsubscript{\ensuremath{r}}\fi}
\newunicodechar{ₖ}{\ifmmode_{k}\else\textsubscript{\ensuremath{k}}\fi}
\newunicodechar{ᵦ}{\ifmmode_{\beta}\else\textsubscript{\ensuremath{\beta}}\fi}
\newunicodechar{ₓ}{\ifmmode_{x}\else\textsubscript{\ensuremath{x}}\fi}
\newfontfamily\popdisplay{ArchivoBlack-Regular.ttf}[Path=../../../fonts/]
\newfontfamily\poplogo{SpaceGrotesk-Bold.ttf}[Path=../../../fonts/]
\newfontfamily\popsemi{PlusJakartaSans-SemiBold.ttf}[Path=../../../fonts/]
\newfontfamily\popxbold{PlusJakartaSans-ExtraBold.ttf}[Path=../../../fonts/]
\newfontfamily\popxboldls{PlusJakartaSans-ExtraBold.ttf}[Path=../../../fonts/,LetterSpace=3.0]

\setlist[enumerate]{leftmargin=1.7em,itemsep=5pt,topsep=4pt}
\setlist[itemize]{leftmargin=1.7em,itemsep=4pt,topsep=4pt,label={\color{poporange}\textbullet}}
\setlength{\parindent}{0pt}
\setlength{\parskip}{5pt}
\renewcommand{\arraystretch}{1.25}

% pgfplots : style Pop par défaut
\pgfplotsset{
  every axis/.append style={axis line style={popink,line width=1pt},tick style={popink},
    grid style={popline,line width=0.5pt},label style={font=\small},tick label style={font=\footnotesize}},
  popcurve/.style={poporange,line width=1.8pt,no markers,smooth},
}
% Même style disponible dans un \draw TikZ simple (hors pgfplots)
\tikzset{popcurve/.style={poporange,line width=1.8pt}}

% ============================================================
% Pill / squiggle
% ============================================================
\newcommand{\poppill}[3]{%
  \tikz[baseline=-0.55ex]\node[draw=popink,line width=1.2pt,rounded corners=2.6pt,%
    fill=#2,inner xsep=6.5pt,inner ysep=3pt]{%
    \popxboldls\fontsize{7}{7}\selectfont\color{#3}\MakeUppercase{#1}};%
}
\newcommand{\popsquiggle}[1]{%
  \tikz[baseline=-0.4ex]{
    \draw[#1,line width=2.6pt,line cap=round]
      plot [smooth] coordinates {(0,0.09) (0.34,0.01) (0.68,0.21) (1.02,0.01) (1.36,0.10)};
  }%
}
% Mot-clé surligné (marqueur orange sous le texte)
\newcommand{\cle}[1]{{\popxbold\color{popdark}#1}}

% ============================================================
% En-tête / pied
% ============================================================
\pagestyle{fancy}
\fancyhf{}
\renewcommand{\headrulewidth}{0pt}
\renewcommand{\footrulewidth}{0pt}
\lhead{\raisebox{-0.15ex}{\poplogo\fontsize{13}{13}\selectfont\color{popink}OnBuch\color{poporange}+}}
\rhead{\poppill{\DOCMATIERE\ \textbullet\ \DOCNIVEAU\ \textbullet\ Leçon \DOCLECON}{white}{popink}}
\lfoot{\popxbold\fontsize{7.6}{7.6}\selectfont\color{popmuted}\MakeUppercase{Cours signé OnBuch+}\ \textbullet\ {\popsemi\DOCTITRE}}
\rfoot{\tikz[baseline=-0.6ex]\node[rounded corners=8.5pt,fill=popink,minimum height=17pt,minimum width=17pt,inner xsep=5pt,inner sep=0]{\popxbold\fontsize{8}{8}\selectfont\color{popcream}\thepage\,/\,\pageref*{LastPage}};}

% ============================================================
% Titres
% ============================================================
\newcommand{\chaptitre}[2]{%
  \begin{tcolorbox}[enhanced,colback=poporange,colframe=popink,boxrule=2.6pt,arc=11pt,
      drop shadow=popink,left=15pt,right=15pt,top=12pt,bottom=11pt,before skip=3pt,after skip=18pt]
    \poppill{#2}{popink}{popcream}\par\vspace{9pt}
    {\raggedright\hyphenpenalty=10000\exhyphenpenalty=10000\popdisplay\fontsize{22}{24}\selectfont\color{popcream}\MakeUppercase{#1}\par}
  \end{tcolorbox}
}
% Section numérotée : pastille orange avec le numéro + titre Archivo
\newcounter{popsec}
\newcommand{\coursec}[1]{%
  \stepcounter{popsec}\setcounter{popsub}{0}%
  \par\vspace{16pt}\noindent
  \begin{minipage}{\linewidth}
  \tikz[baseline=-0.7ex]\node[draw=popink,line width=1.6pt,fill=poporange,rounded corners=4pt,minimum size=22pt,inner sep=0]{\popdisplay\fontsize{12}{12}\selectfont\color{popcream}\thepopsec};%
  \hspace{8pt}{\popdisplay\fontsize{15}{17}\selectfont\color{popink}\MakeUppercase{#1}}\par
  \vspace{4pt}\noindent{\color{poporange}\rule{36pt}{3.6pt}}
  \end{minipage}\par\nopagebreak\vspace{8pt}
}
\newcounter{popsub}
\newcommand{\courssub}[1]{%
  \stepcounter{popsub}%
  \par\vspace{9pt}\noindent{\popxbold\fontsize{11.5}{13}\selectfont\color{popink}{\color{poporange}\thepopsec.\thepopsub}\hspace{6pt}#1}\par\nopagebreak\vspace{4pt}
}

% ============================================================
% Boîtes pédagogiques — même recette : fond blanc, bordure épaisse colorée,
% ombre dure encre, badge plein. Titre optionnel [#1].
% ============================================================
\tcbset{popbase/.style={breakable,enhanced,colback=white,boxrule=2.3pt,arc=9pt,
  shadow={3pt}{-3pt}{0pt}{popink},left=14pt,right=14pt,top=11pt,bottom=11pt,before skip=11pt,after skip=15pt,
  fontupper=\popsemi\fontsize{10.6}{14.5}\selectfont\color{popink},
  fontlower=\popsemi\fontsize{10.6}{14.5}\selectfont\color{popink}}}
% Coche dessinée (le glyphe ✓ n'existe pas dans les polices du document)
\newcommand{\popcheck}[1]{\tikz[baseline=-0.5ex]\draw[#1,line width=1.6pt,line cap=round,line join=round](-0.13,0.0)--(-0.04,-0.09)--(0.13,0.1);}
\newcommand{\popboxhead}[3]{\poppill{#1}{#2}{white}\ifx\relax#3\relax\else\hspace{6pt}{\popxbold\fontsize{10.6}{12}\selectfont\color{popink}#3}\fi\par\vspace{7pt}}

\newtcolorbox{aretenir}[1][]{popbase,colframe=popgreen,before upper={\popboxhead{À retenir}{popgreen}{#1}}}
\newtcolorbox{exemplebox}[1][]{popbase,colframe=poporange,before upper={\popboxhead{Exemple}{poporange}{#1}}}
\newtcolorbox{definition}[1][]{popbase,colframe=popblue,before upper={\popboxhead{Définition}{popblue}{#1}}}
\newtcolorbox{propriete}[1][]{popbase,colframe=poppurple,before upper={\popboxhead{Propriété}{poppurple}{#1}}}
\newtcolorbox{methode}[1][]{popbase,colframe=popgold,before upper={\popboxhead{Méthode}{popgold}{#1}}}
\newtcolorbox{attention}[1][]{popbase,colframe=poppink,before upper={\popboxhead{Attention}{poppink}{#1}}}
\newtcolorbox{experience}[1][]{popbase,colframe=popblue,colback=popblueL,before upper={\popboxhead{Expérience}{popblue}{#1}}}
% « Le savais-tu ? » : carte violette pleine (contexte, culture, Cameroun)
\newtcolorbox{savaistu}[1][]{popbase,colback=poppurple,colframe=popink,
  fontupper=\popsemi\fontsize{10.4}{14.2}\selectfont\color{white},
  before upper={\poppill{Le savais-tu\,?}{white}{poppurple}\ifx\relax#1\relax\else\hspace{6pt}{\popxbold\color{white}#1}\fi\par\vspace{7pt}}}
% Exercice résolu : énoncé en haut, \tcblower puis la solution sur fond crème
\newcounter{popexo}\newcounter{popexr}
\newtcolorbox{exoresolu}[1][]{popbase,colframe=poporange,colbacklower=popcream,
  segmentation style={popink,line width=1.2pt,dashed},
  before upper={\stepcounter{popexr}\popboxhead{Exercice résolu \thepopexr}{poporange}{#1}},
  before lower={\poppill{Solution}{popink}{popcream}\par\vspace{6pt}}}
% Exercice (non résolu en place) — la correction est dans la partie Corrigés
\newtcolorbox{exercice}[1][]{popbase,colframe=popink,
  before upper={\stepcounter{popexo}\popboxhead{Exercice \thepopexo}{popink}{#1}}}
% Corrigé
\newtcolorbox{corrige}[1][]{popbase,colframe=popgreen,colback=popgreenL,
  before upper={\popboxhead{Corrigé}{popgreen}{#1}}}
% Objectifs / prérequis / bilan
\newtcolorbox{objectifs}{popbase,colframe=popink,colback=poporangeL,
  before upper={\popboxhead{Objectifs de la leçon}{popink}{}}}
\newtcolorbox{prerequis}{popbase,colframe=popink,colback=popblueL,
  before upper={\popboxhead{Prérequis}{popblue}{}}}
\newtcolorbox{bilan}{popbase,colframe=popink,colback=white,boxrule=2.8pt,
  before upper={\popboxhead{Fiche bilan}{poporange}{L'essentiel à connaître pour l'examen}}}
% Figure encadrée avec légende
\newtcolorbox{popfigure}[1][]{enhanced,breakable=false,colback=white,colframe=popline,boxrule=1.4pt,arc=8pt,
  left=10pt,right=10pt,top=10pt,bottom=8pt,before skip=10pt,after skip=12pt,halign=center,
  lower separated=false,halign lower=center,
  fontlower=\popsemi\fontsize{9}{11.5}\selectfont\color{popmuted},
  #1}
\newcommand{\legende}[1]{\par\vspace{6pt}{\popsemi\fontsize{9}{11.5}\selectfont\color{popmuted}#1}}

% ============================================================
% Signature OnBuch+
% ============================================================
\newcommand{\poplogoinline}[1]{{\poplogo\fontsize{#1}{#1}\selectfont\color{popink}OnBuch\color{poporange}+}}
\newcommand{\signatureonbuch}{%
  \begin{tcolorbox}[enhanced,colback=popink,colframe=popink,boxrule=2.2pt,arc=10pt,
      drop shadow=poporange,left=14pt,right=14pt,top=10pt,bottom=10pt,before skip=0pt,after skip=0pt]
    \tikz[baseline=-0.7ex]\node[circle,fill=popgreen,draw=popcream,line width=1.3pt,minimum size=22pt,inner sep=0]{\popcheck{white}};%
    \hspace{9pt}\begin{minipage}[c]{0.62\linewidth}
      {\popxbold\fontsize{12}{13}\selectfont\color{popcream}Signé\ \textbullet\ {\poplogo OnBuch\color{poporange}+}}\par\vspace{2pt}
      {\raggedright\popsemi\fontsize{8}{10.5}\selectfont\color{popline}\MakeUppercase{Équipe pédagogique OnBuch+}\\\MakeUppercase{Conforme au programme officiel MINESEC}\par}
    \end{minipage}\hfill
    {\poplogo\fontsize{22}{22}\selectfont\color{popcream}OnBuch\color{poporange}+}
  \end{tcolorbox}%
}

% ============================================================
% Page de garde
% #1 titre · #2 matière · #3 niveau · #4 module · #5 n° leçon · #6 durée ·
% #7 description / objectifs (texte ou itemize)
% ============================================================
\newcommand{\pagegarde}[7]{%
  \thispagestyle{empty}
  \newgeometry{a4paper,margin=1.7cm,top=1.6cm,bottom=1.4cm}%
  \begin{tikzpicture}[remember picture,overlay]
    \fill[poporange!18] ([xshift=-3.2cm,yshift=-3.6cm]current page.north east) circle (4.6cm);
    \fill[poppurple!14] ([xshift=2.4cm,yshift=7.6cm]current page.south west) circle (3.8cm);
  \end{tikzpicture}%
  {\poplogo\fontsize{30}{30}\selectfont\color{popink}OnBuch\color{poporange}+}\hfill
  \raisebox{0.6ex}{\poppill{Cours signé \textbullet\ Premium}{popink}{popcream}}\par
  \vspace{1pt}\popsquiggle{poppurple}\par
  \vspace{8pt}
  \begin{tcolorbox}[enhanced,colback=poporange,colframe=popink,boxrule=2.8pt,arc=13pt,
      drop shadow=popink,left=18pt,right=18pt,top=12pt,bottom=13pt,after skip=10pt]
    \poppill{#2 \textbullet\ #3}{popink}{popcream}\hspace{5pt}\poppill{#4}{white}{popink}\par\vspace{8pt}
    {\popdisplay\fontsize{14}{14}\selectfont\color{popink}LEÇON #5}\par\vspace{4pt}
    {\raggedright\hyphenpenalty=10000\exhyphenpenalty=10000\popdisplay\fontsize{25}{27}\selectfont\color{popcream}\MakeUppercase{#1}\par}
    \vspace{8pt}
    {\popxbold\fontsize{9.5}{9.5}\selectfont\color{popink}\MakeUppercase{Durée conseillée : #6}}
  \end{tcolorbox}
  \begin{tcolorbox}[enhanced,colback=white,colframe=popink,boxrule=2.2pt,arc=10pt,
      drop shadow=popink,left=16pt,right=16pt,top=10pt,bottom=10pt,after skip=8pt]
    \poppill{Dans cette leçon}{poppurple}{white}\par\vspace{5pt}
    \popsemi\fontsize{9.3}{12.6}\selectfont\color{popink}#7
  \end{tcolorbox}
  \noindent\begin{minipage}[t]{0.487\linewidth}
    \begin{tcolorbox}[enhanced,colback=popgreen,colframe=popink,boxrule=2.2pt,arc=10pt,
        drop shadow=popink,left=11pt,right=11pt,top=10pt,bottom=10pt,height=3.1cm,valign=center]
      \tcbox[colback=white,colframe=popink,boxrule=1.3pt,arc=5pt,boxsep=0pt,nobeforeafter,
        left=3pt,right=3pt,top=3pt,bottom=3pt,valign=center]{\qrcode[height=1.9cm]{https://play.google.com/store/apps/details?id=com.luvvix.onbuch&hl=fr}}%
      \hspace{8pt}\begin{minipage}[c]{0.5\linewidth}
        {\popdisplay\fontsize{11}{12.5}\selectfont\color{white}\MakeUppercase{Télécharge OnBuch}}\par\vspace{4pt}
        {\raggedright\popsemi\fontsize{8.4}{11}\selectfont\color{white}Gratuit \textbullet\ Google Play\\Tuteur IA, annales, concours, résultats\par}
      \end{minipage}
    \end{tcolorbox}
  \end{minipage}\hfill
  \begin{minipage}[t]{0.487\linewidth}
    \begin{tcolorbox}[enhanced,colback=poppurple,colframe=popink,boxrule=2.2pt,arc=10pt,
        drop shadow=popink,left=11pt,right=11pt,top=10pt,bottom=10pt,height=3.1cm,valign=center]
      \tikz[baseline=-0.7ex]\node[circle,fill=white,draw=popink,line width=1.3pt,minimum size=1.6cm,inner sep=0]{\popdisplay\fontsize{24}{24}\selectfont\color{poppurple}+};%
      \hspace{8pt}\begin{minipage}[c]{0.6\linewidth}
        {\popdisplay\fontsize{11}{12.5}\selectfont\color{white}\MakeUppercase{Passe à OnBuch+}}\par\vspace{4pt}
        {\raggedright\popsemi\fontsize{8.4}{11}\selectfont\color{white}Tous les cours signés, Tuteur IA illimité, fiches PDF — onglet OnBuch+ de l'app\par}
      \end{minipage}
    \end{tcolorbox}
  \end{minipage}
  \vspace{8pt}
  \popcontenu
  \vfill
  \signatureonbuch
  \restoregeometry
  \newpage
}

% Rangée « Ce cours contient » (page de garde)
\newcommand{\poptile}[3]{% #1 couleur #2 gros texte #3 légende
  \begin{tcolorbox}[enhanced,colback=#1,colframe=popink,boxrule=2pt,arc=9pt,shadow={2.5pt}{-2.5pt}{0pt}{popink},
      left=6pt,right=6pt,top=9pt,bottom=9pt,halign=center,valign=center,height=1.75cm,width=\linewidth,nobeforeafter]
    {\popdisplay\fontsize{12.5}{13}\selectfont\color{white}\MakeUppercase{#2}}\par\vspace{3pt}
    {\centering\popsemi\fontsize{7.8}{9.5}\selectfont\color{white}#3\par}
  \end{tcolorbox}}
\newcommand{\popcontenu}{%
  \par\noindent{\popxboldls\fontsize{7.5}{7.5}\selectfont\color{popmuted}CE COURS SIGNÉ CONTIENT}\par\vspace{7pt}
  \noindent\begin{minipage}[t]{0.235\linewidth}\poptile{popblue}{Cours}{complet, illustré, pas à pas}\end{minipage}\hfill
  \begin{minipage}[t]{0.235\linewidth}\poptile{poporange}{Méthodes}{et exercices résolus}\end{minipage}\hfill
  \begin{minipage}[t]{0.235\linewidth}\poptile{poppink}{Exercices}{type examen + corrigés}\end{minipage}\hfill
  \begin{minipage}[t]{0.235\linewidth}\poptile{popgreen}{Fiche bilan}{l'essentiel à retenir}\end{minipage}\par}

% Sommaire
\newtcolorbox{sommaire}{enhanced,colback=white,colframe=popink,boxrule=2.2pt,arc=10pt,
  drop shadow=popink,left=18pt,right=18pt,top=13pt,bottom=13pt,before skip=6pt,after skip=20pt,
  before upper={\poppill{Sommaire}{poppurple}{white}\par\vspace{9pt}\popsemi\fontsize{10.6}{14.5}\selectfont\color{popink}}}

% Signature de fin de cours — poussée tout en bas de la dernière page
\newcommand{\findecours}{%
  \par\vfill\nopagebreak
  \begin{center}\popsquiggle{poporange}\end{center}\vspace{4pt}
  \signatureonbuch
}
\AtBeginDocument{\def\llbracket{\mathopen{[\mkern-3.5mu[}}\def\rrbracket{\mathclose{]\mkern-3.5mu]}}} % intervalles d'entiers (glyphe ⟦ absent de la police)
% Glyphes absents de Fira Math : redéfinis à partir de glyphes disponibles
\AtBeginDocument{%
  \def\setminus{\mathbin{\backslash}}\let\smallsetminus\setminus
  \def\vdots{\mathord{\vbox{\baselineskip3.2pt\lineskiplimit0pt\kern4pt\hbox{.}\hbox{.}\hbox{.}}}}%
  \def\ddots{\mathinner{\mkern1mu\raise6.5pt\hbox{.}\mkern2mu\raise3.6pt\hbox{.}\mkern2mu\raise0.7pt\hbox{.}\mkern1mu}}%
  \def\triangle{\mathord{\scalebox{0.9}{$\symup{\Delta}$}}}%
  \def\nabla{\mathord{\raisebox{\depth}{\scalebox{1}[-1]{$\symup{\Delta}$}}}}%
  \def\checkmark{\popcheck{popdark}}%
  \def\star{\mathbin{\ast}}\def\bigstar{\mathord{\ast}}%
  \def\diamond{\mathbin{\scalebox{0.6}{$\Diamond$}}}%
  \def\bigcup{\mathop{\mathchoice{\raisebox{-0.3em}{\scalebox{1.7}{$\cup$}}}{\scalebox{1.25}{$\cup$}}{\cup}{\cup}}\displaylimits}%
  \def\bigcap{\mathop{\mathchoice{\raisebox{-0.3em}{\scalebox{1.7}{$\cap$}}}{\scalebox{1.25}{$\cap$}}{\cap}{\cap}}\displaylimits}%
  \def\lesseqqgtr{\mathrel{\lessgtr}}\def\bigoplus{\mathop{\oplus}\displaylimits}%
}
\providecommand{\ssi}{si et seulement si} % abréviation fréquente
\AtBeginDocument{\providecommand{\wideparen}{\overparen}} % arc de cercle

%% FIGFIX-BEGIN v5 — correctifs globaux des figures (docs/onbuch-generation/tools/figfix)
\usepackage{adjustbox}\usepackage{etoolbox}\tcbuselibrary{raster}
% toute figure TikZ/pgfplots plus large que la ligne est réduite à la largeur disponible
\BeforeBeginEnvironment{tikzpicture}{\begin{adjustbox}{max width=\linewidth}}
\AfterEndEnvironment{tikzpicture}{\end{adjustbox}}
% légendes pgfplots lisibles par-dessus les courbes
\pgfplotsset{every axis/.append style={legend style={fill=white,fill opacity=0.92,text opacity=1,draw=black!15,inner sep=3pt}}}
% étiquettes d'arêtes (syntaxe edge["..."]) avec fond blanc
\tikzset{every edge quotes/.append style={fill=white,inner sep=1.2pt}}
%% FIGFIX-END
\begin{document}

\pagegarde{Généralités sur les fonctions numériques}{Mathématiques}{1re Tech}{Mathématiques – classe de Première technique}{N12}{2 séances (fiche de progression harmonisée)}{Cette leçon consolide les notions de fonction numérique : opérations, composition, parité, périodicité, tracé de courbes et fonctions associées. Elle prépare l'élève à modéliser et résoudre des problèmes concrets du quotidien camerounais.
\vspace{4pt}
\begin{multicols}{2}\small
\begin{itemize}
\item À la fin de cette leçon, je sais définir une fonction, son ensemble de définition et son ensemble d'arrivée.
\item Je sais effectuer les opérations usuelles (somme, produit, quotient) et la composition de deux fonctions.
\item Je sais reconnaître une fonction paire ou impaire et en déduire les symétries de sa courbe.
\item Je sais identifier la période d'une fonction périodique et l'utiliser pour prolonger sa courbe.
\item Je sais tracer la courbe représentative d'une fonction en utilisant un tableau de valeurs, les asymptotes et les variations.
\item Je sais construire les fonctions associées : translation, homothétie, symétrie axiale et inverse (si elle existe).
\end{itemize}\end{multicols}}

\begin{sommaire}
\begin{multicols}{2}
\begin{enumerate}
\item 1. Rappels et vocabulaire sur les fonctions numériques
\item 2. Opérations sur les fonctions et composition
\item 3. Parité, symétrie axiale et centrale
\item 4. Périodicité
\item 5. Courbe d'une fonction : tracé et analyse
\item 6. Fonctions associées à une fonction donnée
\item Activité d'intégration
\item Méthodes et exercices résolus
\item Exercices
\item Corrigés des exercices
\item Fiche bilan
\end{enumerate}
\end{multicols}
\end{sommaire}

\begin{objectifs}
\begin{itemize}
\item À la fin de cette leçon, je sais définir une fonction, son ensemble de définition et son ensemble d'arrivée.
\item Je sais effectuer les opérations usuelles (somme, produit, quotient) et la composition de deux fonctions.
\item Je sais reconnaître une fonction paire ou impaire et en déduire les symétries de sa courbe.
\item Je sais identifier la période d'une fonction périodique et l'utiliser pour prolonger sa courbe.
\item Je sais tracer la courbe représentative d'une fonction en utilisant un tableau de valeurs, les asymptotes et les variations.
\item Je sais construire les fonctions associées : translation, homothétie, symétrie axiale et inverse (si elle existe).
\item Je sais rédiger une démonstration guidée (ex. : composition de deux fonctions paires est paire) et justifier chaque étape.
\item Je sais appliquer ces outils à une situation réelle (facture d'électricité, horaires de bus, etc.) et présenter ma démarche.
\end{itemize}
\end{objectifs}

\begin{prerequis}
\begin{itemize}
\item Notion de fonction (définition, image, antécédent) vue en 4ème.
\item Calcul littéral : factorisation, simplification de fractions rationnelles.
\item Repérage dans le plan : coordonnées, symétrie axiale et centrale.
\item Notion de limite et d'asymptote (abordée en début de Première).
\item Utilisation d'une calculatrice graphique ou d'un logiciel de tracé (GeoGebra, TikZ).
\end{itemize}
\end{prerequis}

\newpage

\coursec{Rappels et vocabulaire sur les fonctions numériques}

\courssub{Situation d'accroche et problématique}

Imaginez un ménage à Yaoundé qui reçoit sa facture d'électricité \textsc{eneo} chaque mois. Le montant à payer n'est pas un nombre au hasard : il dépend directement de la quantité d'énergie consommée (en kWh), elle‑même influencée par la saison — climatiseurs en saison sèche, éclairage plus long en saison des pluies. Le tarif appliqué compose une partie fixe (l'abonnement) et une partie proportionnelle à la consommation, avec parfois des tranches progressives. Mathématiquement, ce mécanisme est une \textbf{fonction} : à chaque consommation possible (l'entrée) correspond un unique montant à payer (la sortie).

Plus généralement, en technique (électricité, mécanique, gestion, génie civil), on modélise sans cesse des grandeurs qui dépendent les unes des autres : la tension aux bornes d'une résistance en fonction du courant (loi d'Ohm), la distance de freinage en fonction de la vitesse, le coût de production en fonction du nombre d'unités fabriquées. Maîtriser le vocabulaire et les outils de base des fonctions numériques, c'est se donner le langage commun pour décrire, analyser et prévoir ces situations concrètes.

Dans cette première section, nous posons les fondations rigoureuses : définition précise, notations, lecture graphique, et premiers exemples usuels que vous retrouverez tout au long de l'année.

\courssub{Définition rigoureuse : application, ensemble de définition, ensemble d'arrivée}

\begin{definition}[Fonction numérique d'une variable réelle]
Une \textbf{fonction numérique} $f$ est une application qui associe à \textbf{chaque} élément $x$ d'un ensemble $D$ (appelé \textbf{ensemble de définition} ou \textbf{domaine}) un \textbf{unique} élément $y$ d'un ensemble $A$ (appelé \textbf{ensemble d'arrivée}), ici toujours $\mathbb{R}$ (les nombres réels).
On note :
\[
f \,:\, D \longrightarrow \mathbb{R} \qquad\text{ou}\qquad \begin{array}{rcl} f : D & \to & \mathbb{R} \\ x & \mapsto & f(x) \end{array}
\]
L'élément $y = f(x)$ est appelé \textbf{image} de $x$ par $f$. Réciproquement, $x$ est un \textbf{antécédent} de $y$ par $f$.
L'\textbf{ensemble image} (ou \textbf{ensemble des valeurs}) de $f$ est l'ensemble $f(D) = \{ f(x) \mid x \in D \} \subset \mathbb{R}$.
\end{definition}

\begin{propriete}[Critère de fonctionnalité]
Une relation (ou un graphe dans le plan) définit une fonction si et seulement si \textbf{chaque} abscisse $x \in D$ possède \textbf{une et une seule} ordonnée $y$.
En termes ensemblistes : $\forall x \in D,\; \exists!\, y \in \mathbb{R}$ tel que $y = f(x)$.
\end{propriete}

\begin{attention}[Piège : ensemble de définition $\neq$ ensemble image]
L'ensemble de définition $D$ est l'ensemble des \textbf{entrées} autorisées (les $x$). L'ensemble image $f(D)$ est l'ensemble des \textbf{sorties} obtenues (les $y$).
\emph{Exemple :} $f(x) = x^2$ avec $D = \mathbb{R}$. Alors $f(D) = [0, +\infty[$. Ne confondez jamais les deux : $D$ se lit sur l'axe des abscisses, $f(D)$ sur l'axe des ordonnées.
\end{attention}

\begin{savaistu}[Histoire de la notation $f(x)$]
C'est le mathématicien suisse \textbf{Leonhard Euler} qui, en \textbf{1734}, a introduit la notation $f(x)$ pour désigner la valeur de la fonction $f$ au point $x$. Avant lui, on utilisait des notations géométriques lourdes ou des lettres majuscules sans parenthèse. Cette notation compacte a révolutionné l'écriture du calcul différentiel et intégral, et elle est aujourd'hui universelle — y compris dans les cahiers de chantier des ingénieurs camerounais !
\end{savaistu}

\courssub{Notations, vocabulaire et premiers exemples usuels}

\begin{exemplebox}[Fonction affine : tarif électrique simplifié]
Soit la fonction $f$ qui associe à une consommation $x$ (en kWh) le montant $y$ (en FCFA) à payer selon le tarif : abonnement fixe $500$ FCFA + $50$ FCFA par kWh.
\begin{itemize}
\item Ensemble de définition : $D = [0, +\infty[$ (on ne consomme pas une quantité négative).
\item Expression analytique : $f(x) = 50x + 500$.
\item Image de $10$ : $f(10) = 50 \times 10 + 500 = 1000$ FCFA.
\item Antécédent de $2000$ : on résout $50x + 500 = 2000 \Rightarrow 50x = 1500 \Rightarrow x = 30$ kWh.
\item Ensemble image : $f(D) = [500, +\infty[$.
\end{itemize}
C'est une \textbf{fonction affine} (forme $ax+b$ avec $a \neq 0$). Si $a=0$, c'est une fonction constante.
\end{exemplebox}

\begin{exemplebox}[Fonctions usuelles du programme]
Voici les cinq familles de fonctions de référence que vous devez connaître parfaitement (ensemble de définition $D$, ensemble image $f(D)$, allure de la courbe) :

\begin{center}
\begin{tabularx}{\linewidth}{|l|X|X|X|}\hline
\textbf{Nom} & \textbf{Expression $f(x)$} & \textbf{Ensemble de définition $D$} & \textbf{Ensemble image $f(D)$} \\ \hline
Affine & $ax+b$ ($a,b \in \mathbb{R}$) & $\mathbb{R}$ & $\mathbb{R}$ si $a \neq 0$ ; $\{b\}$ si $a=0$ \\ \hline
Quadratique (carré) & $x^2$ & $\mathbb{R}$ & $[0, +\infty[$ \\ \hline
Rationnelle (inverse) & $\dfrac{1}{x}$ & $\mathbb{R}^* = ]-\infty,0[ \cup ]0,+\infty[$ & $\mathbb{R}^*$ \\ \hline
Racine carrée & $\sqrt{x}$ & $[0, +\infty[$ & $[0, +\infty[$ \\ \hline
Valeur absolue & $|x|$ & $\mathbb{R}$ & $[0, +\infty[$ \\ \hline
\end{tabularx}
\end{center}
\end{exemplebox}

\begin{exoresolu}[Déterminer l'ensemble de définition et calculer des images]
Soit la fonction $f$ définie par $f(x) = \dfrac{2x+1}{x-3}$.
\begin{enumerate}
\item Déterminer l'ensemble de définition $D_f$.
\item Calculer $f(0)$, $f(2)$, $f(-1)$.
\item $4$ a-t-il un antécédent ? Si oui, le(s) trouver.
\end{enumerate}
\tcblower
\textbf{Solution détaillée.}
\begin{enumerate}
\item Le dénominateur ne doit pas s'annuler : $x-3 \neq 0 \iff x \neq 3$.
Donc $D_f = \mathbb{R} \setminus \{3\} = ]-\infty, 3[ \cup ]3, +\infty[$.
\item $f(0) = \dfrac{2\times0+1}{0-3} = -\dfrac{1}{3}$.
$f(2) = \dfrac{2\times2+1}{2-3} = \dfrac{5}{-1} = -5$.
$f(-1) = \dfrac{2\times(-1)+1}{-1-3} = \dfrac{-1}{-4} = \dfrac{1}{4}$.
\item On cherche $x \in D_f$ tel que $f(x) = 4$.
$\dfrac{2x+1}{x-3} = 4 \iff 2x+1 = 4(x-3) \iff 2x+1 = 4x-12 \iff 2x = 13 \iff x = 6,5$.
$6,5 \neq 3$, donc $x=6,5$ est bien dans $D_f$. \textbf{Oui, l'antécédent de $4$ est $6,5$.}
\end{enumerate}
\end{exoresolu}

\courssub{Représentation graphique : la courbe $C_f$ dans un repère orthonormé}

\begin{definition}[Courbe représentative]
Dans un repère orthonormé $(O;\vec{\imath},\vec{\jmath})$, la \textbf{courbe représentative} de la fonction $f$, notée $C_f$, est l'ensemble des points $M(x\,;\,y)$ tels que $x \in D$ et $y = f(x)$.
\[
C_f = \{ M(x\,;\,f(x)) \mid x \in D \}
\]
Chaque point $M$ a pour abscisse $x$ (l'entrée) et pour ordonnée $f(x)$ (la sortie).
\end{definition}

\begin{popfigure}
\begin{tikzpicture}[scale=0.85]
% Axes
\draw[->, thick] (-3.5,0) -- (3.5,0) node[right] {$x$};
\draw[->, thick] (0,-1.5) -- (0,4.5) node[above] {$y$};
% Graduations
\foreach \x in {-3,-2,-1,1,2,3} \draw (\x,0.05) -- (\x,-0.05) node[below] {$\x$};
\foreach \y in {-1,1,2,3,4} \draw (0.05,\y) -- (-0.05,\y) node[left] {$\y$};
\node[below left] at (0,0) {$O$};
% Courbe y = x^2
\draw[domain=-2.2:2.2, smooth, variable=\x, popcurve, thick] plot ({\x},{\x*\x});
% Point générique M(x, f(x))
\coordinate (M) at (1.5, 2.25);
\draw[dashed, popmuted] (M) -- (1.5,0) node[below] {$x$};
\draw[dashed, popmuted] (M) -- (0,2.25) node[left] {$f(x)$};
\fill[popblue] (M) circle (2.5pt) node[above right] {$M(x\,;\,f(x))$};
% Flèches d'orientation
\draw[poporange, <->, thick] (-2.5, 0.5) -- (-2.5, 3.5) node[midway, left] {$f(D)$};
\draw[popgreen, <->, thick] (0.5, -1) -- (2.5, -1) node[midway, below] {$D$};
\end{tikzpicture}
\legende{Courbe $C_f$ de la fonction $f(x)=x^2$. L'ensemble de définition $D$ se lit sur l'axe des abscisses, l'ensemble image $f(D)$ sur l'axe des ordonnées.}
\end{popfigure}

\begin{methode}[Lire une image et des antécédents sur la courbe $C_f$]
Donnée la courbe $C_f$ d'une fonction $f$ dans un repère orthonormé :
\begin{enumerate}
\item \textbf{Image d'une valeur $a$ (si $a \in D$) :}
Repérez $a$ sur l'axe des abscisses $\to$ montez (ou descendez) verticalement jusqu'à rencontrer la courbe $\to$ lisez l'ordonnée du point d'intersection : c'est $f(a)$.
\item \textbf{Antécédent(s) d'une valeur $b$ :}
Repérez $b$ sur l'axe des ordonnées $\to$ tracez l'horizontale $y = b$ $\to$ repérez \textbf{tous} les points d'intersection avec la courbe $\to$ lisez leurs abscisses : ce sont les antécédents de $b$.
\item \textbf{Ensemble de définition $D$ :}
C'est la projection de la courbe sur l'axe des abscisses (l'ombre de la courbe au sol).
\item \textbf{Ensemble image $f(D)$ :}
C'est la projection de la courbe sur l'axe des ordonnées (l'ombre de la courbe sur le mur).
\end{enumerate}
\end{methode}

\begin{exemplebox}[Lecture graphique sur $f(x) = x^2$]
Sur la courbe ci‑dessus ($f(x)=x^2$, $D=\mathbb{R}$) :
\begin{itemize}
\item Image de $1,5$ : on monte à $x=1,5$, on trouve $y=2,25$. Donc $f(1,5)=2,25$.
\item Antécédents de $4$ : l'horizontale $y=4$ coupe la courbe en deux points d'abscisses $-2$ et $2$. Donc $4$ a deux antécédents : $-2$ et $2$.
\item $D = \mathbb{R}$ (la courbe s'étend indéfiniment à gauche et à droite).
\item $f(D) = [0, +\infty[$ (la courbe ne descend jamais sous l'axe des abscisses).
\end{itemize}
\end{exemplebox}

\begin{attention}[Erreurs fréquentes en lecture graphique]
\begin{itemize}
\item \textbf{Confondre abscisse et ordonnée :} « L'image de $2$ est $4$ » $\neq$ « L'antécédent de $2$ est $4$ ».
\item \textbf{Oublier qu'un nombre peut avoir plusieurs antécédents} (ex. $x^2$ ou $|x|$) ou aucun (ex. $-1$ pour $x^2$).
\item \textbf{Lire l'ensemble image sur l'axe des abscisses} (ou inversément). Projetez toujours \emph{perpendiculairement} à l'axe concerné.
\item \textbf{Ne pas vérifier que la valeur demandée appartient à $D$} avant de chercher son image (ex. chercher $f(3)$ pour $1/(x-3)$).
\end{itemize}
\end{attention}

\courssub{Exercices d'application immédiate}

\begin{exercice}[Vocabulaire et ensembles]
Soit la fonction $g$ définie par $g(x) = \sqrt{4-x}$.
\begin{enumerate}
\item Déterminer l'ensemble de définition $D_g$.
\item Calculer $g(0)$, $g(-5)$, $g(4)$.
\item $2$ a-t-il un antécédent par $g$ ? $3$ a-t-il un antécédent ?
\item Quel est l'ensemble image $g(D_g)$ ?
\end{enumerate}
\end{exercice}

\begin{corrige}[Exercice 1]
\begin{enumerate}
\item La radicande doit être $\ge 0$ : $4-x \ge 0 \iff x \le 4$. Donc $D_g = ]-\infty, 4]$.
\item $g(0) = \sqrt{4} = 2$. $g(-5) = \sqrt{4-(-5)} = \sqrt{9} = 3$. $g(4) = \sqrt{0} = 0$.
\item Pour $2$ : $\sqrt{4-x} = 2 \iff 4-x = 4 \iff x = 0 \in D_g$. \textbf{Oui, antécédent : $0$}.
Pour $3$ : $\sqrt{4-x} = 3 \iff 4-x = 9 \iff x = -5 \in D_g$. \textbf{Oui, antécédent : $-5$}.
\item $g(x) = \sqrt{4-x} \ge 0$ pour tout $x \in D_g$. Quand $x \to -\infty$, $g(x) \to +\infty$. Quand $x=4$, $g(4)=0$. La fonction est continue et décroissante sur $D_g$. Donc $g(D_g) = [0, +\infty[$.
\end{enumerate}
\end{corrige}

\begin{exercice}[Lecture graphique]
On donne ci‑dessous la courbe $C_h$ d'une fonction $h$ définie sur $[-4, 4]$.
\begin{center}
\begin{tikzpicture}[scale=0.7]
\draw[->] (-5,0) -- (5,0) node[right]{$x$};
\draw[->] (0,-3) -- (0,4) node[above]{$y$};
\foreach \x in {-4,-3,-2,-1,1,2,3,4} \draw (\x,0.05) -- (\x,-0.05) node[below]{\small$\x$};
\foreach \y in {-2,-1,1,2,3} \draw (0.05,\y) -- (-0.05,\y) node[left]{\small$\y$};
\node[below left] at (0,0) {$O$};
% Courbe par morceaux : deux demi-paraboles
\draw[popcurve, thick, domain=-4:0, smooth, variable=\x] plot ({\x},{-0.5*(\x+2)*(\x+2)+3});
\draw[popcurve, thick, domain=0:4, smooth, variable=\x] plot ({\x},{0.5*(\x-2)*(\x-2)-1});
% Points caractéristiques
\fill[popblue] (-2,3) circle (2pt) node[above left] {$A(-2;3)$};
\fill[popblue] (2,-1) circle (2pt) node[below right] {$B(2;-1)$};
\fill[popblue] (0,1) circle (2pt) node[left] {$C(0;1)$};
\fill[popblue] (-4,1) circle (2pt) node[above left] {$D(-4;1)$};
\fill[popblue] (4,1) circle (2pt) node[below right] {$E(4;1)$};
\end{tikzpicture}
\end{center}
Répondez par vrai/faux et justifiez par une lecture graphique :
\begin{enumerate}
\item $h(-2) = 3$
\item $h(0) = -1$
\item L'antécédent de $1$ est $4$.
\item L'ensemble de définition est $[-4, 4]$.
\item L'ensemble image est $[-1, 3]$.
\item La fonction $h$ est injective (chaque valeur a au plus un antécédent).
\end{enumerate}
\end{exercice}

\begin{corrige}[Exercice 2]
\begin{enumerate}
\item \textbf{Vrai.} Le point $A(-2;3)$ appartient à la courbe.
\item \textbf{Faux.} Le point d'abscisse $0$ est $C(0;1)$, donc $h(0)=1$.
\item \textbf{Faux.} L'horizontale $y=1$ coupe la courbe en trois points : $C(0;1)$, $E(4;1)$ et $D(-4;1)$. Donc $1$ a \textbf{trois} antécédents : $-4$, $0$ et $4$.
\item \textbf{Vrai.} La courbe existe exactement pour $x \in [-4, 4]$.
\item \textbf{Vrai.} Les ordonnées varient de $-1$ (minimum en $B$) à $3$ (maximum en $A$).
\item \textbf{Faux.} La valeur $1$ a trois antécédents ($-4$, $0$, $4$), la valeur $0$ en a deux (entre $-2$ et $0$, et entre $0$ et $2$).
\end{enumerate}
\end{corrige}

\begin{savaistu}[Modélisation en génie civil : flèche d'une poutre]
En résistance des matériaux, la flèche (déformation verticale) $f(x)$ d'une poutre simplement appuyée de longueur $L$, soumise à une charge uniforme, suit une loi quartique : $f(x) = \frac{q}{24EI}(x^4 - 2Lx^3 + L^3x)$. Ici $x$ est la distance à l'appui gauche ($D=[0,L]$), $q$ la charge, $E$ le module de Young, $I$ le moment d'inertie. C'est une fonction polynomiale de degré $4$. L'ingénieur en cherche le maximum (milieu de portée) pour vérifier que la déformation reste dans les normes. Les outils de cette leçon (définition, image, courbe, extremum) sont exactement ceux qu'il utilise quotidiennement.
\end{savaistu}

\coursec{Opérations sur les fonctions et composition}

Dans la section précédente, nous avons rappelé ce qu'est une \cle{fonction numérique} : une machine qui, à un nombre $x$ de son ensemble de définition, associe un unique nombre $f(x)$. Nous savons maintenant lire une image, un antécédent, un ensemble de définition. Mais dans la vie professionnelle, on ne rencontre presque jamais une fonction toute seule. Un électricien calcule un prix de revient en \emph{additionnant} le coût du matériel et le coût de la main-d'œuvre ; un gestionnaire \emph{multiplie} un prix unitaire par une quantité ; un atelier \emph{enchaîne} deux transformations : d'abord la découpe, puis la soudure. Toutes ces situations reviennent à \textbf{combiner des fonctions entre elles}. C'est l'objet de cette section : apprendre à additionner, multiplier, diviser et surtout \emph{composer} des fonctions, en surveillant à chaque fois la question capitale : \textbf{pour quels nombres $x$ le calcul a-t-il un sens ?}

\courssub{Somme, différence, produit et quotient de deux fonctions}

\textbf{Intuition.} Imaginons un petit atelier de menuiserie à Yaoundé. Pour fabriquer $x$ tables, le coût du bois est donné par une fonction $f(x)$ et le coût de la main-d'œuvre par une fonction $g(x)$. Le coût total est naturellement $f(x)+g(x)$ : c'est une \emph{nouvelle} fonction, que l'on note $f+g$. De même, si le patron veut connaître le bénéfice, il calcule la différence entre la recette et le coût : $f-g$. Rien de mystérieux : on effectue, pour chaque valeur de $x$, l'opération sur les images.

\begin{definition}[Somme, différence, produit, quotient de deux fonctions]
Soient $f$ et $g$ deux fonctions numériques définies respectivement sur $D_f$ et $D_g$. On pose $D = D_f \cap D_g$ (l'ensemble des nombres réels où \emph{les deux} fonctions sont définies). On définit sur $D$ :
\begin{itemize}
\item la \cle{somme} : $(f+g)(x) = f(x) + g(x)$ ;
\item la \cle{différence} : $(f-g)(x) = f(x) - g(x)$ ;
\item le \cle{produit} : $(f \times g)(x) = f(x) \times g(x)$, noté aussi $(fg)(x)$ ;
\item le \cle{quotient} : pour tout $x \in D$ tel que $g(x) \neq 0$, $\left(\dfrac{f}{g}\right)(x) = \dfrac{f(x)}{g(x)}$.
\end{itemize}
\end{definition}

\begin{aretenir}[Règle d'or des ensembles de définition]
Pour la somme, la différence et le produit, l'ensemble de définition est $D_f \cap D_g$ : il faut que \textbf{les deux} calculs soient possibles. Pour le quotient, il faut \textbf{en plus} exclure les valeurs de $x$ qui annulent le dénominateur :
\[ D_{f/g} = \{\, x \in D_f \cap D_g \ \mid \ g(x) \neq 0 \,\}. \]
\end{aretenir}

\begin{exemplebox}[Coût total d'un atelier]
Soient $f(x) = 3x + 2$ (coût du bois, en milliers de francs CFA) et $g(x) = x + 5$ (coût de la main-d'œuvre), toutes deux définies sur $\mathbb{R}$. Alors :
\[ (f+g)(x) = (3x+2) + (x+5) = 4x + 7, \qquad (f \times g)(x) = (3x+2)(x+5) = 3x^2 + 17x + 10. \]
Pour $x = 10$ tables : $(f+g)(10) = 47$, soit un coût total de \SI{47000}{\text{F CFA}}.
\end{exemplebox}

\begin{attention}[Le piège du quotient]
L'erreur la plus fréquente au Probatoire est d'oublier la condition $g(x) \neq 0$ dans un quotient. Par exemple, si $f(x) = x+1$ et $g(x) = x - 3$, alors $\left(\dfrac{f}{g}\right)(x) = \dfrac{x+1}{x-3}$ n'est définie que pour $x \neq 3$, même si le numérateur, lui, ne pose aucun problème. Écris \emph{toujours} la condition avant de simplifier.
\end{attention}

\courssub{Composition de deux fonctions}

\textbf{Intuition.} La composition, c'est l'\emph{enchaînement} de deux machines : le résultat de la première devient l'entrée de la seconde. Prenons une situation de la vie courante. La consommation de gaz d'un restaurant dépend du nombre $x$ de repas servis : disons $g(x) = 0{,}5x$ bouteilles. Le montant de la facture dépend du nombre de bouteilles : disons $f(u) = 6500u$ francs. Pour connaître la facture directement en fonction du nombre de repas, on enchaîne : on calcule d'abord $g(x)$, puis on applique $f$ au résultat. On obtient $f(g(x)) = 6500 \times 0{,}5x = 3250x$ francs. Cette nouvelle fonction est la \cle{composée} de $g$ par $f$.

\begin{definition}[Fonction composée]
Soient $g$ une fonction définie sur $D_g$ et $f$ une fonction définie sur $D_f$. La \cle{composée} de $g$ par $f$, notée $f \circ g$ (lire « $f$ rond $g$ »), est la fonction définie par :
\[ (f \circ g)(x) = f\bigl(g(x)\bigr). \]
Son ensemble de définition est :
\[ D_{f \circ g} = \{\, x \in D_g \ \mid \ g(x) \in D_f \,\}. \]
Autrement dit : $x$ doit être dans l'ensemble de définition de $g$ (première machine), \textbf{et} le résultat $g(x)$ doit être dans l'ensemble de définition de $f$ (seconde machine).
\end{definition}

\begin{attention}[L'ordre compte !]
Dans $f \circ g$, c'est $g$ qui agit \textbf{en premier}, même si elle est écrite à droite. On lit le schéma de droite à gauche :
\[ x \xrightarrow{\ g\ } g(x) \xrightarrow{\ f\ } f\bigl(g(x)\bigr). \]
En général, $f \circ g \neq g \circ f$ : la composition \textbf{n'est pas commutative}. Par exemple, avec $f(x) = x^2$ et $g(x) = x + 1$ :
\[ (f \circ g)(x) = (x+1)^2 = x^2 + 2x + 1 \qquad \text{alors que} \qquad (g \circ f)(x) = x^2 + 1. \]
Les deux résultats sont différents : ne jamais intervertir l'ordre sans vérification.
\end{attention}

\begin{methode}[Déterminer l'ensemble de définition d'une composée $f \circ g$]
\begin{enumerate}
\item Écrire les ensembles de définition $D_g$ et $D_f$ séparément.
\item Poser la première condition : $x \in D_g$.
\item Poser la seconde condition : $g(x) \in D_f$, puis la traduire par une équation ou une inéquation en $x$ (par exemple, si $f(u) = \sqrt{u}$, exiger $g(x) \geq 0$ ; si $f(u) = \dfrac{1}{u}$, exiger $g(x) \neq 0$).
\item Résoudre, puis prendre l'\textbf{intersection} des deux conditions : c'est $D_{f \circ g}$.
\item Calculer enfin l'expression simplifiée de $(f \circ g)(x)$.
\end{enumerate}
\end{methode}

\begin{exoresolu}[Composée avec une racine carrée]
Soient $f(x) = \sqrt{x}$ et $g(x) = x^2 + 1$. Déterminer $(f \circ g)(x)$ et son ensemble de définition, puis $(g \circ f)(x)$ et son ensemble de définition.
\tcblower
\textbf{Étape 1 : ensembles de définition.} $D_f = [0\,;\,+\infty[$ (la racine carrée n'existe que pour un nombre positif ou nul) et $D_g = \mathbb{R}$.

\textbf{Étape 2 : calcul de $f \circ g$.}
\[ (f \circ g)(x) = f\bigl(g(x)\bigr) = f(x^2+1) = \sqrt{x^2+1}. \]
Conditions : $x \in D_g = \mathbb{R}$ (toujours vrai) et $g(x) = x^2 + 1 \in D_f$, c'est-à-dire $x^2 + 1 \geq 0$. Or pour tout réel $x$, $x^2 \geq 0$ donc $x^2 + 1 \geq 1 > 0$ : la condition est toujours satisfaite. Conclusion :
\[ D_{f \circ g} = \mathbb{R} \qquad \text{et} \qquad (f \circ g)(x) = \sqrt{x^2+1}. \]

\textbf{Étape 3 : calcul de $g \circ f$.}
\[ (g \circ f)(x) = g\bigl(f(x)\bigr) = g(\sqrt{x}) = (\sqrt{x})^2 + 1 = x + 1. \]
Attention au piège : même si $x+1$ est défini pour tout réel, l'ensemble de définition se détermine \emph{avant} simplification ! Il faut $x \in D_f = [0\,;\,+\infty[$ et $f(x) = \sqrt{x} \in D_g = \mathbb{R}$ (toujours vrai). Donc :
\[ D_{g \circ f} = [0\,;\,+\infty[. \]
On retrouve bien que $f \circ g \neq g \circ f$ : les expressions \emph{et} les ensembles de définition diffèrent.
\end{exoresolu}

\begin{attention}[Simplifier ne change pas l'ensemble de définition]
Dans l'exercice précédent, $(g \circ f)(x) = x + 1$ semble défini partout, mais $D_{g \circ f} = [0\,;\,+\infty[$ à cause de la racine carrée de départ. Règle absolue : on détermine l'ensemble de définition \textbf{à partir de l'écriture initiale}, jamais après simplification.
\end{attention}

\begin{propriete}[Associativité de la composition]
Soient $h$, $g$, $f$ trois fonctions telles que les compositions soient définies. Alors :
\[ (f \circ g) \circ h = f \circ (g \circ h). \]
On peut donc écrire sans ambiguïté $f \circ g \circ h$ : c'est la fonction qui applique d'abord $h$, puis $g$, puis $f$.
\end{propriete}

\begin{experience}[Démonstration guidée de l'associativité]
Montrons que pour tout $x$ où les deux membres sont définis, $\bigl((f \circ g) \circ h\bigr)(x) = \bigl(f \circ (g \circ h)\bigr)(x)$.
\begin{enumerate}
\item Par définition de la composée de $h$ par $(f \circ g)$ : $\bigl((f \circ g) \circ h\bigr)(x) = (f \circ g)\bigl(h(x)\bigr)$.
\item Par définition de $f \circ g$ appliquée au nombre $h(x)$ : $(f \circ g)\bigl(h(x)\bigr) = f\bigl(g(h(x))\bigr)$.
\item De l'autre côté : $\bigl(f \circ (g \circ h)\bigr)(x) = f\bigl((g \circ h)(x)\bigr) = f\bigl(g(h(x))\bigr)$.
\item Les deux expressions sont identiques : l'égalité est démontrée.
\end{enumerate}
Cette démonstration est formatrice : elle montre que l'associativité n'est qu'une question de \emph{parenthésage}, l'ordre d'application ($h$, puis $g$, puis $f$) restant le même.
\end{experience}

\courssub{Exemple concret : tarification et consommation}

Dans les métiers de la gestion et de l'énergie, la composition est un outil quotidien. Une entreprise de distribution d'eau facture selon un tarif dégressif : le prix à payer est une fonction $f$ du volume $u$ consommé (en m$^3$). Mais le volume consommé dépend lui-même du nombre $x$ de jours d'activité d'un atelier, via une fonction $g$. La facture en fonction du nombre de jours est alors $(f \circ g)(x)$.

\begin{exoresolu}[Facture d'eau d'un atelier]
Un atelier de Briqueterie à Yaoundé consomme $g(x) = 2x + 3$ m$^3$ d'eau pour $x$ jours de travail. Le fournisseur facture $f(u) = 400u + 1500$ francs pour $u$ m$^3$ (1500 F de frais fixes).
\begin{enumerate}
\item Déterminer la fonction donnant la facture en fonction du nombre de jours.
\item Calculer la facture pour 12 jours de travail.
\end{enumerate}
\tcblower
\textbf{1.} Les deux fonctions sont définies sur $\mathbb{R}$, et pour $x \geq 0$ on a $g(x) \geq 0$, ce qui est cohérent avec le contexte. La facture en fonction du nombre de jours est :
\[ (f \circ g)(x) = f(2x+3) = 400(2x+3) + 1500 = 800x + 1200 + 1500 = 800x + 2700. \]
\textbf{2.} Pour $x = 12$ jours : $(f \circ g)(12) = 800 \times 12 + 2700 = 9600 + 2700 = \SI{12300}{\text{F CFA}}$.

Vérification par étapes : en 12 jours, l'atelier consomme $g(12) = 27$ m$^3$ ; la facture est $f(27) = 400 \times 27 + 1500 = 10800 + 1500 = \SI{12300}{\text{F CFA}}$. On retrouve le même résultat : la composition enchaîne bien les deux calculs.
\end{exoresolu}

\begin{popfigure}
\begin{center}
\begin{tikzpicture}[scale=0.7]
\draw[->] (-0.5,0)--(3.5,0) node[right]{$x$};
\draw[->] (0,-0.5)--(0,10) node[above]{$y$};
\draw[domain=0:3,smooth,variable=\x,red] plot ({\x},{\x*\x});
\draw[domain=0:1.5,smooth,variable=\x,blue] plot ({\x},{2*\x+1});
\draw[domain=0:1.3,smooth,variable=\x,popgreen,thick] plot ({\x},{(2*\x+1)*(2*\x+1)});
\node[red] at (2.5,8) {$f(x)=x^2$};
\node[blue] at (2.5,3) {$g(x)=2x+1$};
\node[popgreen] at (1.5,8) {$(f\circ g)(x)=(2x+1)^2$};
\end{tikzpicture}
\end{center}
\legende{La composée $f \circ g$ (en vert) s'obtient en appliquant d'abord $g$ (bleu), puis $f$ (rouge) : pour chaque $x$, on reporte $g(x)$ sur la courbe de $f$.}
\end{popfigure}

\begin{savaistu}[Pourquoi dit-on « f rond g » ?]
La notation $f \circ g$ se lit « $f$ \emph{rond} $g$ », à cause du petit rond $\circ$ placé entre les deux lettres. Ce symbole a été popularisé au XX\textsuperscript{e} siècle pour bien distinguer la composition du produit $f \times g$. Retiens l'ordre de lecture : dans $f \circ g$, on applique $g$ \emph{d'abord} — un peu comme dans l'expression française « prendre le bus \emph{puis} marcher » : l'action la plus proche de $x$ est celle qui s'effectue en premier. Au Cameroun, les ingénieurs en télécommunications utilisent constamment ce principe : le signal subit une première transformation (codage), puis une seconde (amplification) — c'est exactement une composition de fonctions !
\end{savaistu}

\begin{exercice}[À toi de jouer]
Soient $g(x) = \dfrac{1}{x-2}$ et $f(x) = x^2 - 3$.
\begin{enumerate}
\item Déterminer l'ensemble de définition de $f \circ g$ et calculer $(f \circ g)(x)$.
\item Déterminer l'ensemble de définition de $g \circ f$ et calculer $(g \circ f)(x)$.
\item Comparer les deux résultats : que constates-tu ?
\end{enumerate}
\end{exercice}

\begin{corrige}[Exercice 1]
\textbf{1.} $D_g = \mathbb{R} \setminus \{2\}$ et $D_f = \mathbb{R}$. Pour $f \circ g$ : il faut $x \neq 2$ et $g(x) \in \mathbb{R}$ (toujours vrai). Donc $D_{f \circ g} = \mathbb{R} \setminus \{2\}$ et
\[ (f \circ g)(x) = \left(\frac{1}{x-2}\right)^2 - 3 = \frac{1}{(x-2)^2} - 3. \]
\textbf{2.} Pour $g \circ f$ : il faut $x \in D_f = \mathbb{R}$ et $f(x) \in D_g$, c'est-à-dire $f(x) \neq 2$, soit $x^2 - 3 \neq 2$, soit $x^2 \neq 5$, soit $x \neq \sqrt{5}$ et $x \neq -\sqrt{5}$. Donc $D_{g \circ f} = \mathbb{R} \setminus \{-\sqrt{5}\,;\,\sqrt{5}\}$ et
\[ (g \circ f)(x) = \frac{1}{(x^2-3)-2} = \frac{1}{x^2-5}. \]
\textbf{3.} Les deux composées n'ont ni la même expression ni le même ensemble de définition : on confirme que $f \circ g \neq g \circ f$ en général.
\end{corrige}

\begin{aretenir}[L'essentiel de la section]
\begin{itemize}
\item Somme, différence, produit : définis sur $D_f \cap D_g$ ; quotient : exclure en plus les zéros du dénominateur.
\item $(f \circ g)(x) = f(g(x))$ : $g$ agit en premier ; $D_{f \circ g} = \{x \in D_g \mid g(x) \in D_f\}$.
\item La composition est \textbf{associative} : $(f \circ g) \circ h = f \circ (g \circ h)$, mais \textbf{non commutative} en général.
\item L'ensemble de définition se détermine \emph{avant} toute simplification de l'expression.
\end{itemize}
Ces outils seront indispensables dans la suite de la leçon : la parité, la périodicité et l'étude des courbes reposent toutes sur la manipulation des expressions $f(-x)$, $f(x+T)$, qui sont précisément des compositions.
\end{aretenir}

\coursec{Parité, symétrie axiale et centrale}

Dans la section précédente, nous avons appris à combiner des fonctions entre elles (somme, produit, composition). Nous allons maintenant étudier une propriété remarquable que possèdent certaines fonctions : la \cle{parité}. Cette propriété, lorsqu'elle existe, se lit directement sur la courbe sous forme d'une \cle{symétrie}. Connaître la parité d'une fonction avant de tracer sa courbe est un gain de temps considérable : il suffit d'étudier la fonction sur la moitié de son ensemble de définition, puis de compléter par symétrie. Au Probatoire technique, savoir tester la parité d'une fonction est une compétence régulièrement évaluée.

\courssub{Fonction paire}

\textbf{Intuition.} Regardez la courbe de la fonction $f(x)=x^2$. Les points d'abscisses $2$ et $-2$ ont la même image : $f(2)=4$ et $f(-2)=4$. La courbe se « replie » sur elle-même le long de l'axe des ordonnées, comme une feuille de papier pliée en deux. C'est exactement le phénomène que traduit la définition suivante.

\begin{definition}[Fonction paire]
Soit $f$ une fonction numérique définie sur un ensemble $D_f$ \cle{symétrique par rapport à $0$} (c'est-à-dire : si $x\in D_f$, alors $-x\in D_f$).

On dit que $f$ est \cle{paire} lorsque, pour tout $x\in D_f$ :
\[ f(-x)=f(x). \]
\end{definition}

\begin{attention}[Condition préalable souvent oubliée]
Avant même de calculer $f(-x)$, il faut vérifier que l'ensemble de définition est \textbf{symétrique par rapport à $0$}. Par exemple, la fonction $f(x)=x^2$ définie sur $[-1\,;\,2]$ n'est \textbf{pas} paire, car $-2\notin[-1\,;\,2]$ alors que $2\in[-1\,;\,2]$. Cette vérification est exigée dans la rédaction au Probatoire.
\end{attention}

\begin{propriete}[Interprétation graphique d'une fonction paire]
Une fonction $f$ est paire si et seulement si sa courbe représentative $\mathcal{C}_f$ est \cle{symétrique par rapport à l'axe des ordonnées} $(Oy)$.
\end{propriete}

\textbf{Justification.} Le point $M(x\,;\,f(x))$ appartient à $\mathcal{C}_f$. Son symétrique par rapport à l'axe $(Oy)$ est le point $M'(-x\,;\,f(x))$. Or $f$ paire signifie $f(-x)=f(x)$, donc $M'=(-x\,;\,f(-x))$ appartient aussi à $\mathcal{C}_f$. Réciproquement, si la courbe est symétrique par rapport à $(Oy)$, alors pour tout $x$, le point $(-x\,;\,f(x))$ est sur la courbe, donc $f(-x)=f(x)$.

\begin{exemplebox}[La fonction cosinus est paire]
Pour tout réel $x$, on a $\cos(-x)=\cos x$. La fonction cosinus est donc paire : sa courbe est symétrique par rapport à l'axe des ordonnées.
\end{exemplebox}

\begin{popfigure}
\begin{tikzpicture}[scale=0.8]
\draw[->] (-3,0)--(3,0) node[right]{$x$};
\draw[->] (0,-2)--(0,2) node[above]{$y$};
\draw[domain=-2.5:2.5,smooth,variable=\x,popblue,thick] plot ({\x},{cos(\x r)});
\draw[dashed,popmuted] (-2.5,0)--(2.5,0);
\draw[dashed,poporange] (1.5,0)--(1.5,{cos(1.5 r)}) node[above right,popdark]{\scriptsize $M(x\,;\,\cos x)$};
\draw[dashed,poporange] (-1.5,0)--(-1.5,{cos(1.5 r)}) node[above left,popdark]{\scriptsize $M'(-x\,;\,\cos x)$};
\fill[poporange] (1.5,{cos(1.5 r)}) circle (2pt);
\fill[poporange] (-1.5,{cos(1.5 r)}) circle (2pt);
\node[popblue] at (2.8,-1) {$y=\cos x$};
\end{tikzpicture}
\legende{La fonction cosinus est paire : les points $M$ et $M'$ sont symétriques par rapport à l'axe $(Oy)$.}
\end{popfigure}

\courssub{Fonction impaire}

\textbf{Intuition.} Considérons maintenant $f(x)=x^3$. On a $f(2)=8$ et $f(-2)=-8$ : les images sont \emph{opposées}. La courbe se superpose à elle-même après un demi-tour autour de l'origine. C'est la symétrie centrale de centre $O$.

\begin{definition}[Fonction impaire]
Soit $f$ une fonction numérique définie sur un ensemble $D_f$ symétrique par rapport à $0$.

On dit que $f$ est \cle{impaire} lorsque, pour tout $x\in D_f$ :
\[ f(-x)=-f(x). \]
\end{definition}

\begin{propriete}[Interprétation graphique d'une fonction impaire]
Une fonction $f$ est impaire si et seulement si sa courbe représentative $\mathcal{C}_f$ est \cle{symétrique par rapport à l'origine} $O$ du repère.
\end{propriete}

\textbf{Justification.} Le symétrique du point $M(x\,;\,f(x))$ par rapport à l'origine est le point $M'(-x\,;\,-f(x))$. Si $f$ est impaire, $f(-x)=-f(x)$, donc $M'=(-x\,;\,f(-x))$ appartient à $\mathcal{C}_f$. La réciproque se montre de la même façon.

\begin{exemplebox}[La fonction sinus est impaire]
Pour tout réel $x$, on a $\sin(-x)=-\sin x$. La fonction sinus est donc impaire : sa courbe est symétrique par rapport à l'origine.
\end{exemplebox}

\begin{popfigure}
\begin{tikzpicture}[scale=0.8]
\draw[->] (-3.5,0)--(3.5,0) node[right]{$x$};
\draw[->] (0,-1.8)--(0,1.8) node[above]{$y$};
\draw[domain=-3.14:3.14,smooth,variable=\x,popgreen,thick] plot ({\x},{sin(\x r)});
\draw[dashed,poporange] (2,0)--(2,{sin(2 r)}) node[above right,popdark]{\scriptsize $M(x\,;\,\sin x)$};
\draw[dashed,poporange] (-2,0)--(-2,{sin(-2 r)}) node[below left,popdark]{\scriptsize $M'(-x\,;\,-\sin x)$};
\fill[poporange] (2,{sin(2 r)}) circle (2pt);
\fill[poporange] (-2,{sin(-2 r)}) circle (2pt);
\draw[popmuted,dashed] (2,{sin(2 r)})--(-2,{sin(-2 r)});
\fill[popdark] (0,0) circle (1.5pt) node[below left]{\scriptsize $O$};
\node[popgreen] at (3,-1.3) {$y=\sin x$};
\end{tikzpicture}
\legende{La fonction sinus est impaire : $M$ et $M'$ sont symétriques par rapport à l'origine $O$.}
\end{popfigure}

\courssub{Méthode : tester la parité d'une fonction}

\begin{methode}[Étudier la parité d'une fonction]
Pour déterminer si une fonction $f$ est paire, impaire, ou ni l'une ni l'autre :
\begin{enumerate}
\item \textbf{Déterminer} l'ensemble de définition $D_f$ et \textbf{vérifier} qu'il est symétrique par rapport à $0$. Si ce n'est pas le cas, la fonction n'est ni paire ni impaire : on s'arrête là.
\item \textbf{Calculer} $f(-x)$ pour un réel $x$ quelconque de $D_f$, en simplifiant l'expression obtenue.
\item \textbf{Comparer} $f(-x)$ avec $f(x)$ et $-f(x)$ :
\begin{itemize}
\item si $f(-x)=f(x)$ pour tout $x$, alors $f$ est \textbf{paire} ;
\item si $f(-x)=-f(x)$ pour tout $x$, alors $f$ est \textbf{impaire} ;
\item sinon, $f$ n'est \textbf{ni paire ni impaire}.
\end{itemize}
\item \textbf{Conclure} en précisant la conséquence graphique (symétrie par rapport à $(Oy)$ ou à $O$).
\end{enumerate}
\end{methode}

\begin{exoresolu}[Parité d'une fonction polynomiale]
Énoncé. Soit $f$ la fonction définie sur $\mathbb{R}$ par $f(x)=x^3-2x$. Étudier la parité de $f$ et en déduire un élément de symétrie de sa courbe.
\tcblower
\textbf{Solution détaillée.}

\textbf{Étape 1.} $f$ est une fonction polynomiale, donc $D_f=\mathbb{R}$, qui est bien symétrique par rapport à $0$.

\textbf{Étape 2.} Calculons $f(-x)$ pour un réel $x$ quelconque :
\begin{align*}
f(-x) &= (-x)^3-2(-x) \\
&= -x^3+2x \\
&= -(x^3-2x) \\
&= -f(x).
\end{align*}

\textbf{Étape 3.} Pour tout réel $x$, $f(-x)=-f(x)$ : la fonction $f$ est \textbf{impaire}.

\textbf{Étape 4.} Conséquence graphique : la courbe de $f$ est symétrique par rapport à l'origine $O$ du repère. Il suffit donc de l'étudier sur $[0\,;\,+\infty[$, puis de compléter par symétrie centrale.
\end{exoresolu}

\begin{attention}[Une fonction n'est pas forcément paire ou impaire]
La plupart des fonctions ne sont \textbf{ni paires ni impaires}. Par exemple, pour $f(x)=x+1$ définie sur $\mathbb{R}$ :
\[ f(-x)=-x+1, \]
qui n'est égal ni à $f(x)=x+1$, ni à $-f(x)=-x-1$. Pour le prouver rapidement, un \textbf{contre-exemple} suffit : $f(1)=2$ et $f(-1)=0$ ; comme $f(-1)\neq f(1)$ et $f(-1)\neq -f(1)$, $f$ n'est ni paire ni impaire. Sa courbe n'a donc ni axe $(Oy)$ ni origine comme élément de symétrie.

Autre piège : ne concluez jamais à partir d'\emph{un seul} calcul. Montrer que $f(-2)=f(2)$ ne prouve pas que $f$ est paire ; l'égalité doit être démontrée pour \emph{tout} $x$.
\end{attention}

\courssub{Exemples classiques à connaître}

\begin{propriete}[Parité des fonctions puissances]
Soit $n$ un entier naturel non nul et $f(x)=x^n$.
\begin{itemize}
\item Si $n$ est \textbf{pair} ($n=2,4,6,\dots$), alors $f$ est \textbf{paire}.
\item Si $n$ est \textbf{impair} ($n=1,3,5,\dots$), alors $f$ est \textbf{impaire}.
\end{itemize}
De même, les fonctions trigonométriques de référence vérifient : $x\mapsto\cos x$ est \textbf{paire}, $x\mapsto\sin x$ est \textbf{impaire}.
\end{propriete}

\textbf{Démonstration (formatrice).} Pour $n$ pair, on peut écrire $n=2k$ avec $k\in\mathbb{N}$, d'où :
\[ f(-x)=(-x)^{2k}=\bigl((-x)^2\bigr)^k=(x^2)^k=x^{2k}=f(x). \]
Pour $n$ impair, on écrit $n=2k+1$ :
\[ f(-x)=(-x)^{2k+1}=(-x)^{2k}\times(-x)=x^{2k}\times(-x)=-x^{2k+1}=-f(x). \]
C'est d'ailleurs de là que viennent les mots « paire » et « impaire » : ils font référence à la parité de l'exposant des monômes.

\begin{center}
\begin{tabularx}{\linewidth}{|l|X|X|X|}
\hline
\rowcolor{popblueL} \textbf{Fonction} & \textbf{Parité} & \textbf{Symétrie de la courbe} & \textbf{Exemple} \\
\hline
$x^n$ avec $n$ pair & paire & axe $(Oy)$ & $f(x)=x^2$, $g(x)=x^4$ \\
\hline
$x^n$ avec $n$ impair & impaire & origine $O$ & $f(x)=x^3$, $g(x)=x$ \\
\hline
$\cos x$ & paire & axe $(Oy)$ & $\cos(-x)=\cos x$ \\
\hline
$\sin x$ & impaire & origine $O$ & $\sin(-x)=-\sin x$ \\
\hline
$\dfrac{1}{x}$ & impaire & origine $O$ & $f(-x)=-\dfrac{1}{x}=-f(x)$ \\
\hline
\end{tabularx}
\end{center}

\courssub{Opérations et parité}

Les règles suivantes, très pratiques, évitent de refaire le calcul de $f(-x)$ dans de nombreux cas. Elles ressemblent aux règles de calcul sur les nombres pairs et impairs.

\begin{propriete}[Parité et opérations sur les fonctions]
Soient $f$ et $g$ deux fonctions définies sur un même ensemble symétrique par rapport à $0$.
\begin{itemize}
\item La \textbf{somme} de deux fonctions paires est paire ; la somme de deux fonctions impaires est impaire.
\item Le \textbf{produit} de deux fonctions de \emph{même} parité (paire $\times$ paire, ou impaire $\times$ impaire) est \textbf{pair}.
\item Le \textbf{produit} d'une fonction paire par une fonction impaire est \textbf{impair}.
\item Si $f$ est paire (resp. impaire) et ne s'annule pas, alors $\dfrac{1}{f}$ est paire (resp. impaire).
\end{itemize}
\end{propriete}

\textbf{Démonstration d'un cas (formatrice).} Montrons que le produit d'une fonction paire $f$ et d'une fonction impaire $g$ est impair. Posons $h=f\times g$. Pour tout $x$ :
\[ h(-x)=f(-x)\times g(-x)=f(x)\times\bigl(-g(x)\bigr)=-f(x)g(x)=-h(x). \]
Donc $h$ est impaire. Les autres cas se démontrent de la même manière.

\begin{exemplebox}[Application directe]
\begin{itemize}
\item $f(x)=x^2+\cos x$ est \textbf{paire} (somme de deux fonctions paires).
\item $g(x)=x^3+\sin x$ est \textbf{impaire} (somme de deux fonctions impaires).
\item $h(x)=x^2\sin x$ est \textbf{impaire} (paire $\times$ impaire).
\item $k(x)=x\sin x$ est \textbf{paire} (impaire $\times$ impaire).
\end{itemize}
\end{exemplebox}

\begin{exoresolu}[Utiliser les règles de parité]
Énoncé. Sans calculer $h(-x)$ en détail, déterminer la parité de la fonction $h$ définie sur $\mathbb{R}$ par $h(x)=x^3\cos x - 5\sin x$.
\tcblower
\textbf{Solution détaillée.}

La fonction $x\mapsto x^3$ est impaire et $x\mapsto\cos x$ est paire : leur produit $x\mapsto x^3\cos x$ est donc \textbf{impair} (impaire $\times$ paire).

La fonction $x\mapsto 5\sin x$ est impaire (produit d'une constante, fonction paire, par une fonction impaire).

$h$ est donc la somme de deux fonctions impaires : $h$ est \textbf{impaire}.

Vérification par le calcul :
\begin{align*}
h(-x) &= (-x)^3\cos(-x)-5\sin(-x) \\
&= -x^3\cos x +5\sin x \\
&= -(x^3\cos x - 5\sin x) = -h(x).
\end{align*}
La courbe de $h$ est symétrique par rapport à l'origine.
\end{exoresolu}

\begin{attention}[La somme d'une paire et d'une impaire]
La somme d'une fonction paire et d'une fonction impaire n'est, en général, \textbf{ni paire ni impaire}. Par exemple $f(x)=x^2+x$ : $f(-x)=x^2-x$, qui n'est égal ni à $f(x)$ ni à $-f(x)$. Ne généralisez pas abusivement les règles du produit à la somme.
\end{attention}

\begin{savaistu}[La parité en électricité]
En électrotechnique, les signaux alternatifs sont modélisés par des fonctions sinusoïdales. Décomposer un signal périodique en somme de fonctions paires (cosinus) et impaires (sinus) est à la base de l'\emph{analyse de Fourier}, un outil indispensable pour l'étude des courants alternatifs, du son et des télécommunications. Reconnaître la parité d'un signal permet de diviser par deux le travail de calcul des techniciens et des ingénieurs.
\end{savaistu}

\begin{exercice}[À vous de jouer]
\begin{enumerate}
\item Étudier la parité des fonctions suivantes, définies sur $\mathbb{R}$ : $f(x)=3x^4-x^2+1$ ; $g(x)=x^5-4x^3$ ; $h(x)=x^2+x$.
\item Soit $u$ la fonction définie sur $\mathbb{R}^*$ par $u(x)=\dfrac{x^2+1}{x}$. Montrer que $u$ est impaire et préciser la conséquence graphique.
\end{enumerate}
\end{exercice}

\begin{corrige}[Exercice 1]
\textbf{1.} Les trois fonctions sont polynomiales, définies sur $\mathbb{R}$, ensemble symétrique par rapport à $0$.

$f(-x)=3(-x)^4-(-x)^2+1=3x^4-x^2+1=f(x)$ : $f$ est \textbf{paire} (tous les exposants sont pairs).

$g(-x)=(-x)^5-4(-x)^3=-x^5+4x^3=-(x^5-4x^3)=-g(x)$ : $g$ est \textbf{impaire} (tous les exposants sont impairs).

$h(-x)=x^2-x$ : ce n'est ni $h(x)$ ni $-h(x)$. Contre-exemple : $h(1)=2$ et $h(-1)=0$. Donc $h$ est \textbf{ni paire ni impaire}.

\textbf{2.} $D_u=\mathbb{R}^*$ est symétrique par rapport à $0$. Pour tout $x\neq 0$ :
\[ u(-x)=\frac{(-x)^2+1}{-x}=\frac{x^2+1}{-x}=-\frac{x^2+1}{x}=-u(x). \]
Donc $u$ est impaire : sa courbe est symétrique par rapport à l'origine $O$.
\end{corrige}

\begin{aretenir}[L'essentiel sur la parité]
\begin{itemize}
\item $f$ est \cle{paire} si $D_f$ est symétrique par rapport à $0$ et si $f(-x)=f(x)$ pour tout $x$ : la courbe est alors symétrique par rapport à \textbf{l'axe des ordonnées}.
\item $f$ est \cle{impaire} si $D_f$ est symétrique par rapport à $0$ et si $f(-x)=-f(x)$ pour tout $x$ : la courbe est alors symétrique par rapport à \textbf{l'origine}.
\item Toujours vérifier d'abord la symétrie de l'ensemble de définition.
\item $x^n$ a la parité de $n$ ; $\cos$ est paire, $\sin$ est impaire.
\item Une fonction peut être ni paire ni impaire : c'est le cas général.
\end{itemize}
\end{aretenir}

La parité fournit des symétries \emph{ponctuelles} de la courbe. Dans la section suivante, nous étudierons la \cle{périodicité}, une propriété de répétition qui, elle aussi, simplifie considérablement l'étude et le tracé des fonctions, en particulier des fonctions trigonométriques.

\coursec{Périodicité}

Dans la section précédente, nous avons étudié les \emph{symétries} d'une fonction : la parité (symétrie par rapport à l'axe des ordonnées) et l'imparité (symétrie par rapport à l'origine). Ces symétries permettent de réduire le domaine d'étude : il suffit d'étudier la fonction sur une moitié de son ensemble de définition, puis de compléter la courbe par symétrie.

Nous allons maintenant découvrir une autre propriété qui réduit encore plus le travail : la \cle{périodicité}. Alors que la parité « replie » la courbe sur elle-même, la périodicité la fait se \emph{répéter} indéfiniment, comme un motif de pagne qui se reproduit à l'identique le long du tissu. Un technicien rencontre cette situation partout : la tension alternative d'une prise électrique, le mouvement d'un piston dans un moteur, la rotation d'une roue de vélo sont des phénomènes qui se répètent à intervalles réguliers.

\courssub{Notion de fonction périodique}

\textbf{Intuition.} Observez l'aiguille d'une horloge : toutes les 12 heures, elle revient exactement dans la même position. Si l'on note $f(t)$ l'angle de l'aiguille à l'instant $t$, alors $f(t+12)=f(t)$ : la fonction $f$ « recommence » toutes les 12 heures. On dit qu'elle est périodique de période 12 heures. De même, la tension du réseau ENEO est alternative : elle reprend la même valeur tous les $\SI{0,02}{s}$ (fréquence de \SI{50}{Hz}).

\begin{definition}[Fonction périodique, période fondamentale]
Soit $f$ une fonction numérique définie sur un ensemble $D_f$ et $T$ un réel \textbf{strictement positif}.

On dit que $f$ est \cle{$T$-périodique} (ou que $T$ est \emph{une} période de $f$) si :
\begin{itemize}
\item pour tout $x\in D_f$, on a $x+T\in D_f$ ;
\item pour tout $x\in D_f$, on a $f(x+T)=f(x)$.
\end{itemize}

Lorsqu'il existe une \textbf{plus petite} valeur strictement positive de $T$ vérifiant ces conditions, on l'appelle la \cle{période fondamentale} de $f$ (ou simplement \emph{la période} de $f$).
\end{definition}

\begin{attention}[Conditions à ne pas oublier]
Deux pièges fréquents au Probatoire :
\begin{itemize}
\item On exige $T>0$. Dire « $f$ est périodique de période $-2\pi$ » est une erreur de rédaction : on prend toujours la valeur positive.
\item Il faut vérifier que $x+T$ reste dans l'ensemble de définition. Par exemple, pour la fonction tangente, si $x\in D_{\tan}$, alors $x+\pi\in D_{\tan}$ : la condition est bien satisfaite.
\end{itemize}
\end{attention}

\begin{propriete}[Multiples d'une période]
Si $f$ est $T$-périodique, alors pour tout entier relatif $k\in\mathbb{Z}$ et tout $x\in D_f$ tel que $x+kT\in D_f$ :
\[ f(x+kT)=f(x). \]
Autrement dit, tout multiple entier d'une période est encore une période.
\end{propriete}

\begin{experience}[Démonstration guidée de la propriété]
Démontrons que $f(x+kT)=f(x)$ pour tout $k\in\mathbb{Z}$.
\begin{enumerate}
\item \textbf{Pour $k$ positif.} On raisonne par récurrence sur $k\in\mathbb{N}$. Pour $k=0$, $f(x+0)=f(x)$ : c'est vrai. Supposons $f(x+kT)=f(x)$ pour un entier $k$. Alors, en appliquant la définition de la $T$-périodicité au point $x+kT$ :
\[ f\bigl(x+(k+1)T\bigr)=f\bigl((x+kT)+T\bigr)=f(x+kT)=f(x). \]
La propriété est donc vraie pour tout $k\in\mathbb{N}$.
\item \textbf{Pour $k$ négatif.} Posons $y=x+kT$. Comme $-k\in\mathbb{N}$, on a $f(y-kT)=f(y)$, c'est-à-dire $f(x)=f(x+kT)$.
\item \textbf{Conclusion.} Pour tout $k\in\mathbb{Z}$, $f(x+kT)=f(x)$. $\square$
\end{enumerate}
\end{experience}

\begin{exemplebox}[Périodes des fonctions trigonométriques]
Les fonctions trigonométriques sont les exemples fondamentaux de fonctions périodiques :
\begin{itemize}
\item la fonction \textbf{sinus} est $2\pi$-périodique : $\sin(x+2\pi)=\sin x$ pour tout $x\in\mathbb{R}$ ;
\item la fonction \textbf{cosinus} est $2\pi$-périodique : $\cos(x+2\pi)=\cos x$ pour tout $x\in\mathbb{R}$ ;
\item la fonction \textbf{tangente} est $\pi$-périodique : $\tan(x+\pi)=\tan x$ pour tout $x$ de son ensemble de définition.
\end{itemize}
Sur le cercle trigonométrique, ajouter $2\pi$ à un angle revient à faire un tour complet : on retombe sur le même point, donc sur les mêmes valeurs de sinus et de cosinus. De même, ajouter $\pi$ à un angle revient à passer au point diamétralement opposé : le rapport $\dfrac{\sin x}{\cos x}$ est inchangé, d'où la $\pi$-périodicité de la tangente.
\end{exemplebox}

\courssub{Interprétation graphique : répétition du motif}

Graphiquement, la $T$-périodicité se traduit ainsi : \textbf{la courbe de $f$ est invariante par les translations de vecteur $kT\,\vec{\imath}$} ($k\in\mathbb{Z}$). Concrètement, il suffit de tracer la courbe sur \emph{un seul} intervalle de longueur $T$ (appelé \cle{motif} ou \emph{patron}), puis de reproduire ce motif par translations successives vers la droite et vers la gauche.

\begin{popfigure}
\begin{tikzpicture}[scale=0.7]
\draw[->] (-1,0)--(20.5,0) node[right]{$x$};
\draw[->] (0,-1.6)--(0,1.6) node[above]{$y$};
\foreach \k in {0,1,2}{\draw[domain=\k*6.2832:\k*6.2832+6.2832,smooth,variable=\x,popblue,thick] plot ({\x},{sin(\x r)});}
\draw[popline,dashed] (6.2832,-1.6)--(6.2832,1.6);
\draw[popline,dashed] (12.5664,-1.6)--(12.5664,1.6);
\node[below] at (6.2832,0) {$2\pi$};
\node[below] at (12.5664,0) {$4\pi$};
\node[below] at (18.8496,0) {$6\pi$};
\node[popblue] at (19.5,1.2) {$y=\sin x$};
\end{tikzpicture}
\legende{La courbe de la fonction sinus ($T=2\pi$) : le motif tracé sur $[0\,;2\pi]$ se répète indéfiniment par translations.}
\end{popfigure}

\begin{aretenir}[Conséquence pratique pour l'étude d'une fonction]
Si $f$ est $T$-périodique, on l'étudie sur un intervalle de longueur $T$, par exemple $[0\,;T]$ ou $\left[-\dfrac{T}{2}\,;\dfrac{T}{2}\right]$. Si de plus $f$ est paire ou impaire, on peut réduire encore l'intervalle d'étude à $\left[0\,;\dfrac{T}{2}\right]$, puis compléter par symétrie et par translations.
\end{aretenir}

\courssub{Utiliser la périodicité pour calculer des images}

La périodicité permet de ramener le calcul d'une image en un point quelconque à un calcul sur l'intervalle principal $[0\,;T[$ : on « retire » autant de périodes que nécessaire.

\begin{methode}[Ramener un calcul à l'intervalle principal]
Pour calculer $f(x)$ lorsque $f$ est $T$-périodique et que $x$ est « loin » de l'intervalle principal :
\begin{enumerate}
\item Chercher l'entier $k\in\mathbb{Z}$ tel que $x-kT\in[0\,;T[$ (division euclidienne de $x$ par $T$).
\item Écrire $f(x)=f(x-kT)$ grâce à la propriété $f(x+kT)=f(x)$.
\item Calculer $f$ au point $x-kT$, qui appartient à l'intervalle principal.
\end{enumerate}
\end{methode}

\begin{exoresolu}[Calcul d'images par périodicité]
\textbf{Énoncé.} Soit $f$ la fonction $2$-périodique définie sur $[0\,;2[$ par $f(x)=x^2-2x$. Calculer $f(7)$, $f(10{,}5)$ et $f(-3)$.

\tcblower
\textbf{Solution détaillée.}

La fonction $f$ est $2$-périodique, donc $f(x+2k)=f(x)$ pour tout $k\in\mathbb{Z}$.

\textbf{Calcul de $f(7)$.} On écrit $7=3\times 2+1$, donc $k=3$ :
\[ f(7)=f(7-3\times 2)=f(1)=1^2-2\times 1=-1. \]

\textbf{Calcul de $f(10{,}5)$.} On écrit $10{,}5=5\times 2+0{,}5$, donc $k=5$ :
\[ f(10{,}5)=f(0{,}5)=(0{,}5)^2-2\times 0{,}5=0{,}25-1=-0{,}75. \]

\textbf{Calcul de $f(-3)$.} Ici $x$ est négatif : on \emph{ajoute} des périodes pour revenir dans $[0\,;2[$. On a $-3+2\times 2=1$, donc $k=-2$ :
\[ f(-3)=f(-3+4)=f(1)=-1. \]

Ainsi $f(7)=-1$, $f(10{,}5)=-0{,}75$ et $f(-3)=-1$.
\end{exoresolu}

\courssub{Période d'une fonction transformée}

\textbf{Intuition.} Si l'on « comprime » horizontalement la courbe d'une fonction périodique, le motif se répète plus vite : la période diminue. C'est ce qui se passe lorsqu'on remplace $x$ par $ax$.

\begin{propriete}[Période de $x\mapsto f(ax)$]
Soit $f$ une fonction $T$-périodique et $a$ un réel strictement positif. Alors la fonction $g$ définie par $g(x)=f(ax)$ est périodique de période $\dfrac{T}{a}$.
\end{propriete}

\begin{experience}[Démonstration guidée]
Vérifions que $\dfrac{T}{a}$ est une période de $g$. Pour tout $x$ de l'ensemble de définition :
\[ g\left(x+\frac{T}{a}\right)=f\left(a\left(x+\frac{T}{a}\right)\right)=f(ax+T)=f(ax)=g(x), \]
car $f$ est $T$-périodique. Donc $g$ est bien $\dfrac{T}{a}$-périodique. $\square$
\end{experience}

\begin{exemplebox}[Exemple chiffré fondamental]
La fonction $f(x)=\sin(2x)$ a pour période $\dfrac{2\pi}{2}=\pi$. Vérifions-le directement :
\[ f(x+\pi)=\sin\bigl(2(x+\pi)\bigr)=\sin(2x+2\pi)=\sin(2x)=f(x). \]
De même, $g(x)=\cos(3x)$ a pour période $\dfrac{2\pi}{3}$, et $h(x)=\tan(4x)$ a pour période $\dfrac{\pi}{4}$.

\emph{Application métier :} en électricité, une tension $u(t)=U_{\max}\sin(\omega t)$ a pour période $T=\dfrac{2\pi}{\omega}$. Pour le réseau camerounais, $T=\SI{0,02}{s}$, d'où $\omega=\dfrac{2\pi}{0{,}02}=100\pi\approx 314$ rad/s.
\end{exemplebox}

\courssub{Période d'une somme de fonctions périodiques}

Lorsqu'on additionne deux fonctions périodiques (par exemple deux signaux électriques superposés), une question naturelle se pose : la somme est-elle périodique, et si oui, quelle est sa période ?

\begin{methode}[Déterminer la période d'une somme]
Soient $f$ de période $T_1$ et $g$ de période $T_2$. Pour trouver la période de $f+g$ :
\begin{enumerate}
\item Vérifier que le rapport $\dfrac{T_1}{T_2}$ est un nombre \textbf{rationnel} (quotient de deux entiers).
\item Si oui, chercher le plus petit réel $T>0$ qui est à la fois un multiple entier de $T_1$ et un multiple entier de $T_2$ : c'est le \cle{PPCM} des périodes (plus petit commun multiple).
\item Conclure : $f+g$ est $T$-périodique (attention, la période fondamentale peut parfois être plus petite ; il faut vérifier).
\end{enumerate}
\end{methode}

\begin{exoresolu}[Période d'une somme]
\textbf{Énoncé.} Déterminer la période de la fonction $h$ définie par $h(x)=\sin(2x)+\cos(3x)$.

\tcblower
\textbf{Solution détaillée.}

\begin{enumerate}
\item $x\mapsto\sin(2x)$ a pour période $T_1=\dfrac{2\pi}{2}=\pi$.
\item $x\mapsto\cos(3x)$ a pour période $T_2=\dfrac{2\pi}{3}$.
\item Le rapport $\dfrac{T_1}{T_2}=\dfrac{\pi}{2\pi/3}=\dfrac{3}{2}$ est rationnel : la somme est périodique.
\item Cherchons le PPCM de $\pi$ et $\dfrac{2\pi}{3}$. On a $2\pi=2\times\pi=3\times\dfrac{2\pi}{3}$. Donc $2\pi$ est un multiple commun aux deux périodes, et c'est le plus petit.
\item Vérification : $h(x+2\pi)=\sin(2x+4\pi)+\cos(3x+6\pi)=\sin(2x)+\cos(3x)=h(x)$.
\end{enumerate}
La fonction $h$ est donc $2\pi$-périodique.
\end{exoresolu}

\begin{attention}[Somme de fonctions à périodes incommensurables]
La somme de deux fonctions périodiques \textbf{n'est pas toujours périodique} ! Si le rapport $\dfrac{T_1}{T_2}$ est \emph{irrationnel} (on dit que les périodes sont \emph{incommensurables}), il n'existe aucun réel $T>0$ qui soit multiple des deux périodes à la fois, et la somme n'a aucune période.

\textbf{Exemple :} $f(x)=\sin(x)+\sin(\pi x)$. Les périodes sont $2\pi$ et $2$ ; le rapport $\dfrac{2\pi}{2}=\pi$ est irrationnel. La fonction $f$ n'est \textbf{pas} périodique : son graphique ne se répète jamais exactement.
\end{attention}

\courssub{Lien entre périodicité et parité}

La périodicité et la parité sont deux propriétés qui se combinent très bien, et les fonctions trigonométriques en sont le meilleur exemple :
\begin{itemize}
\item la fonction \textbf{cosinus} est \textbf{paire} ($\cos(-x)=\cos x$) et $2\pi$-périodique : sa courbe est symétrique par rapport à l'axe des ordonnées et se répète tous les $2\pi$ ;
\item la fonction \textbf{sinus} est \textbf{impaire} ($\sin(-x)=-\sin x$) et $2\pi$-périodique : sa courbe est symétrique par rapport à l'origine et se répète tous les $2\pi$ ;
\item la fonction \textbf{tangente} est \textbf{impaire} et $\pi$-périodique.
\end{itemize}

\begin{aretenir}[Stratégie complète de réduction du domaine d'étude]
Pour étudier une fonction à la fois paire (ou impaire) et $T$-périodique :
\begin{enumerate}
\item la périodicité ramène l'étude à un intervalle de longueur $T$, par exemple $\left[-\dfrac{T}{2}\,;\dfrac{T}{2}\right]$ ;
\item la parité ramène l'étude à la moitié positive de cet intervalle, soit $\left[0\,;\dfrac{T}{2}\right]$ ;
\item on complète ensuite la courbe par symétrie, puis par translations de vecteur $kT\,\vec{\imath}$.
\end{enumerate}
Par exemple, pour la fonction cosinus, il suffit de l'étudier sur $[0\,;\pi]$.
\end{aretenir}

\begin{exercice}[Périodicité et calculs d'images]
Soit $f$ la fonction définie sur $\mathbb{R}$ par $f(x)=\cos(4x)$.
\begin{enumerate}
\item Montrer que $f$ est $\dfrac{\pi}{2}$-périodique.
\item Étudier la parité de $f$.
\item Calculer $f\left(\dfrac{9\pi}{4}\right)$ en se ramenant à l'intervalle $\left[0\,;\dfrac{\pi}{2}\right[$.
\end{enumerate}
\end{exercice}

\begin{corrige}[Exercice 1]
\begin{enumerate}
\item Pour tout réel $x$ :
\[ f\left(x+\frac{\pi}{2}\right)=\cos\left(4\left(x+\frac{\pi}{2}\right)\right)=\cos(4x+2\pi)=\cos(4x)=f(x). \]
Donc $f$ est $\dfrac{\pi}{2}$-périodique.
\item $D_f=\mathbb{R}$ est symétrique par rapport à $0$, et $f(-x)=\cos(-4x)=\cos(4x)=f(x)$ : $f$ est \textbf{paire}.
\item On retire des multiples de $\dfrac{\pi}{2}$ : $\dfrac{9\pi}{4}-4\times\dfrac{\pi}{2}=\dfrac{9\pi}{4}-2\pi=\dfrac{\pi}{4}$, et $\dfrac{\pi}{4}\in\left[0\,;\dfrac{\pi}{2}\right[$. Donc :
\[ f\left(\frac{9\pi}{4}\right)=f\left(\frac{\pi}{4}\right)=\cos\left(4\times\frac{\pi}{4}\right)=\cos(\pi)=-1. \]
\end{enumerate}
\end{corrige}

\begin{exercice}[Période d'une somme de signaux]
Un technicien superpose deux tensions alternatives : $u_1(t)=3\sin(100\pi t)$ et $u_2(t)=2\cos(50\pi t)$.
\begin{enumerate}
\item Déterminer les périodes $T_1$ et $T_2$ de $u_1$ et $u_2$.
\item Montrer que la tension totale $u=u_1+u_2$ est périodique et déterminer sa période.
\end{enumerate}
\end{exercice}

\begin{corrige}[Exercice 2]
\begin{enumerate}
\item $T_1=\dfrac{2\pi}{100\pi}=\dfrac{1}{50}=\SI{0,02}{s}$ et $T_2=\dfrac{2\pi}{50\pi}=\dfrac{1}{25}=\SI{0,04}{s}$.
\item Le rapport $\dfrac{T_2}{T_1}=2$ est rationnel : la somme est périodique. Le plus petit multiple commun est $T_2=2T_1=\SI{0,04}{s}$. Vérification :
\[ u(t+0{,}04)=3\sin(100\pi t+4\pi)+2\cos(50\pi t+2\pi)=3\sin(100\pi t)+2\cos(50\pi t)=u(t). \]
La tension totale est donc périodique de période $\SI{0,04}{s}$.
\end{enumerate}
\end{corrige}

\begin{savaistu}[La fonction « partie fractionnaire »]
Pour tout réel $x$, on note $\{x\}$ sa \cle{partie fractionnaire}, c'est-à-dire $\{x\}=x-E(x)$, où $E(x)$ est la partie entière de $x$ (le plus grand entier inférieur ou égal à $x$). Par exemple $\{3{,}7\}=0{,}7$ et $\{-1{,}2\}=-1{,}2-(-2)=0{,}8$.

Cette fonction est \textbf{$1$-périodique} : $\{x+1\}=\{x\}$ pour tout $x$. Sa courbe est une succession de segments parallèles en « dents de scie ». On l'utilise en informatique et en traitement du signal pour modéliser les phénomènes qui « repartent à zéro » à intervalles réguliers, comme le compteur kilométrique journalier d'un véhicule de transport qui repasse à zéro chaque matin, ou les rampes de balayage des anciens écrans d'oscilloscope utilisés dans les ateliers d'électronique.
\end{savaistu}

\coursec{Courbe d'une fonction : tracé et analyse}

Dans la section précédente, nous avons vu que la \cle{périodicité} permet de réduire l'étude d'une fonction à un intervalle de longueur $T$, le reste de la courbe s'obtenant par translations. Nous disposons désormais de tous les outils (ensemble de définition, parité, périodicité, limites, dérivée) pour répondre à la question centrale de ce chapitre : \emph{comment tracer et analyser la courbe représentative d'une fonction numérique ?} C'est une compétence essentielle au Probatoire technique, mais aussi dans les métiers : un électricien lit la courbe de tension, un mécanicien la courbe de couple d'un moteur, un gestionnaire la courbe de coût de production.

\courssub{Démarche générale d'étude et de tracé}

Tracer la courbe d'une fonction ne se fait jamais « au hasard », point par point. On suit une démarche ordonnée, que l'on peut résumer en sept étapes.

\begin{methode}[Étapes du tracé de la courbe d'une fonction]
Pour étudier une fonction $f$ et tracer sa courbe représentative $\mathcal{C}_f$ :
\begin{enumerate}
\item \textbf{Ensemble de définition} : déterminer $D_f$ (valeurs interdites : dénominateur nul, quantité négative sous une racine carrée).
\item \textbf{Réduction du domaine d'étude} : chercher la \cle{parité} (fonction paire ou impaire : on étudie sur la moitié du domaine et on complète par symétrie) et la \cle{périodicité} (on étudie sur une période et on complète par translation).
\item \textbf{Limites aux bornes} : calculer les limites de $f$ aux bornes ouvertes de $D_f$ (bornes infinies et valeurs interdites).
\item \textbf{Asymptotes} : déduire des limites les asymptotes éventuelles (verticales, horizontales, obliques).
\item \textbf{Variations} : calculer $f'(x)$, étudier son signe, dresser le \cle{tableau de variations} complet (avec les limites).
\item \textbf{Points remarquables} : déterminer les intersections avec les axes ($f(x)=0$ et $f(0)$), les extremums, et quelques points bien choisis.
\item \textbf{Tracé} : placer les asymptotes en pointillés, les points remarquables, puis esquisser la courbe en respectant le tableau de variations.
\end{enumerate}
\end{methode}

\begin{attention}[Respecter l'ensemble de définition]
On ne trace \textbf{jamais} la courbe en dehors de l'ensemble de définition. Si $a \notin D_f$, la courbe ne possède \emph{aucun} point d'abscisse $a$ : il est faux de « franchir » la droite $x = a$ avec le crayon. De même, une valeur interdite donne souvent lieu à une \cle{asymptote verticale} que la courbe ne coupe pas. Vérifier $D_f$ avant tout tracé est le premier réflexe exigé au Probatoire.
\end{attention}

\courssub{Asymptotes : intuition et définitions}

\textbf{Intuition.} Une \cle{asymptote} est une droite dont la courbe se rapproche indéfiniment, sans jamais la « quitter » définitivement. Imaginez un fil électrique tendu le long d'un mur : la courbe est le mur, l'asymptote est le fil ; plus on avance, plus l'écart se réduit.

\begin{definition}[Asymptote verticale]
Soit $f$ une fonction et $a$ un réel (borne finie de $D_f$). La droite d'équation $x = a$ est \cle{asymptote verticale} à la courbe $\mathcal{C}_f$ lorsque :
\[ \lim_{x \to a} f(x) = +\infty \quad \text{ou} \quad \lim_{x \to a} f(x) = -\infty \]
(limite à gauche, à droite, ou les deux).
\end{definition}

\begin{propriete}[Recherche d'une asymptote verticale]
Pour détecter une asymptote verticale, on calcule la limite de $f$ en chaque \textbf{valeur interdite} $a$ de l'ensemble de définition. Si cette limite est infinie, la droite $x = a$ est asymptote verticale. Le signe de l'infini (à gauche et à droite) indique de quel côté la courbe « monte » ou « descend » le long de l'asymptote.
\end{propriete}

\begin{definition}[Asymptote horizontale]
La droite d'équation $y = \ell$ est \cle{asymptote horizontale} à $\mathcal{C}_f$ en $+\infty$ (resp. $-\infty$) lorsque :
\[ \lim_{x \to +\infty} f(x) = \ell \quad \left(\text{resp. } \lim_{x \to -\infty} f(x) = \ell\right), \qquad \ell \in \mathbb{R}. \]
\end{definition}

\begin{propriete}[Asymptote oblique]
La droite d'équation $y = ax + b$ (avec $a \neq 0$) est \cle{asymptote oblique} à $\mathcal{C}_f$ en $+\infty$ (ou en $-\infty$) lorsque :
\[ \lim_{x \to \pm\infty} \big(f(x) - (ax + b)\big) = 0. \]
Pour une fonction rationnelle dont le degré du numérateur égale le degré du dénominateur plus un, on obtient $ax + b$ par \textbf{division euclidienne} du numérateur par le dénominateur : le quotient est $ax + b$, et le reste divisé par le dénominateur tend vers $0$.
\end{propriete}

\begin{exemplebox}[Détection d'une asymptote oblique par division]
Soit $f(x) = \dfrac{x^2 - 1}{x - 2}$. Effectuons la division de $x^2 - 1$ par $x - 2$ :
\[ x^2 - 1 = (x - 2)(x + 2) + 3. \]
Donc $f(x) = x + 2 + \dfrac{3}{x - 2}$. Comme $\displaystyle\lim_{x \to \pm\infty} \frac{3}{x-2} = 0$, la droite $y = x + 2$ est asymptote oblique à $\mathcal{C}_f$ en $-\infty$ et en $+\infty$.
\end{exemplebox}

\begin{attention}[Pièges fréquents sur les asymptotes]
\begin{itemize}
\item Une asymptote horizontale en $+\infty$ n'empêche pas une asymptote oblique en $-\infty$ (et inversement) : il faut étudier \textbf{les deux} bornes.
\item Si $\lim f(x) = \pm\infty$ en l'infini, il n'y a \textbf{pas} d'asymptote horizontale ; il faut alors chercher une asymptote oblique.
\item Une courbe peut \textbf{couper} son asymptote horizontale ou oblique : l'asymptote décrit un comportement à l'infini, pas une interdiction de croisement.
\end{itemize}
\end{attention}

\courssub{Du tableau de variations à l'esquisse de la courbe}

Le tableau de variations est le « plan de construction » de la courbe. Il résume : le signe de $f'$ (donc le sens de variation), les limites aux bornes et les extremums.

\begin{methode}[Esquisser la courbe à partir du tableau de variations]
\begin{enumerate}
\item Tracer en pointillés toutes les asymptotes trouvées.
\item Placer les points correspondant aux extremums (lire leurs coordonnées dans le tableau).
\item Placer les intersections avec les axes.
\item Entre deux colonnes du tableau, la courbe est \textbf{monotone} : la tracer en montant ($\nearrow$) ou en descendant ($\searrow$), en rejoignant les limites indiquées.
\item Vérifier la cohérence : la courbe longe les asymptotes sans les franchir aux bornes concernées.
\end{enumerate}
\end{methode}

\courssub{Exemple complet : étude de $f(x) = \dfrac{x^2 - 1}{x - 2}$}

\begin{exoresolu}[Étude complète d'une fonction rationnelle]
Soit $f$ la fonction définie par $f(x) = \dfrac{x^2 - 1}{x - 2}$ et $\mathcal{C}_f$ sa courbe représentative.
\begin{enumerate}
\item Déterminer l'ensemble de définition de $f$.
\item Calculer les limites aux bornes et en déduire les asymptotes.
\item Étudier les variations de $f$ et dresser le tableau de variations.
\item Tracer $\mathcal{C}_f$.
\end{enumerate}
\tcblower
\textbf{1. Ensemble de définition.} $f$ est définie si et seulement si $x - 2 \neq 0$, donc :
\[ D_f = \mathbb{R} \setminus \{2\} = \left]-\infty ; 2\right[ \cup \left]2 ; +\infty\right[. \]
La fonction n'est ni paire ni impaire (le domaine n'est pas symétrique par rapport à $0$), ni périodique.

\textbf{2. Limites et asymptotes.}
\begin{itemize}
\item En $2$ : le numérateur vaut $2^2 - 1 = 3 > 0$. À gauche de $2$, $x - 2 < 0$, donc $\displaystyle\lim_{x \to 2^-} f(x) = -\infty$ ; à droite, $x - 2 > 0$, donc $\displaystyle\lim_{x \to 2^+} f(x) = +\infty$. La droite $x = 2$ est \textbf{asymptote verticale}.
\item En l'infini : $\displaystyle\lim_{x \to \pm\infty} f(x) = \lim_{x \to \pm\infty} \frac{x^2}{x} = \lim_{x \to \pm\infty} x = \pm\infty$. Pas d'asymptote horizontale.
\item Asymptote oblique : la division $x^2 - 1 = (x-2)(x+2) + 3$ donne $f(x) = x + 2 + \dfrac{3}{x-2}$, et $\displaystyle\lim_{x \to \pm\infty}\frac{3}{x-2} = 0$. La droite $y = x + 2$ est \textbf{asymptote oblique} en $-\infty$ et en $+\infty$.
\end{itemize}

\textbf{3. Variations.} $f$ est dérivable sur $D_f$ (quotient de polynômes) et, avec la formule $\left(\dfrac{u}{v}\right)' = \dfrac{u'v - uv'}{v^2}$ :
\begin{align*}
f'(x) &= \frac{2x(x-2) - (x^2-1)\cdot 1}{(x-2)^2} \\
      &= \frac{2x^2 - 4x - x^2 + 1}{(x-2)^2} \\
      &= \frac{x^2 - 4x + 1}{(x-2)^2}.
\end{align*}
Le dénominateur est strictement positif sur $D_f$ ; le signe de $f'$ est celui du trinôme $x^2 - 4x + 1$. Son discriminant est $\Delta = 16 - 4 = 12$, d'où les racines :
\[ x_1 = 2 - \sqrt{3} \approx 0{,}27 \qquad \text{et} \qquad x_2 = 2 + \sqrt{3} \approx 3{,}73. \]
Le trinôme est positif à l'extérieur des racines, négatif entre elles.
Calculons les valeurs aux extremums en utilisant l'expression $f(x) = x + 2 + \dfrac{3}{x-2}$ :
\begin{itemize}
\item $x_1 = 2 - \sqrt{3}$, alors $x_1 - 2 = -\sqrt{3}$.
  $f(x_1) = (2-\sqrt{3}) + 2 + \dfrac{3}{-\sqrt{3}} = 4 - \sqrt{3} - \sqrt{3} = 4 - 2\sqrt{3} \approx 0{,}54$.
\item $x_2 = 2 + \sqrt{3}$, alors $x_2 - 2 = \sqrt{3}$.
  $f(x_2) = (2+\sqrt{3}) + 2 + \dfrac{3}{\sqrt{3}} = 4 + \sqrt{3} + \sqrt{3} = 4 + 2\sqrt{3} \approx 7{,}46$.
\end{itemize}

\begin{center}
\begin{tabular}{|c|ccccccccc|}\hline
$x$ & $-\infty$ & & $2-\sqrt{3}$ & & $2$ & & $2+\sqrt{3}$ & & $+\infty$ \\ \hline
$f'(x)$ & & $+$ & $0$ & $-$ & & $-$ & $0$ & $+$ & \\ \hline
$f$ & $-\infty$ & $\nearrow$ & $4-2\sqrt{3}$ & $\searrow$ & $-\infty \;\Vert\; +\infty$ & $\searrow$ & $4+2\sqrt{3}$ & $\nearrow$ & $+\infty$ \\ \hline
\end{tabular}
\end{center}

\textbf{4. Tracé.} On place les asymptotes $x = 2$ et $y = x + 2$ en pointillés, le maximum local $\left(2-\sqrt{3}\,;\,4-2\sqrt{3}\right)$, le minimum local $\left(2+\sqrt{3}\,;\,4+2\sqrt{3}\right)$, les intersections avec l'axe des abscisses ($f(x)=0 \Leftrightarrow x^2=1$, soit $x=-1$ et $x=1$) et le point $\left(0\,;\,\tfrac{1}{2}\right)$. On esquisse ensuite la courbe en suivant le tableau de variations (voir figure ci-dessous).
\end{exoresolu}

\begin{popfigure}
\begin{tikzpicture}[scale=0.75, >=Stealth]
% Grille légère
\draw[popline, step=1] (-4,-6) grid (5,8);
% Axes
\draw[->, thick] (-4,0) -- (5,0) node[right] {$x$};
\draw[->, thick] (0,-6) -- (0,8) node[above] {$y$};
% Graduations
\foreach \x in {-3,-2,-1,1,2,3,4} \draw (\x,0.1) -- (\x,-0.1) node[below] {\small $\x$};
\foreach \y in {-5,-4,-3,-2,-1,1,2,3,4,5,6,7} \draw (0.1,\y) -- (-0.1,\y) node[left] {\small $\y$};
% Asymptote verticale x=2
\draw[poporange, dashed, thick] (2,-6) -- (2,8) node[above right] {$x=2$};
% Asymptote oblique y=x+2
\draw[poporange, dashed, thick, domain=-3.5:3.5] plot (\x, {\x+2}) node[above right] {$y=x+2$};
% Courbe branche gauche
\draw[popblue, thick, domain=-3:1.9, smooth, variable=\x, samples=200] plot ({\x}, {(\x*\x-1)/(\x-2)});
% Courbe branche droite
\draw[popblue, thick, domain=2.1:3.8, smooth, variable=\x, samples=200] plot ({\x}, {(\x*\x-1)/(\x-2)});
% Points remarquables
\fill[popblue] (-1,0) circle (2pt) node[below left] {$(-1,0)$};
\fill[popblue] (1,0) circle (2pt) node[above right] {$(1,0)$};
\fill[popblue] (0,0.5) circle (2pt) node[left] {$(0,0,5)$};
% Extremums
\fill[poppurple] (0.268, 0.536) circle (2pt) node[above] {$M$};
\fill[poppurple] (3.732, 7.464) circle (2pt) node[above right] {$m$};
\end{tikzpicture}
\legende{Courbe de $f(x)=\dfrac{x^2-1}{x-2}$ avec ses asymptotes verticale $x=2$ et oblique $y=x+2$.}
\end{popfigure}

\courssub{Lecture graphique : exploiter une courbe donnée}

Savoir \emph{lire} une courbe est aussi important que savoir la tracer. Face à une représentation graphique, on doit pouvoir extraire :

\begin{itemize}
\item \textbf{L'ensemble image} : ensemble des ordonnées effectivement atteintes. On projette la courbe sur l'axe des ordonnées. Pour la fonction étudiée ci-dessus, la courbe suggère que les valeurs comprises entre $4 - 2\sqrt{3}$ et $4 + 2\sqrt{3}$ ne sont pas atteintes : l'équation $f(x) = k$ n'a pas de solution pour ces valeurs de $k$.
\item \textbf{La monotonie} : on repère les intervalles où la courbe « monte » (fonction croissante) ou « descend » (fonction décroissante).
\item \textbf{Les extremums} : un \cle{maximum local} est un point où la courbe atteint un « sommet » ; un \cle{minimum local} un point où elle atteint un « creux ». En ces points, la tangente est horizontale ($f'(x) = 0$).
\item \textbf{Les solutions d'équations} : les solutions de $f(x) = k$ sont les abscisses des points d'intersection de $\mathcal{C}_f$ avec la droite horizontale $y = k$ ; les solutions de $f(x) = 0$ sont les abscisses des intersections avec l'axe des abscisses.
\end{itemize}

\begin{aretenir}[Lecture graphique : les réflexes]
\begin{itemize}
\item Image de $f$ = projection de la courbe sur l'axe $(Oy)$.
\item $f(x) = k$ : intersections avec la droite $y = k$.
\item Maximum / minimum local : sommet / creux, tangente horizontale.
\item Signe de $f$ : position de la courbe par rapport à l'axe $(Ox)$.
\end{itemize}
\end{aretenir}

\begin{exemplebox}[Lecture de l'ensemble image]
Pour $f(x) = \dfrac{x^2-1}{x-2}$, le tableau de variations montre que sur $]-\infty ; 2[$, $f$ prend toutes les valeurs de $]-\infty ; 4 - 2\sqrt{3}]$, et sur $]2 ; +\infty[$, toutes les valeurs de $[4 + 2\sqrt{3} ; +\infty[$. L'ensemble image est donc :
\[ f(D_f) = \left]-\infty ; 4 - 2\sqrt{3}\right] \cup \left[4 + 2\sqrt{3} ; +\infty\right[. \]
L'équation $f(x) = 4$, par exemple, n'a aucune solution car $4 \in \left]4 - 2\sqrt{3} ; 4 + 2\sqrt{3}\right[$.
\end{exemplebox}

\begin{exercice}[À vous de jouer]
Soit $g$ la fonction définie par $g(x) = \dfrac{x^2 + 1}{x - 1}$.
\begin{enumerate}
\item Déterminer l'ensemble de définition de $g$.
\item Calculer les limites de $g$ aux bornes et préciser les asymptotes.
\item Montrer que $g(x) = x + 1 + \dfrac{2}{x - 1}$ et en déduire une asymptote oblique.
\item Étudier le signe de $g'(x) = \dfrac{x^2 - 2x - 1}{(x-1)^2}$ et dresser le tableau de variations.
\end{enumerate}
\end{exercice}

\begin{corrige}[Exercice 1]
\textbf{1.} $D_g = \mathbb{R} \setminus \{1\}$.
\textbf{2.} En $1$ : numérateur égal à $2 > 0$, donc $\lim_{1^-} g = -\infty$ et $\lim_{1^+} g = +\infty$ : asymptote verticale $x = 1$. En l'infini, $g(x) \sim x$, donc $\lim_{\pm\infty} g = \pm\infty$.
\textbf{3.} $x^2 + 1 = (x-1)(x+1) + 2$, donc $g(x) = x + 1 + \dfrac{2}{x-1}$ et $\lim_{x\to\pm\infty}\frac{2}{x-1} = 0$ : la droite $y = x + 1$ est asymptote oblique.
\textbf{4.} $g'(x)$ a le signe de $x^2 - 2x - 1$, de racines $1 - \sqrt{2}$ et $1 + \sqrt{2}$. Donc $g$ est croissante sur $]-\infty ; 1-\sqrt{2}]$, décroissante sur $[1-\sqrt{2} ; 1[$ et sur $]1 ; 1+\sqrt{2}]$, croissante sur $[1+\sqrt{2} ; +\infty[$. Maximum local en $1-\sqrt{2}$ (valeur $2-2\sqrt{2}$), minimum local en $1+\sqrt{2}$ (valeur $2+2\sqrt{2}$).
\end{corrige}

\begin{savaistu}[Une courbe « sans lever le crayon »]
Une fonction \cle{continue} sur un intervalle est une fonction dont la courbe se trace \emph{sans lever le crayon} : c'est un arc sans rupture, sans saut. Les fonctions polynômes, les fonctions rationnelles (sur chaque intervalle de leur domaine), les fonctions usuelles sont continues. C'est ce qui justifie, dans notre exemple, que chaque morceau de la courbe (à gauche et à droite de $x = 2$) forme un arc d'un seul tenant. En revanche, la valeur interdite $x = 2$ crée une rupture : la courbe est en \textbf{deux branches} séparées par l'asymptote verticale. Ce principe de continuité est aussi à la base du théorème des valeurs intermédiaires, très utilisé pour localiser les solutions d'équations, par exemple lorsqu'un technicien cherche la valeur de réglage qui annule une erreur de mesure.
\end{savaistu}

\coursec{Fonctions associées à une fonction donnée}

\noindent Après avoir étudié le tracé et l'analyse de la courbe d'une fonction (section 5), nous allons voir comment construire de nouvelles courbes à partir d'une courbe de référence $C_f$. Ces transformations — \cle{translation}, \cle{homothétie}, \cle{symétrie} et passage à la \cle{fonction inverse} — sont des outils constants en technologie (réglage de capteurs, mise à l'échelle de plans, symétrie de pièces mécaniques, conception assistée par ordinateur). Elles permettent de prévoir le comportement d'une fonction modifiée sans refaire tout le tableau de valeurs.

\courssub{Translations : décalages horizontaux et verticaux}

\begin{definition}[Translation horizontale]
Soit $f$ une fonction définie sur un intervalle $I$ et $a\in\mathbb{R}$. La fonction $g$ définie par $g(x)=f(x-a)$ (avec $x-a\in I$) est l'image de $f$ par une \cle{translation horizontale} de vecteur $\vec{u}(a,0)$.
\begin{itemize}
\item Si $a>0$, la courbe $C_g$ se décale de $a$ unités vers la \textbf{droite}.
\item Si $a<0$, la courbe $C_g$ se décale de $|a|$ unités vers la \textbf{gauche}.
\end{itemize}
Le domaine de définition est lui aussi translaté de $a$.
\end{definition}

\begin{definition}[Translation verticale]
Soit $f$ une fonction définie sur $I$ et $b\in\mathbb{R}$. La fonction $h$ définie par $h(x)=f(x)+b$ est l'image de $f$ par une \cle{translation verticale} de vecteur $\vec{v}(0,b)$.
\begin{itemize}
\item Si $b>0$, $C_h$ monte de $b$ unités.
\item Si $b<0$, $C_h$ descend de $|b|$ unités.
\end{itemize}
Le domaine de définition ne change pas.
\end{definition}

\begin{propriete}[Translation composée]
La courbe de la fonction $x\mapsto f(x-a)+b$ s'obtient par translation de la courbe $C_f$ du vecteur $\vec{w}(a,b)$. L'ordre n'a pas d'importance : translation horizontale puis verticale donne le même résultat que l'inverse.
\end{propriete}

\begin{exemplebox}[Exemple : translations d'une parabole]
Soit $f(x)=x^2$ définie sur $\mathbb{R}$. On considère $g(x)=(x-1)^2$ et $h(x)=x^2+2$.
\begin{itemize}
\item $g(x)=f(x-1)$ : translation de $+1$ sur l'axe des abscisses (vers la droite). Sommet en $(1,0)$.
\item $h(x)=f(x)+2$ : translation de $+2$ sur l'axe des ordonnées (vers le haut). Sommet en $(0,2)$.
\item $k(x)=(x-1)^2+2 = f(x-1)+2$ : translation du vecteur $(1,2)$. Sommet en $(1,2)$.
\end{itemize}
\end{exemplebox}

\begin{popfigure}
\begin{tikzpicture}[scale=0.7]
\draw[->] (-3,0)--(3,0) node[right]{$x$};
\draw[->] (0,-1)--(0,5) node[above]{$y$};
\draw[domain=-2:2,smooth,variable=\x,popblue,thick] plot ({\x},{\x*\x});
\draw[domain=-1:3,smooth,variable=\x,poporange,dashed,thick] plot ({\x},{(\x-1)*(\x-1)});
\draw[domain=-1.7:1.7,smooth,variable=\x,popgreen,dotted,thick] plot ({\x},{\x*\x+2});
\node[popblue] at (2.2,3.5) {$f(x)=x^2$};
\node[poporange] at (2.8,1.5) {$g(x)=(x-1)^2$};
\node[popgreen] at (2.2,5) {$h(x)=x^2+2$};
\end{tikzpicture}
\legende{Translations horizontale (orange) et verticale (vert) de la parabole $f(x)=x^2$.}
\end{popfigure}

\courssub{Homothéties : étirements et compressions}

\begin{definition}[Homothétie par rapport à l'axe des abscisses (verticale)]
Soit $f$ définie sur $I$ et $k\in\mathbb{R}^*$. La fonction $g(x)=k\,f(x)$ correspond à une \cle{homothétie verticale} de centre l'axe des abscisses ($Ox$) et de rapport $k$.
\begin{itemize}
\item Si $|k|>1$, la courbe s'\textbf{étire} verticalement (facteur $|k|$).
\item Si $0<|k|<1$, elle se \textbf{comprime} verticalement (facteur $|k|$).
\item Si $k<0$, il y a en plus une symétrie par rapport à l'axe des abscisses.
\end{itemize}
Le domaine de définition ne change pas.
\end{definition}

\begin{definition}[Homothétie par rapport à l'axe des ordonnées (horizontale)]
Soit $f$ définie sur $I$ et $k\in\mathbb{R}^*$. La fonction $h(x)=f(kx)$ correspond à une \cle{homothétie horizontale} de centre l'axe des ordonnées ($Oy$) et de rapport $1/k$ (\textbf{attention} : le facteur horizontal est l'inverse du coefficient de $x$).
\begin{itemize}
\item Si $|k|>1$, compression horizontale (facteur $1/|k|$).
\item Si $0<|k|<1$, étirement horizontal (facteur $1/|k|$).
\item Si $k<0$, symétrie par rapport à l'axe des ordonnées en plus.
\end{itemize}
Le domaine de définition est divisé par $k$ (bornes divisées par $k$).
\end{definition}

\begin{methode}[Tracer $C_{kf}$ et $C_{f(kx)}$ à partir de $C_f$]
\begin{enumerate}
\item Repérer quelques points caractéristiques de $C_f$ : intercepts, extremums, points d'inflexion, asymptotes.
\item Pour $g(x)=kf(x)$ (verticale) : multiplier l'ordonnée de chaque point par $k$ (l'abscisse ne change pas).
\item Pour $h(x)=f(kx)$ (horizontale) : diviser l'abscisse de chaque point par $k$ (l'ordonnée ne change pas).
\item Relier les nouveaux points en respectant la forme générale (convexité, asymptotes, sens de variation).
\end{enumerate}
\end{methode}

\begin{exemplebox}[Exemple : homothéties de la racine carrée]
Soit $f(x)=\sqrt{x}$ ($x\ge0$).
\begin{itemize}
\item $g(x)=2\sqrt{x}$ : homothétie verticale de rapport $2$ (la courbe s'éloigne de l'axe $Ox$).
\item $h(x)=\sqrt{2x}$ : homothétie horizontale de rapport $1/2$ (compression vers l'axe $Oy$).
\end{itemize}
\end{exemplebox}

\begin{popfigure}
\begin{tikzpicture}[scale=0.8]
\draw[->] (0,0)--(5,0) node[right]{$x$};
\draw[->] (0,0)--(0,3.5) node[above]{$y$};
\draw[domain=0:4,smooth,variable=\x,popblue,thick] plot ({\x},{sqrt(max(0,\x))});
\draw[domain=0:2.25,smooth,variable=\x,poporange,dashed,thick] plot ({\x},{2*sqrt(max(0,\x))});
\draw[domain=0:2,smooth,variable=\x,popgreen,dotted,thick] plot ({\x},{sqrt(max(0,2*\x))});
\node[popblue] at (4.2,1.8) {$f(x)=\sqrt{x}$};
\node[poporange] at (2.2,3.2) {$g(x)=2\sqrt{x}$};
\node[popgreen] at (2.2,1.5) {$h(x)=\sqrt{2x}$};
\end{tikzpicture}
\legende{Homothéties verticale (orange) et horizontale (vert) de la fonction racine carrée.}
\end{popfigure}

\courssub{Symétries axiale et centrale}

\begin{definition}[Symétrie axiale par rapport à l'axe des ordonnées]
La fonction $g(x)=f(-x)$ a pour courbe $C_g$ l'image de $C_f$ par la \cle{symétrie axiale} d'axe $Oy$ (l'axe des ordonnées). Cela correspond à remplacer $x$ par $-x$ dans l'expression de $f$. Le domaine de définition est symétrisé par rapport à $0$.
\end{definition}

\begin{definition}[Symétrie centrale par rapport à l'origine]
La fonction $h(x)=-f(-x)$ a pour courbe $C_h$ l'image de $C_f$ par la \cle{symétrie centrale} de centre $O(0,0)$ (rotation de $180^\circ$). C'est la composition des deux symétries axiales (par rapport à $Ox$ puis $Oy$, ou l'inverse).
\end{definition}

\begin{propriete}[Lien avec la parité]
\begin{itemize}
\item $f$ est \textbf{paire} $\iff$ $f(-x)=f(x)$ pour tout $x$ du domaine $\iff$ $C_f$ symétrique par rapport à $Oy$.
\item $f$ est \textbf{impaire} $\iff$ $f(-x)=-f(x)$ pour tout $x$ du domaine $\iff$ $C_f$ symétrique par rapport à $O$.
\end{itemize}
\end{propriete}

\begin{exemplebox}[Symétries de $f(x)=x^3-x$]
$f(-x)=(-x)^3-(-x)=-x^3+x=-(x^3-x)=-f(x)$. Donc $f$ est \cle{impaire} : sa courbe est symétrique par rapport à l'origine.
\end{exemplebox}

\courssub{Fonction inverse (bijectivité)}

\begin{definition}[Fonction inverse]
Soit $f:I\to J$ une fonction \textbf{bijective} (strictement monotone sur un intervalle $I$). Il existe une unique fonction $f^{-1}:J\to I$ telle que :
\[
\forall x\in I,\; f^{-1}(f(x))=x \quad\text{et}\quad \forall y\in J,\; f(f^{-1}(y))=y.
\]
La courbe $C_{f^{-1}}$ est l'image de $C_f$ par la symétrie axiale par rapport à la droite $y=x$ (bisectrice du premier quadrant).
\end{definition}

\begin{propriete}[Construction graphique de $f^{-1}$]
Pour tracer $C_{f^{-1}}$ à partir de $C_f$ :
\begin{enumerate}
\item Tracer la droite $y=x$.
\item Pour quelques points $A(x_A,f(x_A))$ de $C_f$, construire leur symétriques $A'(f(x_A),x_A)$ par rapport à $y=x$ (échanger abscisse et ordonnée).
\item Relier les points $A'$ en respectant le sens de variation (si $f$ est croissante, $f^{-1}$ l'est aussi ; si $f$ est décroissante, $f^{-1}$ l'est aussi).
\end{enumerate}
Le domaine de $f^{-1}$ est l'ensemble d'arrivée de $f$ ; l'ensemble d'arrivée de $f^{-1}$ est le domaine de $f$.
\end{propriete}

\begin{exoresolu}[Fonction inverse de $f(x)=2x+3$]
\textbf{Énoncé.} Déterminer l'expression de $f^{-1}$ et tracer les deux courbes.
\tcblower
\textbf{Solution.}
$f$ est affine, strictement croissante sur $\mathbb{R}$, donc bijective de $\mathbb{R}$ sur $\mathbb{R}$.
On résout $y=2x+3$ par rapport à $x$ :
\[
y=2x+3 \;\Longrightarrow\; 2x=y-3 \;\Longrightarrow\; x=\frac{y-3}{2}.
\]
Ainsi $f^{-1}(y)=\frac{y-3}{2}$, c'est-à-dire $f^{-1}(x)=\frac{x-3}{2}$.
Les deux droites sont symétriques par rapport à $y=x$ (pentes inverses $2$ et $1/2$).
\end{exoresolu}

\begin{popfigure}
\begin{tikzpicture}[scale=0.7]
\draw[->] (-2,0)--(5,0) node[right]{$x$};
\draw[->] (0,-2)--(0,5) node[above]{$y$};
\draw[domain=-1:1,smooth,variable=\x,popblue,thick] plot ({\x},{2*\x+3});
\draw[domain=-1:5,smooth,variable=\x,poporange,dashed,thick] plot ({\x},{(\x-3)/2});
\draw[domain=-2:5,popmuted,thin] plot ({\x},{\x});
\node[popblue] at (1,4.5) {$f(x)=2x+3$};
\node[poporange] at (4.5,1) {$f^{-1}(x)=\frac{x-3}{2}$};
\node[popmuted] at (4.5,4.5) {$y=x$};
\end{tikzpicture}
\legende{Courbes de $f$ et $f^{-1}$ symétriques par rapport à $y=x$.}
\end{popfigure}

\courssub{Combinaison de transformations : l'ordre compte}

\begin{attention}[Piège : l'ordre des transformations]
Les transformations ne commutent pas en général. La règle fondamentale distingue les transformations \textbf{verticales} (à l'extérieur de $f$) et \textbf{horizontales} (à l'intérieur de l'argument $x$) :
\begin{itemize}
\item \textbf{Verticales} ($k f(x) + b$) : s'appliquent dans l'\textbf{ordre naturel} (lecture de gauche à droite : homothétie puis translation).
\item \textbf{Horizontales} ($f(kx + c)$) : s'appliquent dans l'\textbf{ordre inverse} de l'écriture. Il faut factoriser le coefficient de $x$ pour voir clair : $f(k(x + c/k))$.
\end{itemize}
\textbf{Exemple piège :} $g(x) = f(2x-6)$.
\begin{enumerate}
\item On factorise : $g(x) = f(2(x-3))$.
\item Opérations sur $x$ (ordre d'écriture) : $\times 2$ puis $-6$ (ou $-3$ puis $\times 2$).
\item Ordre graphique (inverse) :
  \begin{itemize}
  \item Translation horizontale $+3$ (vers la droite) \textbf{en premier}.
  \item Homothétie horizontale de rapport $1/2$ (compression) \textbf{ensuite}.
  \end{itemize}
Si on fait la compression d'abord, la translation suivante est de $3$ sur la figure \emph{comprimée}, ce qui correspond à $6$ sur l'originale : on obtient $f(2(x-3))$ au lieu de $f(2x-6)$ si on ne factorise pas.
\end{enumerate}
\textbf{Règle pratique} : Pour $x \mapsto a f(bx+c)+d$, factoriser $bx+c = b(x+c/b)$. Appliquer : translation horizontale $c/b$, homothétie horizontale $1/b$, homothétie verticale $a$, translation verticale $d$.
\end{attention}

\begin{exemplebox}[Exemple complet : $f(x)=\sqrt{x}\;\to\; g(x)=2\sqrt{x-3}+1$]
\begin{enumerate}
\item \textbf{Translation horizontale} $+3$ : $f_1(x)=\sqrt{x-3}$ (domaine $x\ge3$).
\item \textbf{Homothétie verticale} $\times2$ : $f_2(x)=2\sqrt{x-3}$.
\item \textbf{Translation verticale} $+1$ : $g(x)=2\sqrt{x-3}+1$.
\end{enumerate}
La courbe $C_g$ est obtenue en décalant $C_f$ de $3$ à droite, en l'étirant verticalement par $2$, puis en la montant de $1$.
\end{exemplebox}

\begin{savaistu}[Le miroir de la bisectrice]
La symétrie par rapport à la droite $y=x$ échange les rôles de l'abscisse et de l'ordonnée. C'est pourquoi, dans un repère orthonormé, la courbe de la fonction inverse est exactement le \emph{miroir} de la courbe initiale. Cette propriété est utilisée en optique (loi de la réflexion : l'angle d'incidence égale l'angle de réflexion par rapport à la normale, modélisable par symétrie) et en infographie pour inverser des transformations géométriques (matrices inverses).
\end{savaistu}

\begin{exercice}[Entraînement]
Soit $f(x)=\frac{1}{x}$ définie sur $\mathbb{R}^*$.
\begin{enumerate}
\item Tracer $C_f$ (rappel : hyperbole équilatère, asymptotes $x=0$ et $y=0$).
\item Déterminer et tracer $g(x)=f(x+2)-1$. Préciser les asymptotes.
\item Déterminer et tracer $h(x)=-f(2x)$. Préciser les transformations successives.
\item $f$ est-elle bijective sur $\mathbb{R}^*$ ? Si oui, donner $f^{-1}$ et tracer sa courbe.
\end{enumerate}
\end{exercice}

\begin{corrige}[Exercice n°1]
\begin{enumerate}
\item $C_f$ : deux branches, asymptotes $x=0$ (axe $Oy$) et $y=0$ (axe $Ox$), passant par $(1,1)$ et $(-1,-1)$. Décroissante sur $]-\infty,0[$ et sur $]0,+\infty[$.
\item $g(x)=\frac{1}{x+2}-1$ : translation du vecteur $(-2,-1)$. Asymptotes : $x=-2$ (verticale) et $y=-1$ (horizontale). Domaine : $\mathbb{R}\setminus\{-2\}$.
\item $h(x)=-\frac{1}{2x}$ : homothétie horizontale de rapport $1/2$ (compression vers $Oy$) puis symétrie axiale par rapport à $Ox$ (ou l'inverse). Asymptotes inchangées ($x=0$, $y=0$). Domaine : $\mathbb{R}^*$.
\item $f$ est injective sur $\mathbb{R}^*$ (si $1/x = 1/y \Rightarrow x=y$) et surjective sur $\mathbb{R}^*$ (pour tout $y\neq 0$, $x=1/y$ est antécédent). Donc $f$ est bijective de $\mathbb{R}^*$ sur $\mathbb{R}^*$. $f^{-1}(x)=\frac{1}{x}=f(x)$ : la courbe est invariante par symétrie par rapport à $y=x$ (elle coïncide avec sa propre inverse).
\end{enumerate}
\end{corrige}

\newpage

\coursec{Activité d'intégration}

\coursec{Généralités sur les fonctions numériques : Activité d'intégration -- Facturation ENEO}

\courssub{Objectifs de l'activité}
Cette activité d'intégration vise à mobiliser l'ensemble des savoirs et savoir-faire de la leçon « Généralités sur les fonctions numériques » dans un contexte professionnel authentique. Elle permet à l'élève de :
\begin{itemize}
    \item \textbf{Modéliser} une situation réelle (facturation d'électricité) à l'aide de fonctions numériques (périodique, affine par morceaux).
    \item \textbf{Manipuler} la composition de fonctions pour établir la fonction facture mensuelle.
    \item \textbf{Analyser} les propriétés graphiques et algébriques : périodicité, symétrie (parité après recentrage), variations, extremums.
    \item \textbf{Résoudre} graphiquement et algébriquement une inéquation issue de la composition.
    \item \textbf{Rédiger} une conclusion argumentée, en langage courant et mathématique, sur l'impact de la saisonnalité sur le budget du ménage.
\end{itemize}
\cle{Modélisation} \cle{Composition de fonctions} \cle{Périodicité} \cle{Symétrie centrale} \cle{Résolution d'inéquation} \cle{Contexte professionnel}

\courssub{Situation}
\textbf{Contexte :} À Douala, la société \textbf{ENEO} (Energy of Cameroon) facture l'électricité selon un tarif progressif par tranches de consommation. Un ménage du quartier Bonapriso souhaite anticiper son budget énergie pour l'année 2025. 

Après analyse de ses relevés de compteur sur les trois dernières années, le ménage observe que sa consommation mensuelle moyenne (en kWh) suit un cycle annuel régulier lié aux saisons : plus élevée en saison sèche (climatisation, ventilateurs) et plus basse en saison des pluies. Cette consommation est modélisée par la fonction $C$ définie sur $[0\,;\,12]$ (durée en mois, $t=0$ correspondant au 1er janvier) par :
\[
C(t) = 150 + 30 \sin\left(\frac{2\pi t}{12}\right) = 150 + 30 \sin\left(\frac{\pi t}{6}\right).
\]

Le tarif ENEO applicable aux ménages (tranche sociale) pour l'année 2025 est le suivant :
\begin{itemize}
    \item Abonnement mensuel fixe : \textbf{5\,000 FCFA}.
    \item Prix du kWh pour les 200 premiers kWh : \textbf{50 FCFA/kWh}.
    \item Prix du kWh au-delà de 200 kWh : \textbf{80 FCFA/kWh}.
\end{itemize}
On note $T$ la fonction « coût de l'énergie » (en FCFA) en fonction de la consommation $x$ (en kWh) :
\[
T(x) = 
\begin{cases}
5000 + 50x & \text{si } 0 \le x \le 200 \\[4pt]
5000 + 50 \times 200 + 80(x - 200) = 15000 + 80(x - 200) & \text{si } x > 200
\end{cases}
\]

La facture mensuelle totale $F(t)$ (en FCFA) pour le mois $t$ s'obtient en composant $T$ et $C$ : $F = T \circ C$.

\courssub{Consignes}
\emph{Toutes les réponses doivent être justifiées par un calcul, une lecture graphique ou une propriété du cours.}

\begin{enumerate}
    \item \textbf{Étude de la consommation $C(t)$.}
    \begin{enumerate}
        \item Tracer la courbe $\mathcal{C}_C$ de la fonction $C$ sur l'intervalle $[0\,;\,12]$ dans un repère orthonormé (1 cm pour 1 mois en abscisses, 1 cm pour 10 kWh en ordonnées).
        \item Déterminer la \cle{période} de $C$ et l'interpréter dans le contexte.
        \item Étudier la \cle{parité} de $C$ sur $[0\,;\,12]$. La fonction est-elle paire ? Impaire ? Mettre en évidence une symétrie (axiale ou centrale) après un recentrage convenable de l'axe des abscisses. Préciser le centre ou l'axe de symétrie.
        \item Donner les valeurs minimale et maximale de la consommation et les mois correspondants.
    \end{enumerate}

    \item \textbf{Expression de la facture mensuelle $F(t) = T(C(t))$.}
    \begin{enumerate}
        \item Déterminer l'ensemble des valeurs prises par $C(t)$ pour $t \in [0\,;\,12]$. En déduire quelle branche de la fonction $T$ est utilisée pour tout $t \in [0\,;\,12]$.
        \item Établir l'expression littérale simplifiée de $F(t)$ sur $[0\,;\,12]$.
        \item $F$ est-elle une fonction périodique ? Si oui, quelle est sa période ?
    \end{enumerate}

    \item \textbf{Analyse graphique de la facture $F(t)$.}
    \begin{enumerate}
        \item Tracer la courbe $\mathcal{C}_F$ de la fonction $F$ sur $[0\,;\,12]$ dans un nouveau repère (1 cm pour 1 mois, 1 cm pour 1\,000 FCFA).
        \item À l'aide de la courbe (ou par le calcul), déterminer les mois où la facture dépasse \textbf{15\,000 FCFA}.
    \end{enumerate}

    \item \textbf{Synthèse et conclusion.}
    Rédiger un court paragraphe (5 à 8 lignes) adressé au chef de famille, expliquant l'évolution de la facture sur l'année, l'impact de la saisonnalité, et confirmant si le seuil de 15\,000 FCFA est atteint. Utiliser le vocabulaire mathématique approprié (période, amplitude, variation, composition).
\end{enumerate}

\courssub{Ressources à mobiliser}
\begin{propriete}[Fonction sinus : période et amplitude]
Pour $f(t) = a + b \sin(\omega t + \varphi)$ avec $b>0$ :
\begin{itemize}
    \item Période : $T = \frac{2\pi}{\omega}$.
    \item Amplitude : $b$. Valeurs min/max : $a-b$ et $a+b$.
    \item $\sin$ est une fonction \textbf{impaire} ($2\pi$-périodique) : $\sin(-X) = -\sin(X)$.
\end{itemize}
\end{propriete}

\begin{propriete}[Composition de fonctions]
Soit $f : I \to J$ et $g : J \to \mathbb{R}$. La composée $g \circ f$ est définie sur $I$ par $(g \circ f)(x) = g(f(x))$.
\emph{Attention :} L'ensemble de définition de $g \circ f$ est $\{ x \in I \mid f(x) \in J \}$.
\end{propriete}

\begin{propriete}[Parité et symétries]
\begin{itemize}
    \item $f$ \textbf{paire} sur $[-a,a] \iff \forall x, f(-x)=f(x) \iff$ symétrie axiale par rapport à l'axe des ordonnées ($x=0$).
    \item $f$ \textbf{impaire} sur $[-a,a] \iff \forall x, f(-x)=-f(x) \iff$ symétrie centrale par rapport à l'origine $O(0,0)$.
    \item \textbf{Recentrage :} Pour étudier la symétrie d'une fonction définie sur $[0, T]$, on pose $u = t - \frac{T}{2}$. La fonction $h(u) = f(u+\frac{T}{2})$ est définie sur $[-\frac{T}{2}, \frac{T}{2}]$. On étudie alors la parité de $h$.
\end{itemize}
\end{propriete}

\begin{methode}[Résolution d'inéquation $f(x) > k$]
\begin{enumerate}
    \item Résoudre l'équation $f(x) = k$.
    \item Étudier le signe de $f(x) - k$ (tableau de signes ou lecture graphique).
    \item Conclure par l'ensemble solution sous forme d'intervalle ou d'union d'intervalles.
\end{enumerate}
\end{methode}

\courssub{Démarche et corrigé commenté}
\begin{exoresolu}{Corrigé détaillé de l'activité d'intégration}
\paragraph{1. Étude de la consommation $C(t)$.}
\begin{enumerate}
    \item \textbf{Tracé de $\mathcal{C}_C$.} 
    La fonction $C(t) = 150 + 30 \sin(\frac{\pi t}{6})$ est une sinusoïde.
    \begin{itemize}
        \item Période $T = \frac{2\pi}{\pi/6} = 12$ mois. Cela correspond au cycle annuel.
        \item Amplitude $= 30$ kWh. Valeur moyenne (axe d'oscillation) $= 150$ kWh.
        \item Minimum : $150 - 30 = 120$ kWh. Maximum : $150 + 30 = 180$ kWh.
        \item Points remarquables ($t$ en mois, $C$ en kWh) :
        \begin{center}
        \begin{tabular}{|c|c|c|c|c|c|c|c|}\hline
        $t$ & 0 & 3 & 6 & 9 & 12 \\ \hline
        $\frac{\pi t}{6}$ & 0 & $\pi/2$ & $\pi$ & $3\pi/2$ & $2\pi$ \\ \hline
        $\sin(\dots)$ & 0 & 1 & 0 & -1 & 0 \\ \hline
        $C(t)$ & 150 & 180 & 150 & 120 & 150 \\ \hline
        \end{tabular}
        \end{center}
    \end{itemize}
    \begin{popfigure}
    \begin{tikzpicture}
    \begin{axis}[
        width=\linewidth, height=7cm,
        axis lines=middle,
        xlabel={$t$ (mois)}, ylabel={$C(t)$ (kWh)},
        xmin=-0.5, xmax=12.5, ymin=110, ymax=190,
        xtick={0,3,6,9,12}, ytick={120,150,180},
        grid=major, grid style={popmuted},
        domain=0:12, samples=200,
        tick label style={font=\small},
        label style={font=\small}
    ]
    \addplot[popcurve, popblue, thick] {150 + 30*sin(deg(pi*x/6))};
    \addplot[dashed, popmuted] coordinates {(0,150) (12,150)};
    \node[popblue, anchor=south west] at (axis cs: 1, 175) {$\mathcal{C}_C$};
    \node[popmuted, anchor=north east, font=\tiny] at (axis cs: 11.5, 150) {$y=150$ (moyenne)};
    \end{axis}
    \end{tikzpicture}
    \legende{Courbe $\mathcal{C}_C$ de la consommation mensuelle $C(t)=150+30\sin(\pi t/6)$ sur $[0,12]$.}
    \end{popfigure}

    \item \textbf{Période.} $T = 12$ mois. \emph{Interprétation :} Le cycle de consommation se répète chaque année. Le ménage a le même profil de consommation d'une année sur l'autre.

    \item \textbf{Parité et symétrie (après recentrage).}
    $C$ est définie sur $[0,12]$, intervalle non symétrique par rapport à 0. On ne peut pas parler de parité (paire/impaire) directement sur cet intervalle au sens strict de la définition.
    \emph{Recentrage :} On pose $u = t - 6$ (décalage de la moitié de la période). Alors $t = u+6$ avec $u \in [-6, 6]$.
    On définit $h(u) = C(u+6) = 150 + 30 \sin\left(\frac{\pi(u+6)}{6}\right) = 150 + 30 \sin\left(\frac{\pi u}{6} + \pi\right)$.
    Or $\sin(X+\pi) = -\sin X$. Donc $h(u) = 150 - 30 \sin\left(\frac{\pi u}{6}\right)$.
    La fonction $u \mapsto \sin(\frac{\pi u}{6})$ est \textbf{impaire} sur $[-6,6]$.
    Donc $u \mapsto -30 \sin(\frac{\pi u}{6})$ est impaire.
    $h(u) = 150 + [\text{fonction impaire}]$.
    \begin{aretenir}[Symétrie centrale]
    La courbe $\mathcal{C}_C$ admet un \textbf{centre de symétrie} de coordonnées $(6\,;\, 150)$ (mois de juillet, consommation moyenne).
    En effet, $h(-u) + h(u) = 300 = 2 \times 150$.
    Il n'y a \textbf{pas d'axe de symétrie verticale} (fonction non paire après recentrage).
    \end{aretenir}

    \item \textbf{Extremums.}
    \begin{itemize}
        \item Maximum $180$ kWh atteint pour $t=3$ (fin mars / avril, saison sèche chaude).
        \item Minimum $120$ kWh atteint pour $t=9$ (fin septembre / octobre, saison des pluies).
    \end{itemize}
\end{enumerate}

\paragraph{2. Expression de la facture mensuelle $F(t) = T(C(t))$.}
\begin{enumerate}
    \item \textbf{Ensemble des valeurs de $C$.}
    Pour $t \in [0,12]$, $C(t) \in [120\,;\, 180]$.
    Comme $180 < 200$, \textbf{toutes les valeurs de $C(t)$ appartiennent à la première tranche} ($x \le 200$).
    \begin{attention}[Piège fréquent]
    Ne pas écrire la seconde branche de $T$ « par principe » parce qu'elle existe dans l'énoncé. Il faut \textbf{toujours vérifier} si l'ensemble d'arrivée de la fonction intérieure ($C$) est inclus dans l'ensemble de définition de la branche considérée. Ici, $C([0,12]) = [120,180] \subset [0,200]$.
    \end{attention}
    On utilise donc uniquement $T(x) = 5000 + 50x$.

    \item \textbf{Expression de $F$.}
    Pour tout $t \in [0,12]$ :
    \begin{align*}
    F(t) &= T(C(t)) = 5000 + 50 \times C(t) \\
         &= 5000 + 50 \left(150 + 30 \sin\left(\frac{\pi t}{6}\right)\right) \\
         &= 5000 + 7500 + 1500 \sin\left(\frac{\pi t}{6}\right) \\
         &= \boxed{12500 + 1500 \sin\left(\frac{\pi t}{6}\right)} \quad \text{(FCFA)}.
    \end{align*}

    \item \textbf{Périodicité de $F$.}
    $F$ est la composition de $T$ (affine sur l'intervalle considéré) et de $C$ (périodique de période 12).
    $F$ est \textbf{périodique de période 12 mois}. C'est une sinusoïde de même période que $C$.
\end{enumerate}

\paragraph{3. Analyse graphique de la facture $F(t)$.}
\begin{enumerate}
    \item \textbf{Tracé de $\mathcal{C}_F$.}
    $F(t) = 12500 + 1500 \sin(\frac{\pi t}{6})$.
    \begin{itemize}
        \item Axe d'oscillation : $y = 12\,500$ FCFA.
        \item Amplitude : $1\,500$ FCFA.
        \item Minimum : $12\,500 - 1\,500 = 11\,000$ FCFA (à $t=9$).
        \item Maximum : $12\,500 + 1\,500 = 14\,000$ FCFA (à $t=3$).
        \item Période : 12 mois.
    \end{itemize}
    \begin{popfigure}
    \begin{tikzpicture}
    \begin{axis}[
        width=\linewidth, height=7cm,
        axis lines=middle,
        xlabel={$t$ (mois)}, ylabel={$F(t)$ (FCFA)},
        xmin=-0.5, xmax=12.5, ymin=10500, ymax=14500,
        xtick={0,3,6,9,12}, 
        ytick={11000,12500,14000},
        grid=major, grid style={popmuted},
        domain=0:12, samples=200,
        tick label style={font=\small, /pgf/number format/set thousands separator={\,}, /pgf/number format/fixed},
        label style={font=\small}
    ]
    \addplot[popcurve, poporange, thick] {12500 + 1500*sin(deg(pi*x/6))};
    \addplot[dashed, popmuted] coordinates {(0,12500) (12,12500)};
    \addplot[dashed, popred] coordinates {(0,15000) (12,15000)};
    \node[poporange, anchor=south west] at (axis cs: 1, 13800) {$\mathcal{C}_F$};
    \node[popmuted, anchor=north east, font=\tiny] at (axis cs: 11.5, 12500) {$y=12\,500$ (moyenne)};
    \node[popred, anchor=south west, font=\tiny] at (axis cs: 0.5, 15050) {Seuil 15\,000 FCFA};
    \end{axis}
    \end{tikzpicture}
    \legende{Courbe $\mathcal{C}_F$ de la facture mensuelle $F(t)=12500+1500\sin(\pi t/6)$. Le seuil 15\,000 FCFA (ligne rouge) n'est jamais atteint.}
    \end{popfigure}

    \item \textbf{Mois où $F(t) > 15\,000$.}
    \emph{Par lecture graphique :} La courbe $\mathcal{C}_F$ reste strictement en dessous de la ligne $y=15\,000$ (maximum à 14\,000).
    \emph{Par le calcul :} On résout $F(t) > 15000$.
    \[
    12500 + 1500 \sin\left(\frac{\pi t}{6}\right) > 15000
    \iff 1500 \sin\left(\frac{\pi t}{6}\right) > 2500
    \iff \sin\left(\frac{\pi t}{6}\right) > \frac{2500}{1500} = \frac{5}{3} \approx 1,66.
    \]
    Or $\forall X \in \mathbb{R},\; \sin X \le 1$. L'inéquation n'a \textbf{aucune solution}.
    \begin{aretenir}[Conclusion Q3.2]
    \textbf{Aucun mois de l'année} la facture ne dépasse 15\,000 FCFA. La facture maximale est 14\,000 FCFA (en mars/avril).
    \end{aretenir}
\end{enumerate}

\paragraph{4. Synthèse et conclusion (Exemple de rédaction attendue).}
\begin{exemplebox}[Rédaction finale pour le chef de famille]
« Monsieur, l'étude de notre facture d'électricité sur l'année 2025 montre une variation parfaitement périodique de période 12 mois, calquée sur la saisonnalité. La consommation $C(t)$ oscille entre 120 kWh (saison des pluies, septembre-octobre) et 180 kWh (saison sèche, mars-avril) autour d'une moyenne de 150 kWh, avec une symétrie centrale par rapport au mois de juillet (centre $(6;150)$). 

Grâce à la composition avec le tarif ENEO, la facture mensuelle $F(t) = 12\,500 + 1\,500 \sin(\pi t/6)$ varie entre 11\,000 FCFA et 14\,000 FCFA. Le seuil de 15\,000 FCFA n'est \textbf{jamais atteint} car notre consommation reste inférieure à 200 kWh toute l'année, nous maintenant dans la première tranche tarifaire (50 FCFA/kWh). L'amplitude de la facture (3\,000 FCFA d'écart max/min) reflète directement l'amplitude de consommation (60 kWh) amplifiée par le prix unitaire. Le budget annuel prévisionnel est donc stable et maîtrisé. »
\end{exemplebox}
\end{exoresolu}

\courssub{Bilan de l'activité}
Cette activité a permis de mettre en évidence les points clés suivants, conformes aux attendus du Probatoire technique :
\begin{itemize}
    \item \textbf{Modélisation réelle :} Passage d'un énoncé verbal (tarif par morceaux, consommation saisonnière) à un modèle mathématique rigoureux (fonctions $C$, $T$, $F$).
    \item \textbf{Composition de fonctions :} Nécessité de vérifier l'ensemble des valeurs de la fonction intérieure ($C$) pour choisir la bonne branche de la fonction extérieure ($T$). C'est un point \cle{exigible à l'examen} souvent source d'erreurs.
    \item \textbf{Propriétés des fonctions trigonométriques :} Identification de la période, de l'amplitude, exploitation de la parité de la fonction sinus après recentrage (symétrie centrale), détermination des extremums.
    \item \textbf{Lien algèbre / graphique :} Le tracé des courbes $\mathcal{C}_C$ et $\mathcal{C}_F$ valide les calculs algébriques (période, extremums, position relative au seuil 15\,000).
    \item \textbf{Résolution d'inéquation :} Combinaison du raisonnement sur les bornes de la fonction sinus ($\sin x \le 1$) et de la lecture graphique pour conclure à l'ensemble vide.
    \item \textbf{Communication technique :} Rédaction d'une conclusion synthétique, justifiée, utilisant le vocabulaire précis (période, amplitude, composition, tranche tarifaire, symétrie centrale) adaptée à un destinataire non mathématicien.
\end{itemize}
\begin{savaistu}[Le saviez-vous ?]
Au Cameroun, la structure tarifaire d'ENEO comprend effectivement une composante fixe (abonnement) et une composante variable par tranches de consommation (tranches sociales, résidentielles, industrielles). La modélisation par fonctions affines par morceaux est l'outil standard de la \textbf{tarification de l'énergie} et des \textbf{télécommunications} (forfaits data, appels). Maîtriser la composition de fonctions permet aux techniciens (électriciens, gestionnaires) de simuler des factures, de conseiller les clients sur l'optimisation de leur consommation (ex: déplacer des consommations hors pointe) et de vérifier la conformité des factures émises.
\end{savaistu}

\newpage

\coursec{Méthodes et exercices résolus}

\coursec{Rappels et vocabulaire sur les fonctions numériques}

\begin{definition}[Fonction numérique, ensemble de définition, image, antécédent]
Une \cle{fonction numérique} $f$ est une application qui associe à tout nombre réel $x$ appartenant à un ensemble $D_f$ (appelé \cle{ensemble de définition} de $f$) un unique nombre réel noté $f(x)$ (appelé \cle{image} de $x$ par $f$).
Si $y = f(x)$, on dit que $x$ est un \cle{antécédent} de $y$ par $f$.
L'ensemble des images est noté $f(D_f)$ ou $\text{Im}(f)$.
\end{definition}

\begin{propriete}[Fonctions de référence en Première technique]
On doit connaître l'expression, l'ensemble de définition et la forme de la courbe des fonctions suivantes :
\begin{itemize}
\item \textbf{Fonction affine :} $f(x) = ax + b$ \quad ($D_f = \mathbb{R}$, droite).
\item \textbf{Fonction carré :} $f(x) = x^2$ \quad ($D_f = \mathbb{R}$, parabole sommet $O$, axe $Oy$).
\item \textbf{Fonction racine carrée :} $f(x) = \sqrt{x}$ \quad ($D_f = [0\,;\,+\infty[$, demi-parabole).
\item \textbf{Fonction inverse / homographique simple :} $f(x) = \dfrac{k}{x}$ \quad ($D_f = \mathbb{R}^*$, hyperbole, centre $O$).
\item \textbf{Fonction homographique générale :} $f(x) = \dfrac{ax+b}{cx+d}$ \quad ($D_f = \mathbb{R} \setminus \{-d/c\}$).
\end{itemize}
\end{propriete}

\begin{methode}[Trouver l'ensemble de définition $D_f$]
Pour déterminer l'ensemble de définition $D_f$ d'une fonction numérique $f$ donnée par une expression algébrique :
\begin{enumerate}
\item Repérer les \cle{contraintes} dans l'expression de $f(x)$ :
\begin{itemize}
\item un \textbf{dénominateur} : il doit être \textbf{non nul} ($q(x) \neq 0$) ;
\item une \textbf{racine carrée} (ou racine d'indice pair) $\sqrt{u(x)}$ : le radicande doit vérifier $u(x) \geq 0$.
\end{itemize}
\item Traduire chaque contrainte par une équation ou une inéquation et la résoudre dans $\mathbb{R}$.
\item \textbf{Intersecter} les ensembles de solutions (toutes les conditions doivent être vérifiées simultanément).
\item Écrire $D_f$ sous forme d'intervalles ou de réunion d'intervalles.
\item Conclure par une phrase : « $f$ est définie sur $D_f = \ldots$ ».
\end{enumerate}
\textbf{Réflexe :} pas de dénominateur ni de racine carrée (indice pair) $\Rightarrow$ $D_f = \mathbb{R}$.
\end{methode}

\begin{exoresolu}[Ensemble de définition — type Probatoire]
On considère la fonction $f$ définie par $f(x) = \dfrac{2x+1}{x-3} + \sqrt{x+2}$.
Déterminer l'ensemble de définition $D_f$ de $f$.
\tcblower
L'expression de $f$ contient \textbf{deux contraintes} : un dénominateur et une racine carrée.

\textbf{Contrainte 1 : le dénominateur.} On doit avoir $x - 3 \neq 0$, c'est-à-dire $x \neq 3$.

\textbf{Contrainte 2 : la racine carrée.} On doit avoir $x + 2 \geq 0$, c'est-à-dire $x \geq -2$.

\textbf{Contraintes simultanées (intersection).} Les deux conditions doivent être vérifiées en même temps :
\[ x \geq -2 \quad \text{et} \quad x \neq 3. \]
Sur la droite réelle, on garde donc l'intervalle $[-2\,;\,+\infty[$ privé de la valeur $3$ :
\[ D_f = [-2\,;\,3[ \;\cup\; ]3\,;\,+\infty[. \]
\textbf{Conclusion :} la fonction $f$ est définie sur $D_f = [-2\,;\,3[ \cup ]3\,;\,+\infty[$.
\end{exoresolu}

\coursec{Opérations sur les fonctions et composition}

\begin{definition}[Opérations sur les fonctions]
Soient $f$ et $g$ deux fonctions numériques d'ensembles de définition $D_f$ et $D_g$.
On définit les fonctions \textbf{somme}, \textbf{produit} et \textbf{quotient} sur l'ensemble $D = D_f \cap D_g$ (pour le quotient, on exclut en plus les zéros du dénominateur) par :
\begin{align*}
(f+g)(x) &= f(x) + g(x), \\
(fg)(x) &= f(x) \cdot g(x), \\
\left(\frac{f}{g}\right)(x) &= \frac{f(x)}{g(x)} \quad \text{pour } g(x) \neq 0.
\end{align*}
\end{definition}

\begin{methode}[Calculer une composée $g \circ f$]
Pour déterminer l'expression de la composée $g \circ f$ :
\begin{enumerate}
\item Rappeler la définition : $(g \circ f)(x) = g\bigl(f(x)\bigr)$ : on applique d'abord $f$, puis $g$.
\item Remplacer, dans l'expression de $g$, \textbf{chaque} occurrence de la variable par l'expression complète de $f(x)$, entre parenthèses.
\item Développer et réduire soigneusement.
\item \textbf{Préciser l'ensemble de définition :} $D_{g \circ f} = \{ x \in D_f \mid f(x) \in D_g \}$.
\end{enumerate}
\textbf{Attention :} en général $g \circ f \neq f \circ g$ : la composition n'est pas commutative.
\end{methode}

\begin{exoresolu}[Composition d'applications — type Probatoire]
Soient $f$ et $g$ les applications de $\mathbb{R}$ dans $\mathbb{R}$ définies par $f(x) = 2x - 1$ et $g(x) = x^2 + 3$.
\begin{enumerate}
\item Déterminer $(g \circ f)(x)$ et $(f \circ g)(x)$ ainsi que leurs ensembles de définition.
\item Calculer $(g \circ f)(2)$. A-t-on $g \circ f = f \circ g$ ?
\end{enumerate}
\tcblower
\textbf{1. Calcul de $g \circ f$.} Par définition, $(g \circ f)(x) = g\bigl(f(x)\bigr) = g(2x-1)$. On remplace $x$ par $(2x-1)$ dans $g(x) = x^2 + 3$ :
\begin{align*}
(g \circ f)(x) &= (2x-1)^2 + 3 \\
&= 4x^2 - 4x + 1 + 3 \\
&= 4x^2 - 4x + 4.
\end{align*}
Comme $D_f = \mathbb{R}$ et $D_g = \mathbb{R}$, pour tout $x \in \mathbb{R}$, $f(x) \in \mathbb{R} = D_g$. Donc $D_{g \circ f} = \mathbb{R}$.

\textbf{Calcul de $f \circ g$.} Par définition, $(f \circ g)(x) = f\bigl(g(x)\bigr) = f(x^2+3)$. On remplace $x$ par $(x^2+3)$ dans $f(x) = 2x - 1$ :
\begin{align*}
(f \circ g)(x) &= 2(x^2+3) - 1 \\
&= 2x^2 + 6 - 1 \\
&= 2x^2 + 5.
\end{align*}
De même, $D_{f \circ g} = \mathbb{R}$.

\textbf{2. Valeur numérique.} $(g \circ f)(2) = 4(2)^2 - 4(2) + 4 = 16 - 8 + 4 = 12$.
Vérification directe : $f(2) = 2(2) - 1 = 3$, puis $g(3) = 3^2 + 3 = 12$. On retrouve bien $12$.

Comme $4x^2 - 4x + 4 \neq 2x^2 + 5$ (par exemple en $x = 0$ : $(g\circ f)(0) = 4$ et $(f \circ g)(0) = 5$), on a $g \circ f \neq f \circ g$.

\textbf{Conclusion :} $(g \circ f)(x) = 4x^2 - 4x + 4$ sur $\mathbb{R}$, $(f \circ g)(x) = 2x^2 + 5$ sur $\mathbb{R}$ et $(g \circ f)(2) = 12$ ; la composition n'est pas commutative.
\end{exoresolu}

\coursec{Parité, symétrie axiale et centrale}

\begin{definition}[Parité d'une fonction]
Soit $f$ une fonction numérique d'ensemble de définition $D_f$.
\begin{itemize}
\item $f$ est \textbf{paire} si $D_f$ est symétrique par rapport à $0$ (i.e. $x \in D_f \Rightarrow -x \in D_f$) et si $\forall x \in D_f,\; f(-x) = f(x)$.
\item $f$ est \textbf{impaire} si $D_f$ est symétrique par rapport à $0$ et si $\forall x \in D_f,\; f(-x) = -f(x)$.
\item Sinon, $f$ n'est \textbf{ni paire ni impaire}.
\end{itemize}
\end{definition}

\begin{propriete}[Conséquences graphiques — Exigibles au Probatoire]
\begin{itemize}
\item Si $f$ est \textbf{paire} : sa courbe $\mathcal{C}_f$ est symétrique par rapport à l'axe des ordonnées (axe $Oy$).
\item Si $f$ est \textbf{impaire} : sa courbe $\mathcal{C}_f$ est symétrique par rapport à l'origine $O$ (symétrie centrale).
\end{itemize}
\end{propriete}

\begin{methode}[Montrer qu'une fonction est paire, impaire ou ni l'une ni l'autre]
\begin{enumerate}
\item \textbf{Étape indispensable :} vérifier que $D_f$ est \cle{symétrique par rapport à 0}, c'est-à-dire : si $x \in D_f$ alors $-x \in D_f$. Sinon, la fonction n'est ni paire ni impaire (arrêter l'étude).
\item Calculer $f(-x)$ en remplaçant $x$ par $-x$, puis simplifier.
\item Comparer avec $f(x)$ et $-f(x)$ :
\begin{itemize}
\item si $f(-x) = f(x)$ pour tout $x \in D_f$ : $f$ est \textbf{paire} ;
\item si $f(-x) = -f(x)$ pour tout $x \in D_f$ : $f$ est \textbf{impaire} ;
\item sinon : $f$ n'est ni paire ni impaire.
\end{itemize}
\item Conclure en citant la conséquence graphique.
\end{enumerate}
\end{methode}

\begin{exoresolu}[Parité — type Probatoire]
On considère les fonctions définies sur $\mathbb{R}$ par :
\[ f(x) = \frac{x^2}{x^2 + 1} \qquad \text{et} \qquad g(x) = x^3 - 2x. \]
Étudier la parité de $f$ et de $g$, puis préciser les éléments de symétrie de leurs courbes.
\tcblower
\textbf{Ensembles de définition.} Pour $f$ : $x^2 + 1 \geq 1 > 0$, donc le dénominateur ne s'annule jamais et $D_f = \mathbb{R}$. Pour $g$ : polynôme, donc $D_g = \mathbb{R}$. Les deux ensembles sont symétriques par rapport à $0$ : l'étude de parité a un sens.

\textbf{Parité de $f$.} Pour tout réel $x$ :
\[ f(-x) = \frac{(-x)^2}{(-x)^2 + 1} = \frac{x^2}{x^2 + 1} = f(x). \]
Donc $f$ est \textbf{paire}. Sa courbe représentative est symétrique par rapport à l'axe des ordonnées.

\textbf{Parité de $g$.} Pour tout réel $x$ :
\begin{align*}
g(-x) &= (-x)^3 - 2(-x) = -x^3 + 2x \\
&= -(x^3 - 2x) = -g(x).
\end{align*}
Donc $g$ est \textbf{impaire}. Sa courbe représentative est symétrique par rapport à l'origine du repère.

\textbf{Conclusion :} $f$ est paire (symétrie par rapport à l'axe des ordonnées) et $g$ est impaire (symétrie par rapport à l'origine).
\end{exoresolu}

\coursec{Périodicité}

\begin{definition}[Fonction périodique]
Une fonction $f$ est dite \textbf{périodique} s'il existe un nombre réel $T > 0$ tel que :
\begin{itemize}
\item $D_f$ est invariant par translation de $T$ : $x \in D_f \Leftrightarrow x+T \in D_f$ ;
\item $\forall x \in D_f,\; f(x+T) = f(x)$.
\end{itemize}
Le plus petit $T > 0$ vérifiant cette propriété est appelé \cle{période fondamentale} (ou simplement période).
\end{definition}

\begin{methode}[Montrer qu'une fonction est périodique de période $T$]
\begin{enumerate}
\item Vérifier que $D_f$ est invariant par translation de $T$ : si $x \in D_f$ alors $x + T \in D_f$ (et $x-T \in D_f$).
\item Calculer $f(x + T)$ en remplaçant $x$ par $x + T$.
\item Simplifier en utilisant les propriétés connues (ex. $\cos(x + 2\pi) = \cos x$, $\sin(x + 2\pi) = \sin x$, $\tan(x+\pi)=\tan x$).
\item Si $f(x + T) = f(x)$ pour tout $x \in D_f$, conclure : « $f$ est périodique de période $T$ ».
\item Conséquence pratique : il suffit d'étudier $f$ sur un intervalle d'amplitude $T$ (ex. $[0\,;\,T]$ ou $[-T/2\,;\,T/2]$), puis de reproduire la courbe par translations de vecteur $kT\,\vec{i}$ ($k \in \mathbb{Z}$).
\end{enumerate}
\end{methode}

\begin{exoresolu}[Périodicité — type Probatoire]
Soit $f$ la fonction définie sur $\mathbb{R}$ par $f(x) = \sin x + \cos(2x)$.
Montrer que $f$ est périodique de période $2\pi$.
\tcblower
$D_f = \mathbb{R}$, qui est invariant par toute translation : la condition sur l'ensemble de définition est vérifiée.

Calculons $f(x + 2\pi)$ pour tout réel $x$ :
\[ f(x + 2\pi) = \sin(x + 2\pi) + \cos\bigl(2(x + 2\pi)\bigr) = \sin(x + 2\pi) + \cos(2x + 4\pi). \]
Or les fonctions sinus et cosinus sont périodiques de période $2\pi$, donc :
\[ \sin(x + 2\pi) = \sin x \quad \text{et} \quad \cos(2x + 4\pi) = \cos(2x + 2 \times 2\pi) = \cos(2x). \]
Par conséquent :
\[ f(x + 2\pi) = \sin x + \cos(2x) = f(x). \]
\textbf{Conclusion :} pour tout réel $x$, $f(x + 2\pi) = f(x)$, donc $f$ est périodique de période $2\pi$. Il suffit donc de l'étudier sur $[0\,;\,2\pi]$ (ou $[-\pi\,;\,\pi]$) pour connaître toute sa courbe.
\end{exoresolu}

\coursec{Courbe d'une fonction : tracé et analyse}

\begin{definition}[Courbe représentative]
Dans un repère orthonormé $(O;\vec{i},\vec{j})$, la \cle{courbe représentative} d'une fonction $f$, notée $\mathcal{C}_f$, est l'ensemble des points $M(x\,;\,y)$ tels que $x \in D_f$ et $y = f(x)$.
\end{definition}

\begin{methode}[Tracer la courbe $\mathcal{C}_f$ d'une fonction]
\begin{enumerate}
\item Déterminer $D_f$.
\item \textbf{Étudier la parité et la périodicité} (si applicable) pour réduire l'intervalle d'étude.
\item Construire un \textbf{tableau de valeurs} : choisir des valeurs de $x$ dans $D_f$ (zéros, valeurs remarquables, bornes de $D_f$) et calculer $f(x)$.
\item \textbf{Étudier le signe} de $f(x)$ (tableau de signes) pour savoir où la courbe est au-dessus / en-dessous de l'axe des abscisses.
\item \textbf{Étudier les variations} (tableau de variations) : sens de variation (croissante/décroissante), extremums. \emph{En Première technique, on utilise les variations connues des fonctions de référence ou un tableau de valeurs serré ; la dérivée n'est pas au programme.}
\item Repérer les \textbf{points particuliers} : intersections avec les axes ($f(x)=0$ et $f(0)$), sommet (parabole), centre d'hyperbole, asymptotes (pour homographiques).
\item Placer les points dans le repère et relier en respectant les variations et la forme connue.
\end{enumerate}
\end{methode}

\begin{propriete}[Tableau de variations — Modèle de rédaction]
Pour une fonction $f$ définie sur un intervalle $I$, le tableau de variations comporte 3 lignes ($x$, $f'(x)$ ou sens de variation, $f(x)$) et autant de colonnes que de valeurs remarquables + bornes.
\begin{center}
\begin{tabular}{|c|ccccc|}\hline
$x$ & $a$ & & $c$ & & $b$ \\ \hline
Sens de variation & & $\nearrow$ & & $\searrow$ & \\ \hline
$f(x)$ & $f(a)$ & $\nearrow$ & $f(c)$ & $\searrow$ & $f(b)$ \\ \hline
\end{tabular}
\end{center}
\emph{Note :} En l'absence de dérivée, on écrit « croissante » / « décroissante » ou on trace les flèches $\nearrow$ / $\searrow$ sur la ligne centrale.
\end{propriete}

\begin{exoresolu}[Tracé de courbe — type Probatoire]
Soit $f$ définie sur $]-2\,;\,+\infty[$ par $f(x) = \sqrt{x+2} - 1$.
\begin{enumerate}
\item Étudier le signe de $f(x)$.
\item Dresser le tableau de variations de $f$ sur $D_f$.
\item Tracer $\mathcal{C}_f$ dans un repère orthonormé.
\end{enumerate}
\tcblower
\textbf{1. Signe.} $f(x) \geq 0 \Leftrightarrow \sqrt{x+2} \geq 1 \Leftrightarrow x+2 \geq 1 \Leftrightarrow x \geq -1$.
$f(x) = 0$ pour $x = -1$.
Tableau de signes :
\begin{center}
\begin{tabular}{|c|ccccc|}\hline
$x$ & $-2$ & & $-1$ & & $+\infty$ \\ \hline
$f(x)$ & & $-$ & $0$ & $+$ & \\ \hline
\end{tabular}
\end{center}

\textbf{2. Variations.} La fonction $x \mapsto x+2$ est croissante sur $]-2\,;\,+\infty[$. La fonction racine carrée est croissante sur $[0\,;\,+\infty[$. La composition $x \mapsto \sqrt{x+2}$ est donc croissante sur $]-2\,;\,+\infty[$. En retranchant $1$, $f$ reste croissante.
Tableau de variations :
\begin{center}
\begin{tabular}{|c|ccc|}\hline
$x$ & $-2$ & & $+\infty$ \\ \hline
$f(x)$ & $-1$ & $\nearrow$ & $+\infty$ \\ \hline
\end{tabular}
\end{center}

\textbf{3. Tracé.} Points remarquables : $A(-2\,;\,-1)$ (borne de $D_f$, extrémité de courbe), $B(-1\,;\,0)$ (intersection avec $Ox$), $C(2\,;\,1)$, $D(7\,;\,2)$. La courbe est la demi-parabole $y = \sqrt{x}$ translatée de vecteur $(-2\,;\,-1)$.
\end{exoresolu}

\coursec{Fonctions associées à une fonction donnée}

\begin{methode}[Passer de $\mathcal{C}_f$ aux courbes des fonctions associées — Transformations géométriques]
À partir de la courbe $\mathcal{C}_f$ de $f$, on obtient les courbes des fonctions associées par les transformations géométriques suivantes (à connaître par cœur) :
\begin{enumerate}
\item $y = f(x) + k$ : \textbf{translation} de vecteur $k\,\vec{j}$ (verticale, vers le haut si $k>0$).
\item $y = f(x - a)$ : \textbf{translation} de vecteur $a\,\vec{i}$ (horizontale, vers la \textbf{droite} si $a>0$).
\item $y = -f(x)$ : \textbf{symétrie} par rapport à l'axe des abscisses ($Ox$).
\item $y = f(-x)$ : \textbf{symétrie} par rapport à l'axe des ordonnées ($Oy$).
\item $y = a\,f(x)$ ($a \in \mathbb{R}^*$) : \textbf{homothety} de rapport $a$ par rapport à l'axe des abscisses (étirement/rétrécissement vertical). Si $a<0$, symétrie par rapport à $Ox$ composée.
\item $y = f(b\,x)$ ($b \in \mathbb{R}^*$) : \textbf{homothety} de rapport $1/b$ par rapport à l'axe des ordonnées (étirement/rétrécissement horizontal). Si $b<0$, symétrie par rapport à $Oy$ composée.
\item $y = |f(x)|$ : on conserve la partie de $\mathcal{C}_f$ au-dessus de $Ox$ et on \textbf{réfléchit} la partie située en dessous (symétrie par rapport à $Ox$ de la partie négative).
\item $y = f(|x|)$ : on conserve la partie de $\mathcal{C}_f$ pour $x \geq 0$ et on la \textbf{réfléchit} par rapport à $Oy$ (symétrie axiale) pour obtenir la partie $x < 0$ ; la courbe devient paire.
\end{enumerate}
\textbf{Réflexe :} pour tracer, transformer quelques points remarquables de $\mathcal{C}_f$ (sommet, intersections avec les axes, asymptotes), puis esquisser la nouvelle courbe.
\end{methode}

\begin{exoresolu}[Fonctions associées — type Probatoire]
Soit $f$ la fonction définie sur $\mathbb{R}$ par $f(x) = x^2$, de courbe $\mathcal{P}$ (parabole sommet $O$).
\begin{enumerate}
\item Donner l'expression de la fonction $g$ dont la courbe se déduit de $\mathcal{P}$ par la translation de vecteur $2\,\vec{i} - 3\,\vec{j}$.
\item Préciser le sommet de la parabole $\mathcal{C}_g$.
\item On définit $h(x) = |f(x) - 4|$. Décrire la construction de $\mathcal{C}_h$ à partir de $\mathcal{C}_f$.
\end{enumerate}
\tcblower
\textbf{1.} La translation de vecteur $2\,\vec{i}$ (horizontale, vers la droite) transforme $f(x)$ en $f(x - 2)$ ; la translation de vecteur $-3\,\vec{j}$ (verticale, vers le bas) ajoute $-3$. Donc :
\[ g(x) = f(x - 2) - 3 = (x - 2)^2 - 3. \]
En développant : $g(x) = x^2 - 4x + 4 - 3 = x^2 - 4x + 1$.

\textbf{2.} La parabole $\mathcal{P}$ a pour sommet le point $O(0\,;\,0)$. Par la translation de vecteur $2\,\vec{i} - 3\,\vec{j}$, l'image de $O$ est le point $S(0 + 2\,;\,0 - 3)$, c'est-à-dire $S(2\,;\,-3)$.
Vérification : la forme canonique $g(x) = (x-2)^2 - 3$ montre que le minimum de $g$ est $-3$, atteint en $x = 2$.

\textbf{3.} $\mathcal{C}_h$ s'obtient en deux étapes :
\begin{itemize}
\item Tracer $\mathcal{C}_1$ : courbe de $y = f(x) - 4 = x^2 - 4$ (translation de $\mathcal{P}$ de vecteur $-4\,\vec{j}$). Sommet en $(0\,;\,-4)$, intersections avec $Ox$ en $x = \pm 2$.
\item Appliquer la valeur absolue : $h(x) = |x^2 - 4|$. La partie de $\mathcal{C}_1$ au-dessus de $Ox$ (pour $|x| \geq 2$) est conservée. La partie en dessous (pour $|x| < 2$, sommet en $(0\,;\,-4)$) est réfléchie par rapport à $Ox$ : le sommet devient $(0\,;\,4)$, les points $(\pm 2\,;\,0)$ restent fixes. La courbe $\mathcal{C}_h$ a une forme de « W » arrondi.
\end{itemize}

\textbf{Conclusion :} $g(x) = (x-2)^2 - 3 = x^2 - 4x + 1$, sommet $S(2\,;\,-3)$. $\mathcal{C}_h$ est symétrique par rapport à $Oy$ (fonction paire), avec un maximum local en $(0\,;\,4)$ et deux minima en $(\pm 2\,;\,0)$.
\end{exoresolu}

\coursec{Applications concrètes : modélisation par une fonction}

\begin{methode}[Modéliser une situation de la vie courante par une fonction]
\begin{enumerate}
\item \textbf{Identifier} la variable indépendante (souvent notée $x$ : temps, longueur, quantité produite) et la grandeur dépendante (volume, coût, recette, surface...).
\item \textbf{Traduire} les données de l'énoncé en une expression algébrique $f(x)$, en précisant l'ensemble de définition \textbf{adapté au contexte réel} (longueurs strictement positives, quantités entières, temps positif, etc.).
\item \textbf{Exploiter} le modèle mathématique : calculer des images $f(a)$, résoudre l'équation $f(x) = k$ ou l'inéquation $f(x) \geq k$ pour répondre à la question posée.
\item \textbf{Conclure} par une phrase rédigée dans le langage de la situation (avec les unités : cm, cm$^3$, FCFA, s, etc.), et vérifier la cohérence du résultat (ordre de grandeur, signe).
\end{enumerate}
\end{methode}

\begin{exoresolu}[Situation concrète — atelier de menuiserie, type Probatoire]
Un menuisier de Yaoundé découpe une plaque rectangulaire de \SI{60}{cm} de long et \SI{40}{cm} de large. Il retire à chaque coin un carré de côté $x$ (en cm) pour plier la plaque et fabriquer une caisse sans couvercle. On note $V(x)$ le volume de la caisse en cm$^3$.
\begin{enumerate}
\item Justifier que $V(x) = x(60 - 2x)(40 - 2x)$ et préciser l'ensemble de définition de $V$ dans le contexte.
\item Calculer $V(5)$ et $V(10)$. Interpréter les résultats.
\item \textbf{Approfondissement :} On cherche la valeur de $x$ qui maximise le volume. On admet que le maximum est atteint pour $x = 10 - \frac{10\sqrt{3}}{3} \approx 4,2$. Donner le volume maximal arrondi au cm$^3$ près.
\end{enumerate}
\tcblower
\textbf{1. Expression du volume.} Après découpe des quatre carrés de côté $x$ et pliage, la caisse est un pavé droit de :
\begin{itemize}
\item hauteur : $x$ cm ;
\item longueur : $60 - 2x$ cm (on retire $x$ à chaque extrémité de la longueur) ;
\item largeur : $40 - 2x$ cm.
\end{itemize}
Le volume d'un pavé droit est le produit de ses trois dimensions :
\[ V(x) = x(60 - 2x)(40 - 2x). \]
\textbf{Ensemble de définition (contexte).} Les dimensions doivent être strictement positives pour former une caisse :
\[ x > 0, \quad 60 - 2x > 0 \Leftrightarrow x < 30, \quad 40 - 2x > 0 \Leftrightarrow x < 20. \]
La condition la plus restrictive est $x < 20$. Donc $D_V = ]0\,;\,20[$.

\textbf{2. Calculs.}
\[ V(5) = 5 \times (60 - 10) \times (40 - 10) = 5 \times 50 \times 30 = \num{7500}. \]
\[ V(10) = 10 \times (60 - 20) \times (40 - 20) = 10 \times 40 \times 20 = \num{8000}. \]
\textbf{Interprétation :} en retirant des carrés de \SI{5}{cm}, la caisse a un volume de \SI{7500}{cm^3} ; avec des carrés de \SI{10}{cm}, le volume est de \SI{8000}{cm^3}. La caisse est donc plus volumineuse pour $x = 10$ que pour $x = 5$.

\textbf{3. Volume maximal.} Pour $x_0 = 10 - \frac{10\sqrt{3}}{3} \approx 4,226$ :
\begin{align*}
V_{\text{max}} &= x_0(60-2x_0)(40-2x_0) \\
&\approx 4,226 \times 51,548 \times 31,548 \\
&\approx \num{6859} \text{ cm}^3 \quad (\text{arrondi au cm}^3 \text{ près}).
\end{align*}
\textbf{Conclusion :} $V(x) = x(60-2x)(40-2x)$ sur $]0\,;\,20[$, avec $V(5) = \SI{7500}{cm^3}$, $V(10) = \SI{8000}{cm^3}$. Le volume maximal admis est d'environ \SI{6859}{cm^3} pour une découpe d'environ \SI{4,2}{cm}.
\end{exoresolu}

\begin{attention}[Les 5 erreurs qui coûtent des points au Probatoire]
\begin{enumerate}
\item \textbf{Oublier de vérifier la symétrie de $D_f$ avant la parité.} Si l'ensemble de définition n'est pas symétrique par rapport à $0$, la fonction ne peut être ni paire ni impaire : le calcul de $f(-x)$ est alors inutile et la conclusion hâtive est fausse.
\item \textbf{Confondre $g \circ f$ et $f \circ g$.} $(g \circ f)(x) = g(f(x))$ : c'est $f$ qui agit en premier. Écrire $g(f(x)) = f(g(x))$ sans justification est une erreur grave, car la composition n'est pas commutative.
\item \textbf{Oublier une contrainte dans l'ensemble de définition.} Avec un dénominateur \emph{et} une racine carrée, les deux conditions doivent être traitées et \textbf{intersectées}. De plus, une valeur interdite à l'intérieur d'un intervalle oblige à le couper : $[-2\,;\,+\infty[ \setminus \{3\} = [-2\,;\,3[ \cup ]3\,;\,+\infty[$.
\item \textbf{Se tromper de sens dans les translations.} La courbe de $y = f(x - a)$ s'obtient en déplaçant $\mathcal{C}_f$ vers la \textbf{droite} quand $a > 0$ (et non vers la gauche). De même, $y = f(x) + k$ déplace vers le haut si $k > 0$.
\item \textbf{Conclure sans phrase et sans unité dans les problèmes concrets.} Au Probatoire, une réponse numérique non interprétée (volume sans cm$^3$, coût sans FCFA) ou une question de justification laissée sans rédaction (« Justifier que... ») fait perdre des points, même si le calcul est exact.
\end{enumerate}
\end{attention}

\begin{savaistu}[Le saviez-vous ? — Histoire et métier]
Le concept de \textbf{fonction} a été formalisé au XVIII\textsuperscript{e} siècle par Euler (notation $f(x)$) et Dirichlet (définition ensembliste). En technique, les fonctions modélisent tout : la \textbf{loi d'Ohm} $U = RI$ (fonction affine), la \textbf{portée d'un jet d'eau} (parabole), le \textbf{temps de charge d'un condensateur} (exponentielle), le \textbf{coût de production} (fonction polynôme). Maîtriser le vocabulaire (domaine, image, variations, parité) permet de \textbf{lire une courbe comme un plan de construction} : on y voit les maxima (performance), les zéros (seuils), les asymptotes (limites physiques).
\end{savaistu}

\newpage

\coursec{Exercices}

\courssub{Je vérifie mes connaissances}

\begin{exercice}[QCM : Vocabulaire et définitions]
Parmi les affirmations suivantes, une seule est exacte. Laquelle ? Justifier brièvement.
\begin{enumerate}
    \item Soit $f(x)=\dfrac{1}{x-2}$. L'ensemble de définition de $f$ est $\mathbb{R}$.
    \item Soit $g(x)=\sqrt{x+3}$. L'ensemble de définition de $g$ est $[-3, +\infty[$.
    \item La fonction $h(x)=x^2+1$ est une bijection de $\mathbb{R}$ dans $\mathbb{R}$.
    \item La fonction $k(x)=x^3$ est paire sur $\mathbb{R}$.
\end{enumerate}
\end{exercice}

\begin{exercice}[Vrai ou Faux ? Justifier]
Dire si les affirmations suivantes sont vraies ou fausses. Justifier chaque réponse par un calcul ou un contre-exemple.
\begin{enumerate}
    \item Si $f$ et $g$ sont deux fonctions paires, alors $f+g$ est paire.
    \item Si $f$ est une fonction impaire définie sur $\mathbb{R}$, alors $f(0)=0$.
    \item La fonction $f(x)=\sin(x)$ est $2\pi$-périodique et impaire.
    \item Pour toute fonction $f$ définie sur $\mathbb{R}$, la courbe de $x \mapsto f(x+2)$ s'obtient par translation de la courbe de $f$ de vecteur $\vec{u}(2,0)$.
\end{enumerate}
\end{exercice}

\begin{exercice}[Calculs directs : opérations et composition]
Soient les fonctions $f$ et $g$ définies sur $\mathbb{R}$ par $f(x)=2x-1$ et $g(x)=x^2+3$.
\begin{enumerate}
    \item Déterminer les expressions simplifiées des fonctions $f+g$, $f-g$, $fg$ et $\dfrac{f}{g}$ (préciser l'ensemble de définition pour le quotient).
    \item Calculer $(f \circ g)(x)$ et $(g \circ f)(x)$. Ces deux fonctions sont-elles égales ?
\end{enumerate}
\end{exercice}

\begin{exercice}[Définitions et propriétés]
\begin{enumerate}
    \item Donner la définition rigoureuse d'une fonction \cle{paire} et d'une fonction \cle{impaire}. Préciser la condition sur l'ensemble de définition $D$.
    \item Quelle est la symétrie de la courbe représentative d'une fonction paire ? D'une fonction impaire ?
    \item Définir la \cle{périodicité} d'une fonction. Si $T$ est une période de $f$, que vaut $f(x+2T)$ ?
\end{enumerate}
\end{exercice}

\courssub{Je m'entraîne}

\begin{exercice}[Opérations sur les fonctions : domaine et expression]
Soient les fonctions $f(x)=\sqrt{x-1}$ et $g(x)=\dfrac{1}{x-2}$.
\begin{enumerate}
    \item Déterminer les ensembles de définition $D_f$ et $D_g$.
    \item Déterminer l'ensemble de définition et l'expression simplifiée des fonctions $f+g$, $fg$ et $\dfrac{f}{g}$.
    \item Résoudre l'équation $(f+g)(x)=0$.
\end{enumerate}
\end{exercice}

\begin{exercice}[Composition de fonctions et ensemble de définition]
Soient $f(x)=\dfrac{2}{x+1}$ ($x \neq -1$) et $g(x)=\sqrt{x}$ ($x \geqslant 0$).
\begin{enumerate}
    \item Déterminer l'ensemble de définition et l'expression de la fonction composée $f \circ g$.
    \item Déterminer l'ensemble de définition et l'expression de la fonction composée $g \circ f$.
    \item Résoudre l'inéquation $(g \circ f)(x) \leqslant 1$.
\end{enumerate}
\end{exercice}

\begin{exercice}[Étude de la parité et symétrie]
Étudier la parité des fonctions suivantes sur leur ensemble de définition. En déduire, le cas échéant, le type de symétrie de leur courbe représentative.
\begin{enumerate}
    \item $f(x) = x^4 - 3x^2 + 2$
    \item $g(x) = \dfrac{x}{x^2+1}$
    \item $h(x) = x^3 + x + 1$
    \item $k(x) = \sqrt{4-x^2}$
\end{enumerate}
\end{exercice}

\begin{exercice}[Périodicité et tracé de courbe]
La fonction $f$ est $4$-périodique sur $\mathbb{R}$ et paire. Sur l'intervalle $[0, 2]$, elle est définie par $f(x) = x(2-x)$.
\begin{enumerate}
    \item Tracer la courbe $\mathcal{C}_f$ sur l'intervalle $[0, 2]$.
    \item En utilisant la parité, compléter le tracé sur $[-2, 2]$.
    \item En utilisant la périodicité, tracer $\mathcal{C}_f$ sur $[-4, 4]$.
    \item Déterminer l'ensemble des solutions de l'équation $f(x)=1$ sur $[-4, 4]$.
\end{enumerate}
\end{exercice}

\courssub{Je me prépare au Probatoire}

\begin{exercice}[Étude complète d'une fonction homographique — Contexte : Électricité]
\label{exo:probatoire1}
Dans un circuit électrique, la tension $U$ (en volts) aux bornes d'un dipôle non linéaire en fonction de l'intensité $I$ (en ampères) est modélisée par la fonction $f$ définie sur $\mathbb{R}^* = \mathbb{R} \setminus \{3\}$ par :
\[ f(x) = \frac{2x+1}{x-3} \]
\begin{enumerate}
    \item Déterminer les limites de $f$ aux bornes de son ensemble de définition. En déduire les asymptotes (verticale et horizontale) de la courbe $\mathcal{C}_f$.
    \item Calculer la dérivée $f'(x)$ sur $\mathbb{R}^*$. Dresser le tableau de variations de $f$.
    \item Montrer que la courbe $\mathcal{C}_f$ admet un centre de symétrie. Donner ses coordonnées.
    \item Tracer $\mathcal{C}_f$ dans un repère orthonormé (unité : 2 cm pour 1 sur l'axe des abscisses, 1 cm pour 1 sur l'axe des ordonnées).
    \item Résoudre graphiquement l'inéquation $f(x) \geqslant 0$.
\end{enumerate}
\end{exercice}

\begin{exercice}[Parité, périodicité et transformations — Contexte : Traitement de signal]
\label{exo:probatoire2}
Soit la fonction $g$ définie sur $\mathbb{R}$ par $g(x) = x^3 - 3x$.
\begin{enumerate}
    \item Étudier la parité de $g$. En déduire la symétrie de sa courbe $\mathcal{C}_g$.
    \item Calculer $g'(x)$ et dresser le tableau de variations de $g$ sur $\mathbb{R}$.
    \item Tracer $\mathcal{C}_g$ sur $[-3, 3]$.
\end{enumerate}
On définit maintenant une fonction $h$ sur $\mathbb{R}$ par $h(x) = g(x-2)$.
\begin{enumerate}
    \setcounter{enumi}{3}
    \item Exprimer $h(x)$ sous forme développée. $h$ est-elle paire ou impaire ? Justifier.
    \item Comment obtient-on la courbe $\mathcal{C}_h$ à partir de $\mathcal{C}_g$ ? Tracer $\mathcal{C}_h$ sur $[-1, 5]$.
\end{enumerate}
\end{exercice}

\begin{exercice}[Construction de courbe par transformations successives — Contexte : Dessin technique / CAO]
\label{exo:probatoire3}
On considère la fonction de référence $f_0(x) = |x|$ dont la courbe $\mathcal{C}_0$ est connue (bissectrices des 1er et 2e quadrants).
On veut tracer la courbe $\mathcal{C}$ de la fonction $h$ définie sur $\mathbb{R}$ par :
\[ h(x) = 2|x-3| - 1 \]
\begin{enumerate}
    \item Écrire $h(x)$ sous la forme $a f_0(x - b) + c$ en identifiant les réels $a$, $b$, $c$.
    \item Décrire précisément, dans l'ordre, les \textbf{trois transformations géométriques} successives permettant de passer de $\mathcal{C}_0$ à $\mathcal{C}$.
    \item Déterminer les coordonnées du sommet de $\mathcal{C}$ (point anguleux).
    \item Tracer $\mathcal{C}$ dans un repère orthonormé (unité 2 cm). Indiquer les points de rencontre avec les axes.
    \item Résoudre l'équation $h(x) = 3$ et l'inéquation $h(x) < 1$.
\end{enumerate}
\end{exercice}

\newpage

\coursec{Corrigés des exercices}

\begin{corrige}[Exercice 1 — QCM : Vocabulaire et définitions]
La seule affirmation exacte est la \textbf{b)}.
\begin{enumerate}
    \item \textbf{Fausse.} $f(x)=\dfrac{1}{x-2}$ n'est définie que si $x-2 \neq 0$, c'est-à-dire $x \neq 2$. Donc $D_f = \mathbb{R} \setminus \{2\}$.
    \item \textbf{Exacte.} $g(x)=\sqrt{x+3}$ est définie si et seulement si $x+3 \geqslant 0$, soit $x \geqslant -3$. Donc $D_g = [-3, +\infty[$.
    \item \textbf{Fausse.} $h(x)=x^2+1$ n'est pas injective : $h(1)=h(-1)=2$ alors que $1 \neq -1$. De plus, elle n'est pas surjective : l'équation $x^2+1=0$ n'a pas de solution réelle.
    \item \textbf{Fausse.} $k(-x)=(-x)^3=-x^3=-k(x)$ : la fonction cube est \cle{impaire}, pas paire.
\end{enumerate}
\end{corrige}

\begin{corrige}[Exercice 2 — Vrai ou Faux ? Justifier]
\begin{enumerate}
    \item \textbf{Vrai.} Si $f$ et $g$ sont paires, alors pour tout $x$ dans le domaine commun : $(f+g)(-x)=f(-x)+g(-x)=f(x)+g(x)=(f+g)(x)$. Donc $f+g$ est paire.
    \item \textbf{Vrai.} Si $f$ est impaire, $f(-0)=-f(0)$, c'est-à-dire $f(0)=-f(0)$, donc $2f(0)=0$, d'où $f(0)=0$.
    \item \textbf{Vrai.} Pour tout réel $x$ : $\sin(x+2\pi)=\sin(x)$ (périodicité de période $2\pi$) et $\sin(-x)=-\sin(x)$ (imparité).
    \item \textbf{Fausse.} La courbe de $x \mapsto f(x+2)$ s'obtient par translation de la courbe de $f$ de vecteur $\vec{u}(-2,0)$ (vers la \emph{gauche}), et non de vecteur $(2,0)$. Contre-exemple : $f(x)=x^2$ ; la courbe de $f(x+2)=(x+2)^2$ a son sommet en $(-2,0)$, image du sommet $(0,0)$ par la translation de vecteur $(-2,0)$.
\end{enumerate}
\end{corrige}

\begin{corrige}[Exercice 3 — Calculs directs : opérations et composition]
On a $f(x)=2x-1$ et $g(x)=x^2+3$, toutes deux définies sur $\mathbb{R}$.
\begin{enumerate}
    \item \textbf{Opérations :}
    \begin{align*}
    (f+g)(x) &= 2x-1+x^2+3 = x^2+2x+2 \\
    (f-g)(x) &= 2x-1-(x^2+3) = -x^2+2x-4 \\
    (fg)(x) &= (2x-1)(x^2+3) = 2x^3-x^2+6x-3
    \end{align*}
    Pour le quotient : $g(x)=x^2+3 \geqslant 3 > 0$ pour tout réel $x$, donc $g$ ne s'annule jamais. Ainsi $\dfrac{f}{g}$ est définie sur $\mathbb{R}$ tout entier :
    \[ \left(\frac{f}{g}\right)(x) = \frac{2x-1}{x^2+3}, \qquad D_{f/g}=\mathbb{R}. \]
    \item \textbf{Compositions :}
    \[ (f \circ g)(x) = f(g(x)) = 2(x^2+3)-1 = 2x^2+5. \]
    \[ (g \circ f)(x) = g(f(x)) = (2x-1)^2+3 = 4x^2-4x+1+3 = 4x^2-4x+4. \]
    Ces deux fonctions ne sont \textbf{pas égales} : par exemple $(f\circ g)(0)=5$ alors que $(g\circ f)(0)=4$. La composition n'est pas commutative en général.
\end{enumerate}
\end{corrige}

\begin{corrige}[Exercice 4 — Définitions et propriétés]
\begin{enumerate}
    \item \textbf{Fonction paire :} $f$ est paire si son ensemble de définition $D$ est symétrique par rapport à $0$ (c'est-à-dire : pour tout $x \in D$, $-x \in D$) et si pour tout $x \in D$ : $f(-x)=f(x)$.
    
    \textbf{Fonction impaire :} $f$ est impaire si $D$ est symétrique par rapport à $0$ et si pour tout $x \in D$ : $f(-x)=-f(x)$.
    \item La courbe d'une fonction \cle{paire} est symétrique par rapport à \textbf{l'axe des ordonnées} (symétrie axiale). La courbe d'une fonction \cle{impaire} est symétrique par rapport à \textbf{l'origine du repère} (symétrie centrale).
    \item $f$ est \cle{périodique} de période $T$ ($T>0$) si $D$ est stable par l'ajout de $T$ et si pour tout $x \in D$ : $f(x+T)=f(x)$. Par récurrence immédiate : $f(x+2T)=f(x+T)=f(x)$. Donc $f(x+2T)=f(x)$ : $2T$ est aussi une période de $f$.
\end{enumerate}
\end{corrige}

\begin{corrige}[Exercice 5 — Opérations sur les fonctions : domaine et expression]
\begin{enumerate}
    \item $f(x)=\sqrt{x-1}$ exige $x-1 \geqslant 0$, donc $D_f=[1,+\infty[$. $g(x)=\dfrac{1}{x-2}$ exige $x \neq 2$, donc $D_g=\mathbb{R}\setminus\{2\}$.
    \item \textbf{Somme et produit :} le domaine commun est $D_f \cap D_g = [1,+\infty[ \setminus \{2\} = [1,2[ \cup ]2,+\infty[$.
    \[ (f+g)(x) = \sqrt{x-1}+\frac{1}{x-2} = \frac{(x-2)\sqrt{x-1}+1}{x-2}, \qquad (fg)(x)=\frac{\sqrt{x-1}}{x-2}. \]
    \textbf{Quotient :} il faut de plus $g(x) \neq 0$ ; or $g(x)=\dfrac{1}{x-2}$ ne s'annule jamais. Donc $D_{f/g}=[1,2[\cup]2,+\infty[$ et
    \[ \left(\frac{f}{g}\right)(x) = \sqrt{x-1}\,(x-2) = (x-2)\sqrt{x-1}. \]
    \item Sur $[1,2[\cup]2,+\infty[$, $(f+g)(x)=0 \iff \sqrt{x-1}=-\dfrac{1}{x-2}$.
    \begin{itemize}
        \item Si $x>2$ : $\sqrt{x-1} \geqslant 0$ et $-\dfrac{1}{x-2}<0$ : pas de solution.
        \item Si $1 \leqslant x < 2$ : en multipliant par $(x-2)<0$ : $(x-2)\sqrt{x-1}=-1$. Posons $u=\sqrt{x-1} \geqslant 0$, alors $x=u^2+1$ et $x-2=u^2-1$. L'équation devient $(u^2-1)u=-1$, soit $u^3-u+1=0$. Pour $u \in [0,1[$ (car $x<2$), $u^3-u+1 = 1-u+u^3 > 0$ car $1-u>0$ et $u^3 \geqslant 0$. Pas de solution.
    \end{itemize}
    \textbf{Conclusion :} l'équation $(f+g)(x)=0$ n'a \textbf{aucune solution} : $\mathcal{S}=\varnothing$. (On peut aussi remarquer que sur $[1,2[$, $f(x) \geqslant 0$ et $g(x)<0$, et vérifier numériquement que $f+g>0$ : par exemple en $x=1$, $(f+g)(1)=-1$... en fait $(f+g)(1)=0+\frac{1}{-1}=-1<0$ ; mais pour $x$ proche de $2^-$, $g(x) \to -\infty$. Le signe n'est donc pas constant : reprenons.) 
    
    \emph{Analyse correcte :} sur $[1,2[$, posons $\varphi(x)=\sqrt{x-1}+\dfrac{1}{x-2}$. En $x=1$ : $\varphi(1)=-1$. Quand $x \to 2^-$ : $\varphi(x) \to -\infty$. Quand $x \to 1^+$... $\varphi$ reste négative car $\sqrt{x-1} \leqslant 1$ et $\left|\dfrac{1}{x-2}\right|>1$ pour $x \in [1,2[$... plus précisément, pour $x \in [1,2[$ : $\sqrt{x-1}<1$ et $\dfrac{1}{x-2} \leqslant -1$, donc $\varphi(x)<1-1=0$. Sur $]2,+\infty[$ : $\sqrt{x-1}>0$ et $\dfrac{1}{x-2}>0$, donc $\varphi(x)>0$. \textbf{Conclusion : $\mathcal{S}=\varnothing$.}
\end{enumerate}
\end{corrige}

\begin{corrige}[Exercice 6 — Composition de fonctions et ensemble de définition]
\begin{enumerate}
    \item \textbf{$f \circ g$ :} $(f\circ g)(x)=f(\sqrt{x})=\dfrac{2}{\sqrt{x}+1}$. Conditions : $x \geqslant 0$ (pour $g$) et $\sqrt{x}+1 \neq 0$, ce qui est toujours vrai car $\sqrt{x}+1 \geqslant 1 > 0$. Donc $D_{f\circ g}=[0,+\infty[$.
    \item \textbf{$g \circ f$ :} $(g\circ f)(x)=\sqrt{\dfrac{2}{x+1}}$. Conditions : $x \neq -1$ et $\dfrac{2}{x+1} \geqslant 0$. Comme $2>0$, il faut $x+1>0$, soit $x>-1$. Donc $D_{g\circ f}=]-1,+\infty[$.
    \item \textbf{Inéquation} $(g\circ f)(x) \leqslant 1$ sur $]-1,+\infty[$ :
    \[ \sqrt{\frac{2}{x+1}} \leqslant 1 \iff \frac{2}{x+1} \leqslant 1 \iff 2 \leqslant x+1 \iff x \geqslant 1 \]
    (on a pu élever au carré car les deux membres sont positifs, puis multiplier par $x+1>0$). Donc $\mathcal{S}=[1,+\infty[$.
\end{enumerate}
\end{corrige}

\begin{corrige}[Exercice 7 — Étude de la parité et symétrie]
\begin{enumerate}
    \item $f(x)=x^4-3x^2+2$, $D_f=\mathbb{R}$ (symétrique par rapport à $0$). $f(-x)=(-x)^4-3(-x)^2+2=x^4-3x^2+2=f(x)$. $f$ est \textbf{paire} : courbe symétrique par rapport à l'axe des ordonnées.
    \item $g(x)=\dfrac{x}{x^2+1}$, $D_g=\mathbb{R}$ car $x^2+1>0$. $g(-x)=\dfrac{-x}{(-x)^2+1}=-\dfrac{x}{x^2+1}=-g(x)$. $g$ est \textbf{impaire} : courbe symétrique par rapport à l'origine.
    \item $h(x)=x^3+x+1$, $D_h=\mathbb{R}$. $h(-x)=-x^3-x+1$. On n'a ni $h(-x)=h(x)$ (par exemple $h(1)=3 \neq h(-1)=-1$) ni $h(-x)=-h(x)$ (car $-h(x)=-x^3-x-1 \neq h(-x)$). $h$ est \textbf{ni paire ni impaire} : pas de symétrie particulière.
    \item $k(x)=\sqrt{4-x^2}$, $D_k=[-2,2]$ (car $4-x^2 \geqslant 0 \iff -2 \leqslant x \leqslant 2$), ensemble symétrique par rapport à $0$. $k(-x)=\sqrt{4-(-x)^2}=\sqrt{4-x^2}=k(x)$. $k$ est \textbf{paire} : courbe symétrique par rapport à l'axe des ordonnées (c'est d'ailleurs un demi-cercle de centre $O$ et de rayon 2).
\end{enumerate}
\end{corrige}

\begin{corrige}[Exercice 8 — Périodicité et tracé de courbe]
Sur $[0,2]$, $f(x)=x(2-x)=2x-x^2=-(x-1)^2+1$ : parabole de sommet $(1,1)$, passant par $(0,0)$ et $(2,0)$.
\begin{enumerate}
    \item Sur $[0,2]$ : arc de parabole joignant $(0,0)$ à $(2,0)$, de sommet $(1,1)$.
    \item \textbf{Parité :} $f$ est paire, donc la courbe sur $[-2,0]$ est le symétrique de la courbe sur $[0,2]$ par rapport à l'axe des ordonnées : arc joignant $(-2,0)$ à $(0,0)$, de sommet $(-1,1)$.
    \item \textbf{Périodicité :} $f$ est $4$-périodique, donc le motif obtenu sur $[-2,2]$ (de longueur 4) se répète par translations de vecteurs $4\vec{\imath}$ et $-4\vec{\imath}$. Sur $[2,4]$ : $f(x)=f(x-4)$, arc joignant $(2,0)$ à $(4,0)$ de sommet $(3,1)$ ; sur $[-4,-2]$ : arc de sommet $(-3,1)$.
    \item \textbf{Équation $f(x)=1$ :} sur $[0,2]$, $-(x-1)^2+1=1 \iff (x-1)^2=0 \iff x=1$. Par parité : $x=-1$ est aussi solution. Par $4$-périodicité : $x=1+4=5 \notin [-4,4]$, $x=1-4=-3 \in [-4,4]$, $x=-1+4=3 \in [-4,4]$, $x=-1-4=-5 \notin [-4,4]$. Donc $\mathcal{S}=\{-3,-1,1,3\}$ (les sommets des arcs).
\end{enumerate}
\end{corrige}

\begin{corrige}[Exercice 9 — Étude complète d'une fonction homographique (Électricité)]
\emph{Barème indicatif (sur 10) : question 1 : 2 pts ; question 2 : 2,5 pts ; question 3 : 2 pts ; question 4 : 2 pts ; question 5 : 1,5 pt.}

\begin{enumerate}
    \item \textbf{Limites et asymptotes.} $f(x)=\dfrac{2x+1}{x-3}$ sur $\mathbb{R}\setminus\{3\}$.
    \begin{itemize}
        \item En $+\infty$ et $-\infty$ : $\lim_{x \to \pm\infty} f(x) = \lim_{x \to \pm\infty} \dfrac{2x}{x} = 2$. Asymptote \textbf{horizontale} d'équation $y=2$.
        \item En $3$ : numérateur $2(3)+1=7>0$ ; dénominateur $x-3 \to 0^+$ quand $x \to 3^+$ et $0^-$ quand $x \to 3^-$. Donc $\lim_{x \to 3^+} f(x)=+\infty$ et $\lim_{x \to 3^-} f(x)=-\infty$. Asymptote \textbf{verticale} d'équation $x=3$.
    \end{itemize}
    \item \textbf{Dérivée et variations.} Avec $u=2x+1$, $v=x-3$ : $u'=2$, $v'=1$.
    \[ f'(x)=\frac{2(x-3)-(2x+1)}{(x-3)^2}=\frac{2x-6-2x-1}{(x-3)^2}=\frac{-7}{(x-3)^2}. \]
    Pour tout $x \neq 3$, $(x-3)^2>0$, donc $f'(x)<0$ : $f$ est \textbf{strictement décroissante} sur $]-\infty,3[$ et sur $]3,+\infty[$ (attention : pas « sur $\mathbb{R}\setminus\{3\}$ » d'un seul tenant).
    \begin{center}
    \begin{tabular}{|c|ccccc|}\hline
    $x$ & $-\infty$ & & $3$ & & $+\infty$ \\ \hline
    $f'(x)$ & & $-$ & & $-$ & \\ \hline
    $f$ & $2$ & $\searrow$ & & $\searrow$ & $2$ \\ \hline
    \end{tabular}
    \end{center}
    (En $3$ : $f(x) \to -\infty$ à gauche, $+\infty$ à droite.)
    \item \textbf{Centre de symétrie.} Écrivons la forme réduite : $f(x)=\dfrac{2(x-3)+7}{x-3}=2+\dfrac{7}{x-3}$. Posons $X=x-3$ et $Y=y-2$ : alors $Y=\dfrac{7}{X}$, fonction impaire en $X$. Donc le point $\Omega(3,2)$ est \textbf{centre de symétrie} de $\mathcal{C}_f$. Vérification : $\dfrac{f(3+t)+f(3-t)}{2}=\dfrac{1}{2}\left(2+\dfrac{7}{t}+2-\dfrac{7}{t}\right)=2$ pour tout $t \neq 0$.
    \item \textbf{Tracé.} Hyperbole passant par $(0,-\tfrac{1}{3})$ et $(-\tfrac{1}{2},0)$, asymptotes $x=3$ et $y=2$, centrée en $\Omega(3,2)$.
    \item \textbf{Inéquation $f(x) \geqslant 0$.} $f(x)=0 \iff 2x+1=0 \iff x=-\dfrac{1}{2}$. Graphiquement (ou par tableau de signes : $2x+1$ s'annule en $-\tfrac12$, $x-3$ en $3$) : $\mathcal{S}=\left[-\dfrac{1}{2}, 3\right[$.
    
    \textbf{Points de vigilance :} citer la condition $x \neq 3$ avant tout calcul ; ne jamais écrire « $f$ décroissante sur $\mathbb{R}\setminus\{3\}$ » sans séparer les deux intervalles ; pour le centre de symétrie, la vérification $\frac{f(a+t)+f(a-t)}{2}=b$ est la justification attendue ; dans l'inéquation, la valeur $3$ est exclue (valeur interdite).
\end{enumerate}
\end{corrige}

\begin{corrige}[Exercice 10 — Parité, périodicité et transformations (Traitement de signal)]
\emph{Barème indicatif (sur 10) : question 1 : 1,5 pt ; question 2 : 2,5 pts ; question 3 : 2 pts ; question 4 : 2 pts ; question 5 : 2 pts.}
\begin{enumerate}
    \item $g(x)=x^3-3x$, $D_g=\mathbb{R}$ symétrique par rapport à $0$. $g(-x)=(-x)^3-3(-x)=-x^3+3x=-(x^3-3x)=-g(x)$. $g$ est \textbf{impaire} : $\mathcal{C}_g$ est symétrique par rapport à l'origine $O$.
    \item $g'(x)=3x^2-3=3(x^2-1)=3(x-1)(x+1)$. $g'(x)>0$ sur $]-\infty,-1[\cup]1,+\infty[$ et $g'(x)<0$ sur $]-1,1[$. Valeurs : $g(-1)=-1+3=2$ (maximum local), $g(1)=1-3=-2$ (minimum local).
    \begin{center}
    \begin{tabular}{|c|ccccccc|}\hline
    $x$ & $-\infty$ & & $-1$ & & $1$ & & $+\infty$ \\ \hline
    $g'(x)$ & & $+$ & $0$ & $-$ & $0$ & $+$ & \\ \hline
    $g$ & $-\infty$ & $\nearrow$ & $2$ & $\searrow$ & $-2$ & $\nearrow$ & $+\infty$ \\ \hline
    \end{tabular}
    \end{center}
    \item \textbf{Tracé sur $[-3,3]$ :} cubique passant par $(-\sqrt{3},0)$, $(0,0)$, $(\sqrt{3},0)$, sommets locaux $(-1,2)$ et $(1,-2)$, avec $g(3)=27-9=18$ et $g(-3)=-18$.
    \item $h(x)=g(x-2)=(x-2)^3-3(x-2)=x^3-6x^2+12x-8-3x+6=x^3-6x^2+9x-2$. $h$ n'est \textbf{ni paire ni impaire} : $h(0)=-2 \neq 0$ exclut l'imparité, et $h(1)=1-6+9-2=2$ alors que $h(-1)=-1-6-9-2=-18 \neq h(1)$ exclut la parité.
    \item $h(x)=g(x-2)$ : $\mathcal{C}_h$ s'obtient à partir de $\mathcal{C}_g$ par la \textbf{translation de vecteur $2\vec{\imath}$} (décalage de 2 unités vers la droite). Le centre de symétrie $O(0,0)$ de $\mathcal{C}_g$ devient $(2,0)$ pour $\mathcal{C}_h$ ; les extrema deviennent $(1,2)$ et $(3,-2)$. Sur $[-1,5]$, on trace la cubique translatée passant par $(2,0)$, $(2-\sqrt{3},0)$, $(2+\sqrt{3},0)$.
    
    \textbf{Points de vigilance :} pour la parité, toujours commencer par vérifier que l'ensemble de définition est symétrique par rapport à $0$ ; $g(x-2)$ correspond à une translation vers la \emph{droite} (et non vers la gauche) ; développer $(x-2)^3$ avec soin : $(x-2)^3=x^3-6x^2+12x-8$.
\end{enumerate}
\end{corrige}

\begin{corrige}[Exercice 11 — Construction de courbe par transformations successives (Dessin technique / CAO)]
\emph{Barème indicatif (sur 10) : question 1 : 1 pt ; question 2 : 3 pts ; question 3 : 1,5 pt ; question 4 : 2,5 pts ; question 5 : 2 pts.}
\begin{enumerate}
    \item $h(x)=2|x-3|-1=2f_0(x-3)-1$, donc $a=2$, $b=3$, $c=-1$.
    \item \textbf{Transformations successives} (dans l'ordre) :
    \begin{enumerate}
        \item \textbf{Translation} de vecteur $3\vec{\imath}$ : $\mathcal{C}_0 \to \mathcal{C}_1$ courbe de $x \mapsto |x-3|$ ;
        \item \textbf{Affinité orthogonale} d'axe $(Ox)$ et de rapport $2$ (dilatation verticale) : $\mathcal{C}_1 \to \mathcal{C}_2$ courbe de $x \mapsto 2|x-3|$ ;
        \item \textbf{Translation} de vecteur $-\vec{\jmath}$ (1 unité vers le bas) : $\mathcal{C}_2 \to \mathcal{C}$ courbe de $x \mapsto 2|x-3|-1$.
    \end{enumerate}
    \item Le sommet de $\mathcal{C}_0$ est $(0,0)$. Après translation de $3\vec{\imath}$ : $(3,0)$ ; l'affinité le laisse fixe ; après translation verticale de $-1$ : $\boxed{S(3,-1)}$.
    \item \textbf{Tracé} (unité 2 cm) : demi-droites de pentes $-2$ (pour $x<3$) et $+2$ (pour $x>3$) issues de $S(3,-1)$. Intersections avec les axes :
    \begin{itemize}
        \item Axe des ordonnées : $h(0)=2|{-3}|-1=5$, point $(0,5)$ ;
        \item Axe des abscisses : $2|x-3|-1=0 \iff |x-3|=\dfrac{1}{2} \iff x=\dfrac{5}{2}$ ou $x=\dfrac{7}{2}$, points $(2,5\,;\,0)$ et $(3,5\,;\,0)$.
    \end{itemize}
    \item \textbf{Équation $h(x)=3$ :} $2|x-3|-1=3 \iff |x-3|=2 \iff x-3=2$ ou $x-3=-2$, d'où $\mathcal{S}=\{1,5\}$.
    
    \textbf{Inéquation $h(x)<1$ :} $2|x-3|-1<1 \iff |x-3|<1 \iff -1<x-3<1 \iff 2<x<4$, d'où $\mathcal{S}=]2,4[$.
    
    \textbf{Points de vigilance :} respecter l'ordre des transformations (la dilatation verticale doit être faite avant la translation verticale, sinon le décalage serait doublé) ; $|x-3|<1$ équivaut à un encadrement $-1<x-3<1$, pas à deux inéquations disjointes ; bien distinguer les deux demi-droites de la courbe en « V ».
\end{enumerate}
\end{corrige}

\newpage

\coursec{Fiche bilan}

\begin{popfigure}
\begin{center}\tcbox[colback=poporange,colframe=poporange,coltext=white,arc=9pt,boxsep=4pt,left=6pt,right=6pt]{\bfseries\small Fonctions Numériques 1re Tech}\end{center}
\vspace{-2pt}
\begin{tcbraster}[raster columns=3,raster equal height=rows,raster column skip=6pt,raster row skip=6pt,enhanced,blankest]
\begin{tcolorbox}[enhanced,colback=white,colframe=poporange,boxrule=0.9pt,arc=3pt,title={1. Vocabulaire \&~Notation},fonttitle=\bfseries\scriptsize,coltitle=white,colbacktitle=poporange,left=4pt,right=4pt,top=2pt,bottom=2pt,boxsep=1pt,toptitle=1.5pt,bottomtitle=1.5pt,halign=left]
\begin{itemize}[nosep,leftmargin=*,itemsep=1pt,font=\scriptsize]
\item Domaine $D_f$
\item Image, Antécédent
\item Ensemble arrivée
\item Tableau/Valeurs
\end{itemize}
\end{tcolorbox}
\begin{tcolorbox}[enhanced,colback=white,colframe=popgreen,boxrule=0.9pt,arc=3pt,title={2. Opérations \&~Composition},fonttitle=\bfseries\scriptsize,coltitle=white,colbacktitle=popgreen,left=4pt,right=4pt,top=2pt,bottom=2pt,boxsep=1pt,toptitle=1.5pt,bottomtitle=1.5pt,halign=left]
\begin{itemize}[nosep,leftmargin=*,itemsep=1pt,font=\scriptsize]
\item $f\pm g$, $f\times g$, $f/g$
\item Domaine résultat
\item $g \circ f$ : ordre !
\item Non commutativité
\end{itemize}
\end{tcolorbox}
\begin{tcolorbox}[enhanced,colback=white,colframe=poppurple,boxrule=0.9pt,arc=3pt,title={3. Parité \&~Symétrie},fonttitle=\bfseries\scriptsize,coltitle=white,colbacktitle=poppurple,left=4pt,right=4pt,top=2pt,bottom=2pt,boxsep=1pt,toptitle=1.5pt,bottomtitle=1.5pt,halign=left]
\begin{itemize}[nosep,leftmargin=*,itemsep=1pt,font=\scriptsize]
\item Paire : $f(-x)=f(x)$
\item Impaire : $f(-x)=-f(x)$
\item Axe $Oy$ / Centre $O$
\item Domaine symétrique
\end{itemize}
\end{tcolorbox}
\begin{tcolorbox}[enhanced,colback=white,colframe=poppink,boxrule=0.9pt,arc=3pt,title={4. Périodicité},fonttitle=\bfseries\scriptsize,coltitle=white,colbacktitle=poppink,left=4pt,right=4pt,top=2pt,bottom=2pt,boxsep=1pt,toptitle=1.5pt,bottomtitle=1.5pt,halign=left]
\begin{itemize}[nosep,leftmargin=*,itemsep=1pt,font=\scriptsize]
\item $f(x+T)=f(x)$
\item Période $T>0$
\item Reproduction motif
\item Trigonométrie
\end{itemize}
\end{tcolorbox}
\begin{tcolorbox}[enhanced,colback=white,colframe=popgold,boxrule=0.9pt,arc=3pt,title={5. Courbe \&~Analyse},fonttitle=\bfseries\scriptsize,coltitle=white,colbacktitle=popgold,left=4pt,right=4pt,top=2pt,bottom=2pt,boxsep=1pt,toptitle=1.5pt,bottomtitle=1.5pt,halign=left]
\begin{itemize}[nosep,leftmargin=*,itemsep=1pt,font=\scriptsize]
\item Représentation graphique
\item Variations / TV
\item Extrema, Zéros
\item Lecture graphique
\end{itemize}
\end{tcolorbox}
\begin{tcolorbox}[enhanced,colback=white,colframe=popdark,boxrule=0.9pt,arc=3pt,title={6. Fonctions Associées},fonttitle=\bfseries\scriptsize,coltitle=white,colbacktitle=popdark,left=4pt,right=4pt,top=2pt,bottom=2pt,boxsep=1pt,toptitle=1.5pt,bottomtitle=1.5pt,halign=left]
\begin{itemize}[nosep,leftmargin=*,itemsep=1pt,font=\scriptsize]
\item Translation $f(x+a)$
\item Homothety $f(kx)$
\item Symétrie $-f(x)$
\item Valeur absolue $|f|$
\end{itemize}
\end{tcolorbox}
\end{tcbraster}

\legende{Carte mentale : Généralités sur les fonctions numériques (Leçon 12)}
\end{popfigure}

\begin{bilan}
\courssub{Rappels et vocabulaire}
\begin{itemize}
    \item \cle{Fonction} $f : D_f \to \mathbb{R}$ : à chaque $x \in D_f$ (domaine de définition) on associe un unique $y = f(x)$ (image).
    \item \cle{Ensemble de départ} $D_f$, \cle{Ensemble d'arrivée} $\mathbb{R}$, \cle{Ensemble images} $f(D_f) = \{f(x) \mid x \in D_f\}$.
    \item \cle{Antécédent} : $x$ tel que $f(x)=y$. Résolution d'équation $f(x)=k$.
    \item Représentations : algébrique (formule), graphique (courbe $\mathcal{C}_f$ dans repère), tableau de valeurs.
\end{itemize}

\courssub{Opérations sur les fonctions et composition}
\begin{tabularx}{\linewidth}{|l|X|X|}\hline
\rowcolor{popblueL}
\textbf{Opération} & \textbf{Définition} & \textbf{Domaine de définition} \\ \hline
Somme $f+g$ & $(f+g)(x) = f(x)+g(x)$ & $D_f \cap D_g$ \\ \hline
Différence $f-g$ & $(f-g)(x) = f(x)-g(x)$ & $D_f \cap D_g$ \\ \hline
Produit $f \times g$ & $(fg)(x) = f(x)g(x)$ & $D_f \cap D_g$ \\ \hline
Quotient $f/g$ & $(f/g)(x) = \frac{f(x)}{g(x)}$ & $D_f \cap D_g \cap \{x \mid g(x) \neq 0\}$ \\ \hline
Composition $g \circ f$ & $(g \circ f)(x) = g(f(x))$ & $\{x \in D_f \mid f(x) \in D_g\}$ \\ \hline
\end{tabularx}
\begin{attention}[Pièges classiques]
$\circ$ n'est \textbf{pas commutatif} : $g \circ f \neq f \circ g$ en général. L'ordre s'écrit de droite à gauche : on applique $f$ \emph{puis} $g$. Vérifier toujours la condition $f(x) \in D_g$ pour le domaine de $g \circ f$.
\end{attention}

\courssub{Parité et symétrie (Domaine $D_f$ symétrique par rapport à 0 : $x \in D_f \iff -x \in D_f$)}
\begin{tabularx}{\linewidth}{|l|X|X|}\hline
\rowcolor{popblueL}
\textbf{Type} & \textbf{Condition algébrique} & \textbf{Symétrie graphique} \\ \hline
Fonction \cle{paire} & $\forall x \in D_f,\ f(-x) = f(x)$ & Axe $Oy$ (axe des ordonnées) \\ \hline
Fonction \cle{impaire} & $\forall x \in D_f,\ f(-x) = -f(x)$ & Centre $O$ (origine du repère) \\ \hline
\end{tabularx}
\begin{methode}[Étudier la parité]
\begin{enumerate}
    \item Vérifier que $D_f$ est symétrique ($\forall x \in D_f, -x \in D_f$). Si non $\to$ ni paire, ni impaire.
    \item Calculer $f(-x)$ (simplifier).
    \item Comparer à $f(x)$ : égal $\to$ paire ; opposé $\to$ impaire ; autre $\to$ ni l'un ni l'autre.
    \item Conclure avec la phrase type.
\end{enumerate}
\end{methode}

\courssub{Périodicité}
\begin{definition}[Fonction périodique]
$f$ est \cle{périodique} de période $T > 0$ si : $\forall x \in D_f,\ x+T \in D_f$ et $f(x+T) = f(x)$.
La plus petite période $T > 0$ est la \cle{période fondamentale}.
\end{definition}
\begin{propriete}
Si $f$ est périodique de période $T$, sa courbe $\mathcal{C}_f$ est invariante par translation de vecteur $\vec{v} = (T, 0)$. L'étude sur $\mathbb{R}$ se ramène à l'étude sur un intervalle de longueur $T$ (ex: $[0, T[$ ou $[-T/2, T/2[$).
\end{propriete}

\courssub{Courbe d'une fonction : tracé et analyse}
\begin{itemize}
    \item \cle{Courbe représentative} $\mathcal{C}_f = \{(x, f(x)) \mid x \in D_f\}$ dans un repère orthogonal.
    \item \cle{Tableau de variations} : domaine, dérivée (si connue) ou monotonie, sens de variation ($\nearrow$ croissant, $\searrow$ décroissant), extrema, limites aux bornes.
    \item \cle{Lecture graphique} : image $f(a)$ = ordonnée du point d'abscisse $a$ ; antécédent(s) de $k$ = abscisse(s) intersection avec $y=k$ ; signe = position par rapport à l'axe $Ox$.
\end{itemize}

\courssub{Fonctions associées à une fonction donnée}
Soit $f$ définie sur $D_f$. On définit $g$ par transformation sur la variable ou l'image.
\begin{center}
\begin{tabularx}{\linewidth}{|l|l|X|}\hline
\rowcolor{popblueL}
\textbf{Transformation} & \textbf{Expression de $g$} & \textbf{Effet graphique sur $\mathcal{C}_f$} \\ \hline
Translation horizontale & $g(x) = f(x - a)$ & Translation de vecteur $\vec{v}(a, 0)$ \\ \hline
Translation verticale & $g(x) = f(x) + b$ & Translation de vecteur $\vec{v}(0, b)$ \\ \hline
Symétrie axe $Oy$ & $g(x) = f(-x)$ & Symétrie axiale par rapport à $Oy$ \\ \hline
Symétrie origine $O$ & $g(x) = -f(-x)$ & Symétrie centrale par rapport à $O$ \\ \hline
Symétrie axe $Ox$ & $g(x) = -f(x)$ & Symétrie axiale par rapport à $Ox$ \\ \hline
Homothetie horizontale & $g(x) = f(kx),\ k>0$ & Rapport $1/k$ par rapport à $Oy$ \\ \hline
Homothetie verticale & $g(x) = k f(x),\ k>0$ & Rapport $k$ par rapport à $Ox$ \\ \hline
Valeur absolue & $g(x) = |f(x)|$ & Partie $f(x)<0$ reflétée au-dessus de $Ox$ \\ \hline
\end{tabularx}
\end{center}
\begin{attention}[Ordre des transformations]
Pour $g(x) = a f(bx + c) + d$ : factoriser $bx+c = b(x + c/b)$. Ordre : 1. Homothétie horizontale ($1/b$), 2. Translation horizontale ($-c/b$), 3. Homothétie verticale ($a$), 4. Translation verticale ($d$).
\end{attention}
\end{bilan}

\begin{aretenir}[Checklist avant le Probatoire]
\begin{itemize}
    \item $\square$ Savoir donner le domaine $D_f$ d'une fonction (racine, dénominateur, log) et l'écrire en notation d'intervalle.
    \item $\square$ Calculer une image $f(a)$ et résoudre $f(x)=k$ (équation) ou $f(x) > k$ (inéquation) algébriquement ou graphiquement.
    \item $\square$ Déterminer le domaine d'une somme, produit, quotient, composition de deux fonctions.
    \item $\square$ Calculer l'expression simplifiée de $g \circ f$ et $f \circ g$ et distinguer les deux.
    \item $\square$ Étudier la parité d'une fonction : vérifier la symétrie du domaine, calculer $f(-x)$, conclure rigoureusement.
    \item $\square$ Utiliser la parité pour compléter un tableau de variations ou tracer la courbe sur $\mathbb{R}^-$ à partir de $\mathbb{R}^+$.
    \item $\square$ Identifier la période $T$ d'une fonction périodique (donnée ou lue sur graphique) et l'utiliser pour calculer $f(x)$ hors intervalle principal.
    \item $\square$ Construire un tableau de variations complet (valeurs aux bornes, variations, extrema) et en déduire le signe.
    \item $\square$ Tracer la courbe $\mathcal{C}_f$ dans un repère orthogonal (échelle, points caractéristiques, asymptotes si hors programme 1re mais notions utiles).
    \item $\square$ Appliquer les transformations graphiques (translation, symétrie, homothétie, valeur absolue) pour obtenir $\mathcal{C}_g$ à partir de $\mathcal{C}_f$ sans recalculer tout le tableau.
    \item $\square$ Lire graphiquement : images, antécédents, variations, signe, extrema, parité, périodicité.
    \item $\square$ Rédiger une conclusion complète : « La fonction $f$ est paire sur $D_f = \dots$, elle est croissante sur $\dots$, son maximum est $\dots$ ».
\end{itemize}
\end{aretenir}

\findecours

\end{document}
